\documentclass[twoside,10pt]{article}

\usepackage[T1]{fontenc}
\usepackage{lmodern}
\usepackage{amsmath,amsthm,amssymb,amsfonts}
\usepackage{mathtools}
\usepackage{longtable}
\usepackage[small,labelfont=bf,labelsep=period]{caption}
\usepackage[
  letterpaper,
  total={5.5in,8in},
  vmarginratio=1:1,
  hmarginratio=1:1,
  headsep=21pt
]{geometry}
\usepackage{xcolor}
\usepackage[
  breaklinks=false,
  linktocpage=true
]{hyperref}
% Referee copy: uncomment the next line for line numbers.
% \usepackage{lineno}\linenumbers
\hypersetup{
  colorlinks=true,
  linkcolor={purple},
  urlcolor={purple},
  citecolor={blue!80!black},
  pdfauthor={Leslie P. Polzer},
  pdftitle={Semiclassical wave propagation beyond the Ehrenfest time on hyperbolic surfaces},
  pdfsubject={Long-time semiclassical propagation of analytic Lagrangian states on hyperbolic surfaces},
  pdfkeywords={semiclassical analysis, Schroedinger equation, hyperbolic surfaces, Fuchsian groups, analytic amplitudes, Dyson expansion}
}

\medmuskip=4mu
\thickmuskip=8mu
\setlength{\emergencystretch}{2em}

\numberwithin{equation}{section}

\newtheoremstyle{slbody}
  {8pt}
  {8pt}
  {\normalfont\slshape}
  {}
  {\bfseries}
  {.}
  {0.5em}
  {\thmname{#1}\thmnumber{ #2}%
   \thmnote{ {\normalfont\itshape(\ignorespaces#3\unskip)}}}
\theoremstyle{slbody}
\newtheorem{theorem}{Theorem}[section]
\newtheorem{proposition}[theorem]{Proposition}
\newtheorem{lemma}[theorem]{Lemma}
\newtheorem{corollary}[theorem]{Corollary}
% Companion results share the ordinary statement counter.
\newtheorem{theoremR}[theorem]{Theorem}
\newtheorem{corollaryR}[theorem]{Corollary}
\theoremstyle{remark}
\newtheorem{remark}[theorem]{Remark}

\renewcommand{\qedsymbol}{$\blacksquare$}

\newcommand{\D}{\mathbb{D}}
\newcommand{\Hb}{\mathbb{H}}
\newcommand{\norm}[1]{\lVert #1\rVert}
\newcommand{\ue}{\mathrm{e}}
\newcommand{\cU}{\mathcal{U}}
\newcommand{\pcXzero}{X_0}
\newcommand{\pcLog}{\ell}
\newcommand{\pcw}{\omega}
\newcommand{\pcNorm}[1]{\norm{#1}_{\mathrm P}}

\title{Semiclassical wave propagation beyond the Ehrenfest time\\
on hyperbolic surfaces}
\author{Leslie P. Polzer}
\newcommand{\authoraffiliation}{Independent Researcher}
\newcommand{\authoremail}{polzer@fastmail.com}
\date{November 10, 2026}

\makeatletter
\renewcommand{\section}{%
  \@startsection{section}{1}{\z@}%
    {-3.5ex \@plus -1ex \@minus -.2ex}%
    {2.3ex \@plus .2ex}%
    {\normalfont\large\bfseries}}
\renewcommand{\@seccntformat}[1]{%
  \csname the#1\endcsname.\enspace}

\def\tagform@#1{%
  \maketag@@@{(\ignorespaces{\oldstylenums{#1}}\unskip\@@italiccorr)}}
\renewcommand{\eqref}[1]{%
  \textup{{\normalfont(\oldstylenums{\ref{#1}}}\normalfont)}}

\newcommand{\shorttitle}{SEMICLASSICAL PROPAGATION BEYOND THE EHRENFEST TIME}
\newcommand{\shortauthors}{LESLIE P. POLZER}
\newcommand{\ps@bookheader}{%
  \renewcommand\@oddfoot{\hfil}%
  \renewcommand\@evenfoot{\hfil}%
  \renewcommand\@oddhead{%
    \ifnum\value{page}>1
      {\small\hfil\shortauthors}\hfil\thepage
    \else
      \hfil
    \fi}%
  \renewcommand\@evenhead{%
    \thepage\hfil\small\shorttitle\hfil}}
\pagestyle{bookheader}

\renewcommand{\maketitle}{%
  \begin{center}
    {\large\bfseries\@title}\\
    \vskip36pt
    \ifx\@author\@empty
      \mbox{}\par
    \else
      {\scshape\@author}\par
      \medskip
      {\small\authoraffiliation}\par
      \smallskip
      {\small\texttt{\authoremail}}\par
      \smallskip
      {\small\@date}\par
    \fi
    \vskip36pt
  \end{center}}
\makeatother

\renewenvironment{abstract}
  {\quotation\noindent\small{\bfseries\abstractname.\enspace}}
  {\endquotation}

\begin{document}
\maketitle

\begin{abstract}
We study the Schr\"odinger evolution of analytic Lagrangian states attached to unstable manifolds on hyperbolic
quotients $M=\Hb/\Gamma$, for arbitrary Fuchsian groups $\Gamma$. Following R.~Schubert's strategy, we prove
that his explicit semiclassical approximations are accurate in $L^2(M)$ up to times
$\delta(\alpha/\hbar)^{1/2}$, where $\alpha$ is the analytic radius and $\alpha^3\ge\kappa\hbar$.
For fixed $\alpha$, this gives times of order $\hbar^{-1/2}$, beyond the logarithmic Ehrenfest scale.
The estimates control every power of the Laplacian and give exponentially small errors under suitable truncation.
They hold both for Schubert's original amplitude norm and for the reversed order of weights used in his proofs.
The source is his preprint \emph{Semiclassical wave propagation for large times} (arXiv:0705.0134v1, 2007),
withdrawn by its author in 2011 for a gap in the dispersive estimates. Two corrections are necessary:
the leading error is of order $\hbar t$, attained on the hyperbolic plane, and pointwise estimates on arbitrary
quotients require a local multiplicity factor whose square-root growth in the cusp is necessary on the modular surface.
For a family of states on the modular surface, the approximation remains accurate while mass escapes every compact
set after a sufficient time of order $\log(1/\hbar)\log\log(1/\hbar)$, in contrast with classical transport.
The main approximation results and the quantum escape estimates are formalized in Lean~4; the precise scope is given in the formalization appendix.

\vskip5pt
\noindent\textbf{Keywords.}\enspace Semiclassical analysis; time-dependent
Schr\"odinger equation; hyperbolic surfaces; Fuchsian groups; analytic
amplitudes; Dyson expansion.

\vskip5pt
\noindent\textbf{MSC2020 Classification.}\enspace
81Q20; 35Q41; 58J50; 30F35.
\end{abstract}

\vskip50pt
\baselineskip=13.5pt
\section{Introduction}\label{sec:intro}

Semiclassical approximations describe the solutions of the Schr\"odinger equation
$i\hbar\,\partial_t\psi=-\tfrac{\hbar^2}{2}\Delta\psi$ for small $\hbar$ through the geodesic flow. When the flow
is hyperbolic, rigorous approximations are in general under control only up to the Ehrenfest time
$T\sim\lambda^{-1}\log(1/\hbar)$, where $\lambda$ is a Lyapunov exponent (Section~\ref{subsec:context}).

In 2007 R.~Schubert \cite{S} proposed that for one family of states the error can be controlled much
longer. On a hyperbolic surface $M=\Hb/\Gamma$ he considered Lagrangian states $a\,\ue^{i\varphi_b/\hbar}$
whose phase $\varphi_b$ is attached to a boundary point $b$, so that the state is carried by an unstable
manifold of the geodesic flow. For such a state the evolution on the universal cover factorizes exactly
into the transport along the flow and a unitary dispersive part whose generator carries a factor $\hbar$
(equation \eqref{eq:factor} below). Expanding the dispersive part to order $K$ gives an explicit approximation
$u^{(K)}(t)$, and the preprint \cite{S} states that its periodization $u^{(K)}(t)_\Gamma$ is accurate in $L^2(M)$
up to times of order $\hbar^{-1/2}$, with specified dependence on the expansion and derivative orders, for amplitudes $a$ that are analytic in a strip
and decay exponentially. The preprint was withdrawn by its author in 2011 for ``a gap in the proof of the
dispersive estimates''.

Our principal result, Theorem~\ref{pc:main} and its extension in Corollary~\ref{pc:allalpha}, establishes the
$L^2$ approximation for Schubert's original amplitude norm on every Fuchsian quotient. Its corollaries give the leading error,
optimal truncation and pointwise estimates. The leading error has order $\hbar t$, attained on $\Hb$
(Proposition~\ref{prop:g5}); the pointwise estimate on an arbitrary quotient carries a local multiplicity factor $m(z)^{1/2}$,
necessary on the modular surface (Theorem~\ref{cl:main}). Section~\ref{sec:changes} gives the precise comparison with \cite{S}.
The factorization, the approximations $u^{(K)}$, the analytic norms and the mollified flow are Schubert's.
Our contribution is the organization of the remainder and the treatment of the weights (Section~\ref{subsec:method}).

Accuracy of the approximation does not by itself imply classical behaviour of expectation values.
For suitable states attached to the cusp of the modular surface, Proposition~\ref{ev:mod} proves that mass escapes every compact
set throughout the part of the approximation window beyond a sufficient time of order
$\log(1/\hbar)\log\log(1/\hbar)$, whereas classical transport equidistributes.
The main approximation results and parts (i) and (ii) of Proposition~\ref{ev:mod} are formalized in Lean~4.
Appendix~\ref{app:lean} specifies the formal statements and the exclusions, including the cusp sharpness theorem and the cited
classical equidistribution result.

\subsection{Setting and notation}\label{sec:setup}
A Fuchsian group is a discrete subgroup of the orientation-preserving isometry group of $\Hb$.
Torsion is allowed: when elliptic points are present, $M=\Hb/\Gamma$ is understood as a hyperbolic orbifold.
We use the word \emph{quotient} to include both surfaces and these orbifolds, with the measure induced from a fundamental domain.
We use the half-plane coordinates of \cite[Section~2]{S}. Let $\Hb=\{x+iy:y>0\}$ with the metric
$y^{-2}(dx^2+dy^2)$, volume $d\nu=y^{-2}dx\,dy$, and Laplacian $\Delta=y^2(\partial_x^2+\partial_y^2)$; the
spectrum of $-\Delta$ on $L^2(\Hb)$ is $[1/4,\infty)$.
On $\Hb$, $-\Delta$ denotes the Friedrichs extension from $C_c^\infty(\Hb)$; on $M$, it denotes the
Friedrichs extension from periodized test functions. The propagator and the powers of the Laplacian on either space
are defined by spectral calculus; we write $\Delta_M$ when the distinction is needed. Put $X=y\partial_x$, $Y=y\partial_y$; then
$\Delta=X^2+Y^2-Y$. (In \cite{S} the letter $Y$ denotes the generator $i(y\partial_y-\tfrac12)$ of the transport; here it
denotes the vector field $y\partial_y$, except in the proof of Lemma~\ref{cl:asym}, where $Y$ is a height and the vector fields are
written $\mathsf X$, $\mathsf Y$.) The boundary point is $b=\infty$, with $\varphi=\varphi_b=-\log y$ \cite[(2.1)]{S}.
The transport $S(t)a(x,y)=\ue^{-t/2}a(x,\ue^t y)$ is unitary \cite[Lemma~2.1]{S}. On either $\Hb$ or $M$ let $\cU(t)=\ue^{i\hbar t\Delta/2}$, with the underlying space specified by its argument. The
dispersive part is the unitary operator
\begin{equation}\label{eq:factor}
V(t):=\ue^{it/(2\hbar)}S(t)^{*}\,\ue^{-i\varphi/\hbar}\,\cU(t)\,\ue^{i\varphi/\hbar},\quad\text{so that}\quad
\ue^{-i\varphi/\hbar}\,\cU(t)\,\ue^{i\varphi/\hbar}=\ue^{-it/(2\hbar)}S(t)V(t).
\end{equation}
By \cite[Prop.~2.2]{S} it is the propagator of $i\partial_tV=-\tfrac{\hbar}{2}\Delta(t)V$, $V(0)=I$, with
\[
\Delta(t)=S(t)^*\Delta S(t)=y^2(\partial_y^2+\ue^{-2t}\partial_x^2)=\ue^{-2t}X^2+Y^2-Y
\]
\cite[(2.9)]{S}; we use this equation only on amplitudes all of whose derivatives by $X$ and $Y$ are square integrable
(Proposition~\ref{prop:dd-id}).
For a Fuchsian group $\Gamma$ and $u$ on $\Hb$, $u_\Gamma=\sum_{\gamma\in\Gamma}u\circ\gamma^{-1}$.

The approximations are those of \cite[(1.20), (1.27)--(1.29)]{S}. For $k\ge1$ let
\[
P_k(t):=\int_{t>s_1>\dots>s_k>0}\Delta(s_1)\cdots\Delta(s_k)\,ds_1\cdots ds_k,\qquad P_0(t):=1,
\]
and for $K\in\mathbb N_0$
\[
a^{(K)}(t):=\sum_{k=0}^{K}\Bigl(\frac{i\hbar}2\Bigr)^{k}P_k(t)a,\qquad
u^{(K)}(t):=\ue^{-it/(2\hbar)}\bigl(S(t)a^{(K)}(t)\bigr)\ue^{i\varphi/\hbar}.
\]
So the leading approximation is $u^{(0)}(t)=\ue^{-it/(2\hbar)}(S(t)a)\ue^{i\varphi/\hbar}$. For a boundary point $b\ne\infty$ the
phase $\varphi_b$, the transport and the approximations are the images of those for $\infty$ under an isometry that fixes the
base point $x_0$ of the norms and maps $\infty$ to $b$ (Lemma~\ref{lem:base}).

Schubert's amplitude norm is \cite[Def.~1.1]{S}
\[
\norm{a}_{\alpha,\beta}:=\norm{\ue^{\beta\langle d\rangle}\ue^{\alpha\sqrt{-\Delta}}a}_{L^2},\qquad
\langle d\rangle=(1+d(x_0,\cdot)^2)^{1/2}-1 .
\]
Throughout, $\Lambda=\sqrt{-\Delta}$, $X_0=\cosh d(x_0,\cdot)$, $\tilde\psi=\log X_0$ and $W_\gamma=\ue^{\gamma\tilde\psi}$.

\subsection{Main results}\label{sec:repair}\label{subsec:pc-thm}
Our principal theorem uses Schubert's amplitude norm exactly as defined in \cite[Def.~1.1]{S}.
For brevity write $\pcNorm a:=\norm a_{\alpha,\beta}$, with the parameters understood from each statement.
We call this the \emph{printed class} to distinguish it from the reversed order of weights used in the proofs of \cite{S}.
The norm measures analytic regularity and spatial decay; its precise order matters, since the two orders define different classes
(Remark~\ref{pc:classes}). The following estimate controls the expansion order $K$ and the number $N$ of Laplacians simultaneously.

\begin{theorem}[long-time approximation]\label{pc:main}
Let $\Gamma$ be a Fuchsian group (cusps allowed), $\beta>\tfrac12$ and $\kappa>0$. There are constants $C,\delta>0$, depending on $\Gamma$, $\beta$,
$\kappa$ and the base point $x_0$,
such that for all $0<\hbar\le1$, all $\alpha\in(0,\tfrac12]$ with $\alpha^3\ge\kappa\hbar$, all $K,N\in\mathbb N_0$, all $b\in\partial\Hb$, all $a$ with
$\pcNorm{a}<\infty$, and all $0\le t\le\delta(\alpha/\hbar)^{1/2}$,
\[
  \norm{\Delta^N\big[\cU(t)\big(a\,\ue^{i\varphi_b/\hbar}\big)_\Gamma-u^{(K)}(t)_\Gamma\big]}_{L^2(M)}
  \le C^{N+K+1}N^{2N}K!\,(\alpha\hbar)^{-2N}\Big(\frac{\hbar t}{\alpha^2}\Big)^{K+1}\pcNorm{a},
\]
with $N^{2N}:=1$ at $N=0$.
\end{theorem}

The proof is in Section~\ref{subsec:pc-proof}. For fixed $\alpha>0$, the time window has order $\hbar^{-1/2}$.
The following corollaries give the leading error, the exponential estimate obtained by choosing $K$, and the pointwise bounds.

\begin{corollary}[leading approximation]\label{pc:thm13}
Under the hypotheses of Theorem~\ref{pc:main}, with $K=N=0$,
\[
  \norm{\cU(t)\big(a\,\ue^{i\varphi_b/\hbar}\big)_\Gamma-u^{(0)}(t)_\Gamma}_{L^2(M)}\le C\,\frac{\hbar t}{\alpha^2}\,\pcNorm{a},
  \qquad 0\le t\le\delta(\alpha/\hbar)^{1/2}.
\]
\end{corollary}

For fixed $\alpha$, this is an error of order $\hbar t$. Proposition~\ref{prop:g5} shows that this order is attained on $\Hb$.

\begin{corollary}[optimal truncation]\label{pc:cor15}
Under the hypotheses of Theorem~\ref{pc:main}, put $\delta_K:=C\hbar t/\alpha^2$ and let $\epsilon\in(0,1)$, with $\delta\le\kappa^{1/2}/\max(1,C)$.
For every integer $K$ with $(1-\epsilon)/\delta_K\le K\le(1-\epsilon/2)/\delta_K$,
\[
  \norm{\Delta^N\big[\cU(t)\big(a\,\ue^{i\varphi_b/\hbar}\big)_\Gamma-u^{(K)}(t)_\Gamma\big]}_{L^2(M)}
  \le C'\epsilon^{-1/2}C^{N+1}N^{2N}(\alpha\hbar)^{-2N}\ue^{-(1-\epsilon)\alpha^2/(C\hbar t)}\pcNorm{a},
\]
with $C'$ absolute. If the window contains no integer the statement is vacuous; at $t=0$ the error is $0$.
\end{corollary}

For pointwise estimates, let
\[
  m(z):=\#\{\gamma\in\Gamma:d(\tilde z,\gamma\tilde z)<2\},\qquad z\in M,
\]
where $\tilde z$ is any lift of $z$. This finite local multiplicity is independent of the lift (Lemma~\ref{pc:g43}).
It is bounded on quotients of positive injectivity radius and grows like the height deep in an embedded cusp.

\begin{corollary}[pointwise approximation on bounded geometry]\label{pc:cor16bg}
Under the hypotheses of Theorem~\ref{pc:main}, assume in addition that the injectivity radius of $M$ is at least $r_0>0$. Then for every $z\in M$,
\[
  \big|\Delta^N\big[\cU(t)\big(a\,\ue^{i\varphi_b/\hbar}\big)_\Gamma-u^{(K)}(t)_\Gamma\big](z)\big|
  \le C^{N+K+1}(2N)!\,K!\,(\alpha\hbar)^{-2N-4}\Big(\frac{\hbar t}{\alpha^2}\Big)^{K+1}\pcNorm{a},
\]
with $C$ depending also on $r_0$.
\end{corollary}

\begin{corollary}[pointwise approximation on an arbitrary quotient]\label{pc:cor16}
Under the hypotheses of Theorem~\ref{pc:main}, for every $z\in M$,
\begin{multline*}
  \big|\Delta^N\big[\cU(t)\big(a\,\ue^{i\varphi_b/\hbar}\big)_\Gamma-u^{(K)}(t)_\Gamma\big](z)\big|\\
  \le C\,m(z)^{1/2}\,C^{N+K+1}(2N)!\,K!\,(\alpha\hbar)^{-2N-4}\Big(\frac{\hbar t}{\alpha^2}\Big)^{K+1}\pcNorm{a}.
\end{multline*}
In both corollaries the error lies in $\operatorname{Dom}(\Delta_M^N)$, and the left side is the value of the continuous representative.
\end{corollary}

The exponent $1/2$ of $m(z)$ is necessary on the modular surface for a phase attached to the cusp
(Theorem~\ref{cl:main}). The proofs of the pointwise corollaries are in Section~\ref{subsec:pc-cor}.
The restriction $\alpha\le\tfrac12$ in Theorem~\ref{pc:main} can be removed as follows.

\begin{corollary}[arbitrary analytic radius]\label{pc:allalpha}
Let $\Gamma$ be a Fuchsian group, $\beta>\tfrac12$, $\kappa>0$ and $A>0$. There are constants $C,\delta>0$, depending on $\Gamma$, $\beta$,
$\kappa$, the base point $x_0$ and $A$, such that the estimate of Theorem~\ref{pc:main} holds for all $0<\hbar\le1$, all $\alpha\in(0,A]$ with
$\alpha^3\ge\kappa\hbar$, all $K,N\in\mathbb N_0$, all $b\in\partial\Hb$, all $a$ with $\pcNorm a<\infty$ and all
$0\le t\le\delta(\alpha/\hbar)^{1/2}$; here $\pcNorm a$ is the printed norm at the radius $\alpha$. Corollaries~\ref{pc:thm13},
\ref{pc:cor15}, \ref{pc:cor16bg} and~\ref{pc:cor16} hold in the same way for $\alpha\in(0,A]$, with constants depending also on $A$.
\end{corollary}

The proof is in Section~\ref{subsec:pc-all}. Thus the approximation theorem and its corollaries cover every fixed analytic radius.

\medskip\noindent\emph{Companion results and sharpness.}
The reversed-order norm
\[
  \norm{a}_{\mathrm S}:=\norm{\ue^{\alpha\Lambda}\ue^{\beta\langle d\rangle}a}_{L^2(\Hb)}
\]
defines the \emph{swapped class}. Theorems~\ref{thm:main} and~\ref{thm:dd14} and Corollaries~\ref{cor:dd15} and~\ref{cor:dd16}
in Appendix~\ref{sec:dd} give its approximation estimates; their proof supplies the analytic estimates used in Section~\ref{sec:printed}.
Schubert's example $\hbar^{-\eta/2}\ue^{-\hbar^{-\eta}\langle d\rangle}$ belongs to this class for $0<\eta<2/3$
(Lemma~\ref{lem:ex} and the remark following it).
Section~\ref{sec:g5} establishes the sharpness of the leading error and the cusp factor. It also computes the mean of the approximation:
on a quotient of finite area, a nonzero initial mean gives a lower bound on the error of every fixed-order approximation at large times
(Proposition~\ref{prop:mean}). Section~\ref{sec:expect} distinguishes approximation accuracy from classical transport of expectation values.

\subsection{Method of proof}\label{subsec:method}
By \eqref{eq:factor} the error of $u^{(K)}$ is governed by the dispersive part $V(t)$, and $u^{(K)}$ corresponds to the Dyson
expansion of $V(t)$ to order $K$.

\emph{An exact remainder.} We write the remainder of the Dyson expansion with the unitary propagator on the outside
(Proposition~\ref{prop:dd-id}). The error is then $(\hbar/2)^{K+1}$ times an integral over the simplex
$t>s_1>\dots>s_{K+1}>0$ of a periodized propagator applied to
$b'(s)=\Delta_\varphi(s_1)^N\Delta(s_1)\cdots\Delta(s_{K+1})a$, a fixed polynomial in the vector fields $X,Y$ applied
to the amplitude. No resummation of the series is needed, and all derivatives fall on $a$. Analyticity of $a$ in a strip of
width $\alpha$ bounds a word of length $k$ in $X,Y$ applied to $a$ by $C^kk!\,\alpha^{-k}$ (Cauchy estimates,
Lemmas~\ref{lem:fc}--\ref{lem:wb}), and the integration over the simplex then gives the factor
$C^{N+K+1}N^{2N}K!\,(\alpha\hbar)^{-2N}(\hbar t/\alpha^2)^{K+1}$.

\emph{The periodized flow.} It remains to bound the periodized propagator applied to $b'$ in $L^2(M)$, uniformly in $s_1\le t$.
The periodization estimate used here (Lemma~\ref{lem:t2}) needs a weight $\ue^{sd}$ with $s>\tfrac12$, that is, above the
critical weight $\ue^{d/2}$, and for transported functions its bound grows like $\ue^{st}$. It is not known whether the flow
$V$ preserves weighted analytic norms on such time scales; \cite[Conjecture~1]{S} asserts this up to times of order $1/\hbar$. As in \cite{S} we compare $V$ with the flow
$V_\varepsilon$ of the generator mollified at scale $\varepsilon$; the estimate for $V_\varepsilon$ is the step whose proof
in \cite{S} has the gap (Section~\ref{sec:changes}). The mollified flow preserves analytic norms, with a loss of radius, and
power weights $X_0^\gamma$, both up to times of order $\varepsilon/\hbar$ (Propositions~\ref{prop:a0} and~\ref{prop:ws}); the first
is proved by an energy estimate with a clipped weight. Without a weight, the defect of the mollifier is exponentially small, of size
$\ue^{-a/(8\varepsilon)}$, on vectors analytic in a strip of width $a$ (Lemma~\ref{lem:unw}); with the
weight $X_0^\beta$ the bound we use is a power of $\varepsilon^{-1}$ (Lemma~\ref{lem:path}). Interpolation gives smallness at an
intermediate weight, still above the critical one, and this pays for the growth of the periodization bound. With
$\varepsilon=(\alpha\hbar)^{1/2}/2$ all requirements are met for $t\le\delta(\alpha/\hbar)^{1/2}$.

\emph{The original amplitude norm.} An amplitude of the printed class is $a=\ue^{-\alpha\Lambda}f$,
where $f=\ue^{\alpha\Lambda}a$ has the exponential weighted bound $\norm{\ue^{\beta\langle d\rangle}f}=\pcNorm a$.
Poisson smoothing need not preserve a spatial weight above the critical exponent, so this bound on $f$ does not give the
corresponding bound on $a$ (Remark~\ref{pc:classes}). The argument therefore uses the representation
$a=\ue^{-\sigma\Lambda_m}\tilde f$, with $\tilde f$ controlled at the critical weight with a logarithmic factor
(Proposition~\ref{os:data}). We interpolate the linear maps producing the vectors of the proof between unweighted $L^2$
and a power-weighted space (Lemma~\ref{os:transfer}). The smoothing operators are bounded at the critical weight by a
multiplier theorem proved using finite propagation speed (Theorem~\ref{os:cw}); periodization is controlled at the same weight
by Proposition~\ref{os:p1}. These estimates yield Theorem~\ref{pc:main}.

\subsection{Context and related work}\label{subsec:context}
Rigorous semiclassical approximations for
the Schr\"odinger flow are in general controlled only up to the Ehrenfest
time $T\sim\lambda^{-1}\log(1/\hbar)$, where $\lambda$ is a Lyapunov exponent.
For coherent states this is due to Combescure and Robert \cite{CR97} and
Hagedorn and Joye \cite{HJ00}; for Egorov's theorem to Bambusi, Graffi and
Paul \cite{BGP99} and Bouzouina and Robert \cite{BR02}. For quantized
hyperbolic toral automorphisms, Bonechi and De Bi\`evre \cite{BdB00} showed that
coherent states equidistribute on the logarithmic time scale. For Lagrangian
states on Anosov manifolds, Schubert proved equidistribution of expectation
values up to the Ehrenfest time \cite{Sch05}, and with Vallejos and Toscano
described how localized states become Lagrangian states on that time scale
\cite{STV12}. Estimates for the propagator on time scales that are multiples of the
Ehrenfest time are also the main tool in the work of Anantharaman \cite{Ana08} and of
Anantharaman and Nonnenmacher \cite{AN07} on the entropy of semiclassical measures,
and, on compact hyperbolic surfaces, in the proof of Dyatlov and Jin \cite{DJ18} that
semiclassical measures have full support. All of these use propagation estimates on time
scales that are logarithmic in $\hbar$. On such scales, Faure \cite{Fau07} extended the semiclassical
trace formula for a nonlinear hyperbolic map of the torus to times $C$ times the Ehrenfest
time, for every constant $C$. On compact hyperbolic surfaces, Han \cite{Han26} recently obtained
Egorov estimates for Laplace eigenfunctions at all forward logarithmic times; for backward times
he introduced a Schr\"odinger group along the stable horocycles.

For the accuracy of semiclassical approximations of a propagated state beyond the Ehrenfest
time, \cite{S} points to numerical evidence. In the account given there, Tomsovic and Heller
\cite{TH91,TH93} studied the accuracy of semiclassical propagation for the stadium
billiard, found no breakdown at the Ehrenfest time, and argued that the approximation
should remain accurate up to times of order $1/\hbar$ when the classical system has no
singularities. The obstacle to a proof named in \cite{S} is that the approximation is a sum of
oscillating terms whose number grows exponentially with time, and the error terms
are of the same nature. The preprint \cite{S} aimed to control the error far
beyond the Ehrenfest time on
surfaces of constant negative curvature, up to $T\sim\hbar^{-1/2}$, for
Lagrangian states attached to an unstable manifold. The present paper establishes this $L^2$ approximation on arbitrary Fuchsian quotients for the original amplitude norm,
with explicit factorial dependence on the expansion and derivative orders. Section~\ref{sec:changes} records the changes to
the estimates stated in \cite{S}.

The theorems control the error of the approximation; they do not describe the
approximation itself, which is again a sum over exponentially many sheets of the
cover. In the picture that goes back to Berry \cite{Ber77} (for propagated waves see \cite{BB79}, as cited in \cite{S}), waves
propagated in a classically chaotic system behave at large times like random
superpositions of plane waves. Ingremeau and Rivera \cite{IR22} proved that on compact
manifolds of negative curvature, for a wide family of phases, a Lagrangian state evolved
for a long time behaves like a random Gaussian field; we refer to \cite{IR22} for the
precise statement and its time scales. For the equation $i\partial_tu=-\tfrac12\Delta u$ on compact
manifolds with Anosov geodesic flow, Anantharaman and Rivi\`ere \cite{AR12} proved dispersion
and controllability results by means of semiclassical measures; a time of order one for
that equation corresponds to a time of order $1/\hbar$ in the scaling used here. These results
do not give an approximation of the solution.

The approximation here is the periodization $u^{(K)}(t)_\Gamma$ of an explicit
approximation on the universal cover; it is not classical transport of the Lagrangian
measure. The distinction matters on surfaces with cusps. Brooks
\cite{Bro22} showed that, for Lagrangian states attached
to unstable horocycles on the modular surface, expectation values of the
wave-propagated states depart from classical transport along the geodesic flow
at logarithmic times, because the states escape into the cusp. The geometry
behind such an escape is visible in the periodization. In the normalization
$b=\infty$ the classically transported state lies near height $y\asymp\ue^{-t}$.
For a modular element $\gamma$ with lower-left entry $c\ne0$ and lower-right
entry $d$, $\operatorname{Im}\gamma z=y/((cx+d)^2+c^2y^2)$, so the part of the
state within $O(y)$ of the rational $-d/c$ is mapped to height $\asymp\ue^t/c^2$,
which is at least $Y\ge1$ when $c\lesssim Y^{-1/2}\ue^{t/2}$. The periodization
$u^{(K)}(t)_\Gamma$ sums over all these sheets, and beyond the Ehrenfest time the
interference between them, which the classical prediction ignores, governs the
fine behaviour. This geometry is compatible with Brooks's result; it is not a
proof of escape, since it does not compare the mass or phase of the
contributing sheets. Brooks's initial states carry a fixed smooth cutoff in $y$
and so are not analytic amplitudes of the class of Theorem~\ref{pc:main}, and he uses the wave group; we do not compare the
two results directly. Proposition~\ref{ev:mod} proves the corresponding statement for the Schr\"odinger flow and for
amplitudes of the printed class whose periodization along the cusp does not depend on $x$.

\subsection{Relation to the preprint}\label{sec:changes}
We list how the statements above relate to those of \cite{S}, equations (1.15)--(1.36) there.
\begin{itemize}
\item \emph{The gap.} In the proofs of \cite[Theorems~4.1 and~4.4]{S}, the steps (4.16)$\to$(4.17), (4.55)$\to$(4.56) and
(4.62)$\to$(4.63) infer $\norm{\ue^{X}a}\le\norm{\ue^{Y}a}$ from an operator inequality $X\le Y$ between non-commuting
self-adjoint operators. The exponential is not operator monotone: with $X=\mathrm{diag}(0,1)$,
$H=\begin{pmatrix}1&1\\1&1\end{pmatrix}\ge0$ and $v=(-3,1)$, the function $g(s)=\norm{\ue^{X+sH}v}^2$ has
$g'(0)=4(6-\ue^2)<0$. In addition, \cite[Def.~1.1]{S} and the proofs of \cite[Prop.~3.5, Lemma~3.11]{S} take the two weights
of the norm in different orders. The proof of \cite[Lemma~3.11]{S} passes from one order to the other by
\cite[Proposition~7.1, (7.2)]{S}, which states that the two norms are equivalent; the sketch of its proof rests on
\cite[Theorem~7.2]{S}, which is stated for small parameters $|\alpha|,|\beta|\le c$. For $\beta>\tfrac12$, the range the theorems
need, the two norms are not equivalent (Remark~\ref{pc:classes}).
\item \emph{Amplitude class.} Theorem~\ref{pc:main} uses the norm of \cite[Def.~1.1]{S} for $\alpha\le\tfrac12$;
Corollary~\ref{pc:allalpha} extends it to $\alpha\le A$, with constants depending on $A$.
Theorems~\ref{thm:main} and~\ref{thm:dd14} and their corollaries in Appendix~\ref{sec:dd} treat the reversed order used in Schubert's proofs.
\item \emph{Theorem 1.3.} \cite[(1.18)--(1.19)]{S} states the error $C\hbar\norm a_{\alpha,\beta}$ for $0\le t\le\delta\hbar^{-1/2}$.
The error is of order $\hbar t$ (Corollary~\ref{pc:thm13}), and this order is attained on $\Hb$ (Proposition~\ref{prop:g5}). This is also
how the introduction of \cite{S} describes it: the dispersion ``takes place on a scale of order $\hbar t$, and this is
responsible for the error term in (1.18)'', and the expression \cite[(1.26)]{S} ``is of order $\hbar t$''. The window
and the generality in $\Gamma$ are unchanged.
\item \emph{Theorem 1.4.} \cite[(1.30)--(1.31)]{S} states the bound
\[C^{N+K+1}N^{2N}K!\,(\alpha\hbar)^{-2N-4}(\hbar t/\alpha^2)^K\norm a_{\alpha,\beta}\]
for $\alpha^3\ge\beta^2\hbar$ and $0\le t\le\delta\alpha^{3/2}\hbar^{-1/2}$. Theorem~\ref{pc:main} retains the dependence $C^{N+K+1}N^{2N}K!$ on the orders, with the condition
$\alpha^3\ge\kappa\hbar$ for any fixed $\kappa>0$ ($\kappa=\beta^2$ is the printed one), the longer window
$t\le\delta(\alpha/\hbar)^{1/2}$, and a bound that is smaller than the printed one by the factor $(\alpha\hbar)^4\hbar t/\alpha^2$,
which is at most $1$ on the window of \cite{S}.
Theorem~\ref{pc:main} is stated for $\alpha\le\tfrac12$; Corollary~\ref{pc:allalpha} gives every bounded range of radii, with constants depending on its upper bound.
\item \emph{Corollary 1.5.} The window for $K$ printed in \cite[(1.32)]{S} is replaced by the multiplicative window
$K\delta_K\in[1-\epsilon,1-\epsilon/2]$, which is the one used in the proof in \cite{S}. The prefactor is $(\alpha\hbar)^{-2N}$
instead of $(\alpha\hbar)^{-2N-4}$.
\item \emph{Corollary 1.6.} \cite[(1.34)--(1.36)]{S} states the pointwise bound
\[C^{N+K+1}(2N)!\,K!\,(\alpha\hbar)^{-2N-8}(\hbar t)^K\norm a_{\alpha,\beta}\]
for every $\Gamma$. Corollary~\ref{pc:cor16bg} assumes a positive
injectivity radius and has the prefactor $(\alpha\hbar)^{-2N-4}$ and the time factor $(\hbar t/\alpha^2)^{K+1}$. For fixed
$\alpha$ the factor $\alpha^{-2K}$ can be put into $C^K$, and the printed time factor follows; with constants uniform in
$\alpha$ the factor $\hbar t/\alpha^2$ is what Theorem~\ref{pc:main} gives. On surfaces with cusps the bound holds with the factor
$m(z)^{1/2}$ (Corollary~\ref{pc:cor16}). Below cusp height $\asymp(\alpha\hbar)^{-8}$ the printed prefactor
$(\alpha\hbar)^{-2N-8}$ absorbs this factor (Remark~\ref{pc:remcusp}); above that height it cannot be removed on the modular
surface (Theorem~\ref{cl:main}, Remark~\ref{cl:reading}). In \cite[(1.36)]{S}, $\psi^{(K)}$ should read $u^{(K)}$.
\end{itemize}

\subsection{Organization}
Section~\ref{sec:printed} proves the main approximation theorem and its corollaries, including the extension to arbitrary radii.
Section~\ref{sec:g5} treats sharpness and limitations: the leading error on $\Hb$, the mean on quotients of finite area, and the
necessity of the cusp factor on the modular surface. Section~\ref{sec:expect} treats expectation values and escape of mass;
Section~\ref{sec:outlook} collects open problems.
The appendices supply the analytic estimates: the mollified flow (Appendix~\ref{app:kl}), multiplication by weights
(Appendix~\ref{app:m}), weighted propagation (Appendix~\ref{app:ws}), the companion results for the swapped class and their proof
(Appendix~\ref{sec:dd}), periodization (Appendix~\ref{sec:t2}), and the estimates at the critical weight used in the main proof
(Appendix~\ref{app:o1st}). Appendix~\ref{app:lean} describes the formalization.

% Main-text section: Schubert's printed class.  Input file; no preamble.
% Macros introduced here (prefix pc): \pcXzero = X_0, \pcLog = l, \pcw = omega, \pcNorm = the printed norm.
\newcommand{\pcXzero}{X_0}
\newcommand{\pcLog}{\ell}
\newcommand{\pcw}{\omega}
\newcommand{\pcNorm}[1]{\norm{#1}_{\mathrm P}}

\section{Schubert's printed class}\label{sec:printed}

Throughout, $\Gamma$ is a fixed Fuchsian group, $M=\Hb/\Gamma$, $x_0\in\Hb$ is the base point of the norms below, and
$d=d(x_0,\cdot)$. Schubert's printed norm is \cite[(1.15)]{S}
\[
  \pcNorm{a}:=\norm{\ue^{\beta\langle d\rangle}\ue^{\alpha\sqrt{-\Delta}}a}_{L^2(\Hb)},\qquad
  \langle d\rangle=(1+d^2)^{1/2}-1,\qquad \beta>\tfrac12 .
\]

\subsection{The critical weight}\label{subsec:pc-weights}
Let $\pcXzero:=\cosh d$, $\psi:=\log\pcXzero\ge0$, $\pcLog:=\log(2\pcXzero)/\log2\ge1$ and $\Lambda:=\sqrt{-\Delta}$. Fix $c':=\tfrac34$ (any exponent in
$(\tfrac12,1)$ would do) and put
\[
  \pcw:=\pcXzero^{1/2}\pcLog^{\,c'} .
\]
This is the periodization weight: $\sum_\gamma\pcw(\gamma z)^{-2}$ is bounded for every Fuchsian group (Proposition~\ref{os:p1}); for $\pcXzero^{1/2}$ alone
the sum diverges when $M$ has finite area. Since
$\psi\le d\le\langle d\rangle+1$, $\pcw\le C_\beta\ue^{\beta\langle d\rangle}$ for every $\beta>\tfrac12$. So an amplitude of the printed class is
$a=\ue^{-\alpha\Lambda}f$ with $\norm{\pcw f}\le C_\beta\pcNorm a$: its data $f$ lie in the weighted space of the critical weight, and no weight above the
critical one is available. In the swapped class the mollifier part of the error is small because one vector is small without a weight and bounded with
the weight $\pcXzero^\beta$, $\beta>\tfrac12$. Here the same interpolation is carried out for the linear maps that produce these vectors from the data,
between $L^2(\Hb)$ and $L^2(\pcXzero^{2}d\nu)$ (Lemma~\ref{os:transfer}), after the data have been written as $a=\ue^{-\sigma\Lambda_m}\tilde f$ with an even
smoothing operator and $\norm{\pcw\tilde f}\le C\pcNorm a$ (Proposition~\ref{os:data}).

\subsection{The theorem}\label{subsec:pc-thm}
\begin{theorem}[printed class]\label{pc:main}
Let $\Gamma$ be a Fuchsian group (cusps allowed), $\beta>\tfrac12$ and $\kappa>0$. There are constants $C,\delta>0$, depending on $\Gamma$, $\beta$,
$\kappa$ and the base point $x_0$,
such that for all $0<\hbar\le1$, all $\alpha\in(0,\tfrac12]$ with $\alpha^3\ge\kappa\hbar$, all $K,N\in\mathbb N_0$, all $b\in\partial\Hb$, all $a$ with
$\pcNorm{a}<\infty$, and all $0\le t\le\delta(\alpha/\hbar)^{1/2}$,
\[
  \norm{\Delta^N\big[\cU(t)\big(a\,\ue^{i\varphi_b/\hbar}\big)_\Gamma-u^{(K)}(t)_\Gamma\big]}_{L^2(M)}
  \le C^{N+K+1}N^{2N}K!\,(\alpha\hbar)^{-2N}\Big(\frac{\hbar t}{\alpha^2}\Big)^{K+1}\pcNorm{a},
\]
with $N^{2N}:=1$ at $N=0$.
\end{theorem}

\begin{remark}
On the window of \cite[(1.31)]{S}, $t\le\delta\alpha^{3/2}\hbar^{-1/2}$, the bound of Theorem~\ref{pc:main} implies \cite[(1.30)]{S} with the printed
norm, because $(\alpha\hbar)^{4}\hbar t/\alpha^2\le1$ there. The swapped-class version is Theorem~\ref{thm:dd-appx} of Appendix~\ref{sec:dd}; the proof
below uses its frame and its lemmas. The proof gives $\delta=\min(\tfrac1{64},c_A/6)/\sqrt6$, with $c_A$ the constant of Proposition~\ref{prop:a0}
for $\tau_0=1$; this $\delta$ does not depend on $\Gamma$, $\beta$ or $\kappa$.
\end{remark}

\begin{remark}[the two classes, and the range of $\alpha$]\label{pc:classes}
Let $\beta>\tfrac12$ and $0\le\alpha'<\alpha<\pi/2$, and put $h:=\ue^{-\beta\langle d\rangle}\mathbf 1_{B(x_0,1)}$ and $a:=\ue^{-\alpha\Lambda}h$, so that
$\pcNorm a=\norm{\mathbf 1_{B(x_0,1)}}<\infty$ at $(\alpha,\beta)$. For $0<\tau<\pi/2$ the function $\ue^{-\tau\Lambda}h$ is the convolution of $h$ with the
radial kernel $k_\tau$ of Lemma~\ref{os:poisson}, and $d(z,w)\in[d(z)-1,d(z)+1]$ for $w\in B(x_0,1)$. So by \eqref{os:lb}, for $d(z)\ge2$,
\[
  (\ue^{-\tau\Lambda}h)(z)\ge c\,(d(z)+1)^{-3/2}\ue^{-d(z)},
\]
with $c>0$ depending on $\tau$ and $\beta$. Since $\langle d\rangle\ge d-1$, for every $\beta'>\tfrac12$ the square of
$\ue^{\beta'\langle d\rangle}\ue^{-\tau\Lambda}h$ is at least $c'\ue^{(2\beta'-2)d}(d+1)^{-3}$ where $d\ge2$, which is not integrable against the area
element $\sinh\rho\,d\rho\,d\theta$; so this function is not in $L^2(\Hb)$. Two consequences follow.
With $\tau=\alpha$: the amplitude $a$ of the printed class lies in no swapped class with an exponent $\beta'>\tfrac12$, because
$\norm{\ue^{\alpha''\Lambda}\ue^{\beta'\langle d\rangle}a}\ge\norm{\ue^{\beta'\langle d\rangle}a}$ for $\alpha''\ge0$.
With $\tau=\alpha-\alpha'$ and $\beta'=\beta$: the printed norm of $a$ at $(\alpha',\beta)$ is infinite, since
$\ue^{\alpha'\Lambda}a=\ue^{-(\alpha-\alpha')\Lambda}h$. So the printed norms are not monotone in $\alpha$, and Theorem~\ref{pc:main}, which is stated for
$\alpha\le\tfrac12$, does not cover larger radii by a comparison of norms; in the swapped class it does (Section~\ref{sec:repair}).
\end{remark}

\begin{corollary}[Theorem 1.3 for the printed class]\label{pc:thm13}
Under the hypotheses of Theorem~\ref{pc:main}, with $K=N=0$,
\[
  \norm{\cU(t)\big(a\,\ue^{i\varphi_b/\hbar}\big)_\Gamma-u^{(0)}(t)_\Gamma}_{L^2(M)}\le C\,\frac{\hbar t}{\alpha^2}\,\pcNorm{a},
  \qquad 0\le t\le\delta(\alpha/\hbar)^{1/2}.
\]
\end{corollary}

\subsection{Proof of Theorem~\ref{pc:main}}\label{subsec:pc-proof}
We use the notation and the statements of Appendices~\ref{app:kl}, \ref{app:ws}, \ref{sec:dd} and~\ref{app:o1st}; all lemmas are used with $\tau_0=1$.

\medskip\noindent\emph{Step 0 (reductions and parameters).} By Lemma~\ref{lem:base} we may take $x_0=i$ and $b=\infty$: the group enters below only through
the constant $C(\Gamma,c')$ of Proposition~\ref{os:p1}, which is invariant under the conjugation of Lemma~\ref{lem:base}. Since
$a\in\operatorname{Dom}(\ue^{\alpha\Lambda})$, $a$ is $C^\infty$ by Lemma~\ref{lem:lift}. Put
\[
  m:=3,\qquad\sigma:=\alpha/6,\qquad E:=\ue^{-\sigma\Lambda_m},\qquad\varepsilon:=(\sigma\hbar)^{1/2}/2,\qquad x:=(\sigma/\hbar)^{1/2},
\]
so that $\sigma\le\tfrac1{12}$, $\varepsilon\le\tfrac12$, $\varepsilon/\hbar=x/2$, $\sigma/\varepsilon=2x$ and $\sigma^3\ge\kappa'\hbar$ with $\kappa':=\kappa/216$. For
$\mathbf s\in\Sigma_{K+1}(t)$ let $B'_{\mathbf s}:=\Delta_\varphi(s_1)^N\Delta(s_1)\cdots\Delta(s_{K+1})$, so that $b'(\mathbf s)=B'_{\mathbf s}a$ in (D4) of
Proposition~\ref{prop:dd-id}. For $n\ge0$ put
\[
  F_n:=C_5^{N+K+n+1}(2N)!\,(n+2K+2)!\,\sigma^{-(n+2K+2)}(\sigma\hbar)^{-2N},
\]
with a constant $C_5\ge1$ fixed in Step~2(a).

\medskip\noindent\emph{Step 1 (the data).} By Proposition~\ref{os:data} with $c=c'$ (here $\sigma m=\alpha/2\le\tfrac12$), there is $\tilde f$ with
\[
  a=E\tilde f,\qquad\norm{\pcw\tilde f}\le C\,\pcNorm a .
\]

\medskip\noindent\emph{Step 2 (three maps, from the data to the vectors of the proof).} Each of the following maps is linear on $L^2(\Hb)$, and is bounded
there and at the weight $\pcXzero$ on functions supported in one shell. Lemma~\ref{os:transfer} then bounds it at the weight $\pcw$.

\smallskip\noindent(a) \emph{Words.} For a word $Z_J$ of length $n$ let $T_ah:=Z_JB'_{\mathbf s}Eh$. Expand $B'_{\mathbf s}$ as in the proof of
Lemma~\ref{lem:wb}. Each term is a coefficient times a word of length at most $L$ applied to $Eh$, and by Lemma~\ref{os:words}(b), with $\gamma=1$ and
$m=3$ (so $2(\gamma+\tfrac12)^2<m^2$), it is bounded at the weight $\pcXzero$ by $2(L+1)L!\,(QC_0/\sigma)^L\norm{\pcXzero h}$. This is the bound of Lemma~\ref{lem:sc}
with $(2,QC_0,\sigma)$ in place of $(B_\beta,C_3,\alpha)$, so the count in the proof of Lemma~\ref{lem:wb} gives, for a suitable $C_5\ge C_4$,
\[
  \norm{\pcXzero Z_JB'_{\mathbf s}Eu}\le F_n\norm{\pcXzero u},\qquad\norm{Z_JB'_{\mathbf s}Eu}\le F_n\norm u ;
\]
the second bound is the same argument without a weight. Lemma~\ref{os:transfer} with $M_0=M_1=F_n$ gives
\begin{equation}\label{pc:wbw}
  \norm{\pcw Z_JB'_{\mathbf s}Eh}\le K\,F_n\,\norm{\pcw h}.
\end{equation}
In particular every word applied to $a=E\tilde f$ lies in $L^2_{\pcw}$.

\smallskip\noindent(b) \emph{Flow.} For $s',s''\ge0$ with $\hbar|s'-s''|\le\varepsilon$, the map $V_\varepsilon(s',s'')$ has norm $1$ on $L^2$ and, by
Proposition~\ref{prop:ws}, norm at most $\ue^{G_1/2}$ at the weight $\pcXzero$. So
\begin{equation}\label{pc:flow}
  \norm{\pcw V_\varepsilon(s',s'')h}\le K\ue^{G_1/2}\norm{\pcw h}.
\end{equation}

\smallskip\noindent(c) \emph{Defect.} For $0\le t'\le s_1\le t$ let
\[
  T_ch:=\bigl(\Delta(t')-\Delta_\varepsilon(t')\bigr)V_\varepsilon(t',s_1)B'_{\mathbf s}Eh .
\]
\emph{Unweighted bound.} By Lemma~\ref{os:words}, $\norm{\ue^{\sigma\Lambda}Eh}\le\norm h$. By (WB-u) of Lemma~\ref{lem:wb}, with $\sigma$ in the role of
$\alpha$ and the radii $a_1:=\tfrac34\sigma<\sigma$, by Proposition~\ref{prop:a0} from $a_1$ to $a_2:=\tfrac12\sigma$ (valid for
$t\le\tfrac13c_A\varepsilon/\hbar$), and by Lemma~\ref{lem:unw} with $a=a_2$, where $a_2/(8\varepsilon)=x/8$,
\[
  \norm{T_ch}\le M_0\norm h,\qquad M_0:=C\sigma^{-4}\ue^{-x/8}\,C_4^{N+K+1}(2N)!\,(2K+2)!\,\sigma^{-(2K+2)}(\sigma\hbar)^{-2N}.
\]
\emph{Weighted bound.} Let $u$ be supported in one shell and assume $\hbar s_1\le\varepsilon$. By (a), $\pcXzero ZB'_{\mathbf s}Eu\in L^2$ for $Z\in\{1,X,Y,X^2,Y^2\}$, and
Lemma~\ref{lem:path}(b) with $\gamma=1$ and $b'=B'_{\mathbf s}Eu$ gives
\[
  \norm{\pcXzero T_cu}\le M_1\norm{\pcXzero u},\qquad M_1:=C\varepsilon^{-3}F_2 .
\]
Lemma~\ref{os:transfer} requires $M_0\le M_1$; we apply it with $M_0+M_1$ in place of $M_1$. Since $M_0\le C\sigma^{-4}\ue^{-x/8}F_2$,
\begin{equation}\label{pc:defect}
  \norm{\pcw T_ch}\le K\bigl(M_0+M_0^{1/4}M_1^{3/4}\bigr)\norm{\pcw h}\le C\,\varepsilon^{-3}\sigma^{-4}\,\ue^{-x/32}\,F_2\,\norm{\pcw h}.
\end{equation}

\medskip\noindent\emph{Step 3 (Duhamel on $M$ for data of bounded support).} Let $h\in L^2$ be supported in $\{\psi\le R_0\}$ for some $R_0$, and let
$b'_h:=B'_{\mathbf s}Eh$, $g(t'):=V_\varepsilon(t',s_1)b'_h$ for $0\le t'\le s_1$, $w_h:=g(0)$ and $d_h(t'):=T_ch$. Then $\pcXzero h\in L^2$ and, by (a), all
words applied to $b'_h$ lie in $L^2(\pcXzero^2d\nu)$. The weight $\pcXzero^{3/4}$ satisfies (W1) and (W2) (Lemma~\ref{lem:t2}), and by
Lemma~\ref{lem:path}(c) with $\gamma=\tfrac34$ the path $g$ satisfies the hypotheses of Lemma~\ref{lem:surf-duh} for it. As in the proof of
Theorem~\ref{thm:dd-appx},
\[
  \mathcal V_{s_1}b'_h=\mathcal V_0w_h-\tfrac{i\hbar}2\int_0^{s_1}\mathcal V_{t'}d_h(t')\,dt' .
\]
By Remark~\ref{os:indep} each term may be estimated through the weight $\pcw$. By Proposition~\ref{os:p1}, \eqref{pc:flow} and \eqref{pc:wbw} with $n=0$,
\[
  \norm{\mathcal V_0w_h}\le C\norm{\pcw V_\varepsilon(0,s_1)b'_h}\le C\norm{\pcw b'_h}\le C\,F_0\,\norm{\pcw h},
\]
and by Proposition~\ref{os:p1} and \eqref{pc:defect},
\[
  \norm{\mathcal V_{t'}d_h(t')}\le C\ue^{t'/2}(1+t'/\log2)^{c'}\norm{\pcw d_h(t')}\le C\,\ue^{t'}\varepsilon^{-3}\sigma^{-4}\ue^{-x/32}F_2\norm{\pcw h}.
\]
Hence $\norm{\mathcal V_{s_1}b'_h}\le C(F_0+R\,F_2)\norm{\pcw h}$ with
\[
  R:=\hbar t\,\ue^{t}\,\varepsilon^{-3}\sigma^{-4}\,\ue^{-x/32}.
\]
\emph{Absorption.} Let $t\le\delta'x$ with $\delta'\le\tfrac1{64}$, so $\ue^t\le\ue^{x/64}$. Since $\sigma^3\ge\kappa'\hbar$ and $\sigma\le1$,
\[
  \sigma^{-1}\le x/\sqrt{\kappa'},\qquad\varepsilon^{-1}=2x/\sigma\le2x^2/\sqrt{\kappa'},\qquad\hbar t\le\delta'\sigma/x\le\delta'/x .
\]
So $R\le C(\kappa)\,x^{9}\ue^{-x/64}\le C(\kappa)\,x^{-2}$, and $x^{-2}=\hbar/\sigma\le\sigma^2/\kappa'$. Since $F_0\le\sigma^2F_2$,
\begin{equation}\label{pc:star}
  \norm{\mathcal V_{s_1}B'_{\mathbf s}Eh}\le C\,\sigma^2F_2\,\norm{\pcw h}.
\end{equation}

\medskip\noindent\emph{Step 4 (density).} The map $h\mapsto\mathcal V_{s_1}B'_{\mathbf s}Eh$ is linear and bounded from $L^2(\pcw^2d\nu)$ to $L^2(M)$, by
Proposition~\ref{os:p1} and \eqref{pc:wbw} with $n=0$. The functions of bounded $\psi$-support are dense in $L^2(\pcw^2d\nu)$. So \eqref{pc:star} holds for
every $h\in L^2(\pcw^2d\nu)$, in particular for $\tilde f$:
\[
  \norm{\mathcal V_{s_1}b'(\mathbf s)}\le C\,C_5^{N+K+3}(2N)!\,(2K+4)!\,\sigma^{-(2K+2)}(\sigma\hbar)^{-2N}\norm{\pcw\tilde f}.
\]

\medskip\noindent\emph{Step 5 (collecting).} By \eqref{pc:wbw}, every word applied to $a$ lies in $L^2_{\pcw}$, and $\pcw$ satisfies (W1), (W2). So (D4) of
Proposition~\ref{prop:dd-id} holds with the weight $\pcw$: the error is at most $(\hbar/2)^{K+1}$ times the integral of $\norm{\mathcal V_{s_1}b'(\mathbf s)}$ over
$\Sigma_{K+1}(t)$. As in the proof of Theorem~\ref{thm:dd-appx}, with $(2K+4)!/(K+1)!\le C^{K+1}K!$, $(2N)!\le4^NN^{2N}$ and
$\hbar^{K+1}t^{K+1}\sigma^{-(2K+2)}=(\hbar t/\sigma^2)^{K+1}$, Step~4 and Step~1 give the bound of the theorem with $\sigma$ in place of $\alpha$. Since
$\sigma=\alpha/6$, $(\sigma\hbar)^{-2N}(\hbar t/\sigma^2)^{K+1}=6^{2N+2K+2}(\alpha\hbar)^{-2N}(\hbar t/\alpha^2)^{K+1}$, and the factor is absorbed in $C^{N+K+1}$.

\medskip\noindent\emph{Time windows.} Steps~2(b), 2(c) and~3 use $\hbar s_1\le\varepsilon$, that is $t\le x/2$; Proposition~\ref{prop:a0} uses
$t\le\tfrac13c_A\varepsilon/\hbar=c_Ax/6$; the absorption uses $t\le x/64$. So the proof holds for $t\le\delta'x$ with $\delta':=\min(\tfrac1{64},c_A/6)$, and
$\delta'x=(\delta'/\sqrt6)(\alpha/\hbar)^{1/2}$. The hypothesis $\alpha\le\tfrac12$ ensures $\sigma<\alpha\le1$ in Proposition~\ref{os:data} and keeps all
radii below $\tau_0=1$. \qed

\subsection{Corollaries}\label{subsec:pc-cor}
\begin{corollary}[Corollary 1.5 for the printed class]\label{pc:cor15}
Under the hypotheses of Theorem~\ref{pc:main}, put $\delta_K:=C\hbar t/\alpha^2$ and let $\epsilon\in(0,1)$, with $\delta\le\kappa^{1/2}/\max(1,C)$.
For every integer $K$ with $(1-\epsilon)/\delta_K\le K\le(1-\epsilon/2)/\delta_K$,
\[
  \norm{\Delta^N\big[\cU(t)\big(a\,\ue^{i\varphi_b/\hbar}\big)_\Gamma-u^{(K)}(t)_\Gamma\big]}_{L^2(M)}
  \le C'\epsilon^{-1/2}C^{N+1}N^{2N}(\alpha\hbar)^{-2N}\ue^{-(1-\epsilon)\alpha^2/(C\hbar t)}\pcNorm{a},
\]
with $C'$ absolute. If the window contains no integer the statement is vacuous; at $t=0$ the error is $0$.
\end{corollary}
\begin{proof}
The proof of Corollary~\ref{cor:dd15-appx} uses only the form of the bound of Theorem~\ref{thm:dd-appx}:
$C^{N+K+1}N^{2N}K!(\alpha\hbar)^{-2N}(\hbar t/\alpha^2)^{K+1}$ times the norm of $a$, for every $K$. Theorem~\ref{pc:main} gives a bound of
exactly this form with $\pcNorm{a}$ in place of the swapped norm, so that proof applies word for word.
\end{proof}

\begin{lemma}[local Sobolev inequality on a Fuchsian quotient]\label{pc:g43}
Let $m(z):=\#\{\gamma\in\Gamma:\ d(\tilde z,\gamma\tilde z)<2\}$ for $z\in M$ with a lift $\tilde z$; $m$ is finite and does not depend on the lift.
There is an absolute constant $C$ such that every $u\in\operatorname{Dom}(\Delta_M)$ has a continuous representative, and for all $z\in M$,
\[
  |u(z)|\le C\,m(z)^{1/2}\big(\norm{\Delta u}_{L^2(M)}+\norm{u}_{L^2(M)}\big).
\]
\end{lemma}
\begin{proof}
\emph{Step 1 (the lift).} Let $\tilde u$ and $\tilde f$ be the $\Gamma$-invariant lifts of $u$ and of $f:=\Delta_Mu$. Then $\Delta\tilde u=\tilde f$ in the
sense of distributions on $\Hb$: for $\phi\in C_c^\infty(\Hb)$, unfolding gives
$\langle\tilde u,\Delta\phi\rangle_{\Hb}=\langle u,(\Delta\phi)_\Gamma\rangle_M=\langle u,\Delta_M\phi_\Gamma\rangle_M=\langle f,\phi_\Gamma\rangle_M=\langle\tilde f,\phi\rangle_{\Hb}$.

\emph{Step 2 (counting).} For a $\Gamma$-invariant $\tilde h$ and a fundamental domain $F$,
$\int_{B(\tilde z,1)}|\tilde h|^2d\nu=\int_Fn(w)|h(w)|^2d\nu(w)$, where $n(w):=\#\{\gamma:\gamma w\in B(\tilde z,1)\}$. If $\gamma_0w$ and $\gamma w$ both lie in
$B(\tilde z,1)$, then $d(\tilde z,\gamma\gamma_0^{-1}\tilde z)<2$; hence $n(w)\le m(z)$ and $\int_{B(\tilde z,1)}|\tilde h|^2\le m(z)\norm{h}^2_{L^2(M)}$. In
particular $\tilde u$ and $\tilde f$ are locally square integrable. The number $m(z)$ is finite because $\Gamma$ acts properly discontinuously.

\emph{Step 3 (interior regularity).} Let $v\in L^2(B(\tilde z,1))$ with $\Delta v\in L^2(B(\tilde z,1))$ in the sense of distributions. By interior
elliptic regularity, uniform in $\tilde z$ by the homogeneity of $\Hb$, $v\in H^2(B(\tilde z,\tfrac12))$ with
$\norm v_{H^2(B(\tilde z,1/2))}\le C(\norm{\Delta v}_{L^2(B(\tilde z,1))}+\norm v_{L^2(B(\tilde z,1))})$, and since $H^2\subset C^0$ in dimension two, $v$ is
continuous on $B(\tilde z,\tfrac12)$ with $|v(\tilde z)|\le C\norm v_{H^2(B(\tilde z,1/2))}$.

Apply Step~3 to $v=\tilde u$ and use Step~2 for $\tilde h=\tilde u$ and $\tilde h=\tilde f$. The continuous representative of $\tilde u$ is $\Gamma$-invariant
and descends to $M$. \qed
\end{proof}

\begin{corollary}[Corollary 1.6 on bounded geometry, printed class]\label{pc:cor16bg}
Under the hypotheses of Theorem~\ref{pc:main}, assume in addition that the injectivity radius of $M$ is at least $r_0>0$. Then for every $z\in M$,
\[
  \big|\Delta^N\big[\cU(t)\big(a\,\ue^{i\varphi_b/\hbar}\big)_\Gamma-u^{(K)}(t)_\Gamma\big](z)\big|
  \le C^{N+K+1}(2N)!\,K!\,(\alpha\hbar)^{-2N-4}\Big(\frac{\hbar t}{\alpha^2}\Big)^{K+1}\pcNorm{a},
\]
with $C$ depending also on $r_0$.
\end{corollary}

\begin{corollary}[Corollary 1.6 on every Fuchsian quotient, with the cusp factor]\label{pc:cor16}
Under the hypotheses of Theorem~\ref{pc:main}, for every $z\in M$,
\begin{multline*}
  \big|\Delta^N\big[\cU(t)\big(a\,\ue^{i\varphi_b/\hbar}\big)_\Gamma-u^{(K)}(t)_\Gamma\big](z)\big|\\
  \le C\,m(z)^{1/2}\,C^{N+K+1}(2N)!\,K!\,(\alpha\hbar)^{-2N-4}\Big(\frac{\hbar t}{\alpha^2}\Big)^{K+1}\pcNorm{a}.
\end{multline*}
In both corollaries the error lies in $\operatorname{Dom}(\Delta_M^N)$, and the left side is the value of the continuous representative.
\end{corollary}
\begin{proof}
Let $E$ be the bracket in the statement. By Theorem~\ref{pc:main} at $N$ and at $N+1$, $u:=\Delta^NE$ lies in $\operatorname{Dom}(\Delta_M)$. Apply
Lemma~\ref{pc:g43} to $u$, and use $N^{2N}\le(\ue/2)^{2N}(2N)!$, $(N+1)^{2N+2}\le54^{N+1}(2N)!$ and $\alpha\hbar\le1$. This proves
Corollary~\ref{pc:cor16}. Under the hypothesis of Corollary~\ref{pc:cor16bg}, $m(z)\le C(r_0)$, as in the proof of Corollary~\ref{cor:dd16-appx}. \qed
\end{proof}

\begin{remark}[cusps]\label{pc:remcusp}
On a bounded-geometry quotient $m(z)\le C(r_0)$, so Corollary~\ref{pc:cor16bg} is the case of Corollary~\ref{pc:cor16}. In an embedded cusp of
height $H$, deep enough, $m(z)\asymp H$, so the factor $m(z)^{1/2}$ grows like the square root of the cusp height. Lemma~\ref{pc:g43} is sharp in this sense:
the cusp bump $u_S$ of Remark~\ref{rem:g4} at height $H=\ue^{S}$ has $\sup|u|/(\norm{\Delta u}+\norm{u})\asymp H^{1/2}$.

The proofs of the two corollaries use one Laplacian and give the exponent $-2N-2$ of $\alpha\hbar$; we state them with $-2N-4$ for the comparison with
\cite[(1.34)]{S}. Since $m(z)^{1/2}(\alpha\hbar)^{-2N-4}\le(\alpha\hbar)^{-2N-8}$ when $m(z)\le(\alpha\hbar)^{-8}$, Corollary~\ref{pc:cor16} gives the printed
prefactor $(\alpha\hbar)^{-2N-8}$ of \cite[(1.34)]{S} at every such $z$, with $(\hbar t/\alpha^2)^{K+1}$ in place of $(\hbar t)^K$ as in
Corollary~\ref{pc:cor16bg} (the literal $(\hbar t)^K$ only for fixed $\alpha$). In a cusp this is the part below height $\asymp(\alpha\hbar)^{-8}$.
Above that height the factor cannot be removed in general: on $\mathrm{PSL}(2,\mathbb Z)\backslash\Hb$, for $b$ a cusp, Theorem~\ref{cl:main} gives a lower
bound of the same order $m(z)^{1/2}$ up to powers of $\hbar$. The proof of Lemma~\ref{pc:g43} also gives the local
form $|u(z)|\le C\,m(z)^{1/2}\bigl(\norm{\Delta u}_{L^2(B_M(z,1))}+\norm{u}_{L^2(B_M(z,1))}\bigr)$, with $B_M(z,1)$ the image of $B(\tilde z,1)$. So the
factor-free bound would follow if, for $u=\Delta^NE$, both $\norm{u}_{L^2(B_M(z,1))}$ and $\norm{\Delta u}_{L^2(B_M(z,1))}$ were at most
$C\,m(z)^{-1/2}(\alpha\hbar)^{-p}$ times the global bounds at $N$ and $N+1$ (with $p\le6$ for the printed prefactor). Theorem~\ref{cl:main} shows that this
condition fails high in the cusp of the modular surface.
\end{remark}

% Subsection for sec_printed.tex: the cusp factor of Corollary pc:cor16 is necessary.
% Notes source: notes/Cor16_cusps.md, Propositions G4.4 and G4.4P (reviews R22, R22c, R23, R23b). Labels prefix cl.
\subsection{The cusp factor is necessary}\label{subsec:cl}

Throughout this subsection $\Gamma=\mathrm{PSL}(2,\mathbb Z)$, $M=\Gamma\backslash\Hb$ (cusp at $\infty$ of width $1$), $x_0=i$, $X_0=\cosh d(i,\cdot)$,
$b=\infty$ and $\varphi=-\log y$. For a function $F$ on $\Hb$ we write $[F]_\Gamma(z)=\sum_{\gamma\in\Gamma}F(\gamma z)$, and
$E_K(t):=\cU(t)(a\,\ue^{i\varphi/\hbar})_\Gamma-u^{(K)}(t)_\Gamma$. Words $Z_S$ are products of $X$ and $Y$.

\begin{theorem}[the cusp factor is necessary]\label{cl:main}
Let $\beta>\tfrac12$, let $\alpha_0<\tfrac12$ be small enough for Theorems~\ref{thm:dd-appx} and~\ref{pc:main} and for Lemma~\ref{cl:amp}(i), and fix
$\alpha\in(0,\alpha_0]$; the constants below depend on the fixed $\alpha$, $\beta$ and $\kappa$.
\begin{enumerate}
\item[(i)] (Swapped class.) For $a=\ue^{-20X_0}$ there are $\hbar_0,c_*,t_*>0$ such that for $0<\hbar\le\hbar_0$, every
$t$ with $\max(2\log(1/\hbar),t_*)\le t\le\min(\hbar^{-1/5},\delta(\alpha/\hbar)^{1/2})$, and every $z=x+i\ue^t$,
\[
|E_0(t)(z)|\ge c_*\,\ue^{t/2}\hbar^{3/2}t,\qquad m(z)\le5\ue^t .
\]
\item[(ii)] (Printed class.) The same holds for $a=\ue^{-\alpha\Lambda}\ue^{-\kappa X_0}$, with $\kappa\ge\kappa_*$ and $\alpha\le\alpha_*(\kappa):=\min(\alpha_0,c\kappa^{-6})$; here the
proof uses Theorem~\ref{pc:main}.
\end{enumerate}
Consequently, on $\mathrm{PSL}(2,\mathbb Z)\backslash\Hb$ and for this $b$, the factor $m(z)^{1/2}$ in Corollary~\ref{pc:cor16} (and in its swapped-class
version) cannot be removed. For every $p\ge0$ no bound $\sup_M|E_0(t)|\le C(\alpha\hbar)^{-p}\norm a$ with $C$ independent of $\hbar$ holds on the window (take
$t=(2p+4)\log(1/\hbar)$), and for $\gamma<\tfrac12$ no bound $C\,m(z)^{\gamma}\hbar^{-p}\norm a$ with $C$ independent of $\hbar$ holds (take $t=\hbar^{-1/5}$). In particular the printed Corollary~1.6, \cite[(1.34)]{S} at
$N=K=0$, which is claimed for every Fuchsian $\Gamma$ and every $b$, is false when its constants $C,\delta$ do not depend on $\hbar$
(Remark~\ref{cl:reading}): at $t=20\log(1/\hbar)$ the lower bound (ii) exceeds $C(\alpha\hbar)^{-8}\norm a_{\alpha,\beta}$ for small $\hbar$.
\end{theorem}

\begin{remark}[the reading of the printed constants]\label{cl:reading}
\cite[Theorem~1.4 and Corollary~1.6]{S} choose $C,\delta$ after ``for all $\alpha,\beta$ with $\beta>1/2$ and $\alpha^3\ge\beta^2\hbar$''
and do not say whether they may depend on $\hbar$. We read them as independent of $\hbar$, as in a semiclassical estimate and as in all
our theorems. Under that reading the printed Corollary~1.6 is false. The contradiction is asymptotic: it needs one $C$ valid as $\hbar$ tends
to $0$. It does not rule out a bound with $\hbar$-dependent constants, and we do not claim that such a bound holds.
\end{remark}

\begin{remark}[the approximation is still accurate there]\label{cl:relative}
At the points of Theorem~\ref{cl:main}, in case (i) (the swapped-class amplitude $a=\ue^{-20X_0}$), the solution itself is large and the error remains
small relative to it: for
$\hbar\le\hbar_0$ and $t$ in the range of Theorem~\ref{cl:main},
\[
c\,\hbar t\le\frac{|E_0(t)(z)|}{|u^{(0)}(t)_\Gamma(z)|}\le C\,\hbar t,\qquad z=x+i\ue^t .
\]
So what fails is the form of the printed bound (a fixed multiple of $\norm a$, independent of $z$), not the accuracy of the approximation: the
periodized approximation and the solution both grow like $m(z)^{1/2}\hbar^{1/2}$ at these points, and their difference is a factor $O(\hbar t)$
smaller. (This is a statement about this amplitude on the window of Theorem~\ref{cl:main}; it is not a general relative-error theorem.)
\end{remark}
\begin{proof}
Lemma~\ref{cl:asym} with $A=a$ gives
\[
u^{(0)}(t)_\Gamma(z)=\ue^{-it/(2\hbar)}\bigl(\ue^{t/2}\sqrt{\pi\hbar}\,\ue^{i\pi/4}\ue^{it/\hbar}S_a(\hbar)+O(\ue^{t/2}\hbar^{3/2})\bigr),
\]
where $S_a(\hbar)=\sum_{c\ge1}c^{2i/\hbar}\sum_{(a',c)=1}a(a'/c+i/c^2)$.
With $a=\ue^{-\kappa X_0}$ and $\kappa=20$, $a(i)=\ue^{-\kappa}$ and the other terms are at most $\ue^{-\kappa}\ue^{-\kappa(X_0-1)}$; the count in the proof of
Lemma~\ref{cl:amp}(iii) gives $0.99\,\ue^{-\kappa}\le|S_a(\hbar)|\le1.01\,\ue^{-\kappa}$. The proof of Theorem~\ref{cl:main} gives
\[
|E_0(z)|=\ue^{t/2}\sqrt{\pi\hbar}\,\tfrac{\hbar t}2\,|S_0(\hbar)|\,\bigl(1+O(1/t)+O(\hbar t)+O(\hbar^{1/2}t^2)\bigr),
\]
and $0.98\,\kappa\ue^{-\kappa}\le|S_0(\hbar)|\le1.02\,\kappa\ue^{-\kappa}$. Divide.
\end{proof}

\begin{remark}[scope]\label{cl:scope}
Theorem~\ref{cl:main} extends to every $b\in\mathbb Q$, the other cusps of $\mathrm{PSL}(2,\mathbb Z)$: if $b=\gamma\infty$ with $\gamma\in\Gamma$, then
$\varphi_b=\varphi_\infty\circ\gamma^{-1}$ up to a constant, the approximation for $b$ is the $\gamma$-transport of the one for $\infty$, and the error for
$(b,a)$ is a constant phase times the error for $(\infty,a\circ\gamma)$. So take $a_b:=a\circ\gamma^{-1}$, that is $\ue^{-\kappa\cosh d(\gamma i,\cdot)}$ in case
(i) and $\ue^{-\alpha\Lambda}\ue^{-\kappa\cosh d(\gamma i,\cdot)}$ in case (ii) ($\Lambda$ commutes with isometries); these have finite norms about the base
point $i$, since $\langle d(i,\cdot)\rangle\le\langle d(\gamma i,\cdot)\rangle+d(i,\gamma i)$ and the tube argument of Lemma~\ref{cl:amp}(i) is isometry
invariant. Like the printed Corollary~1.6, Corollary~\ref{pc:cor16} is stated uniformly for every Fuchsian $\Gamma$ and every $b$; by Theorem~\ref{cl:main} no
statement of that form can have $m(z)^\gamma$ with $\gamma<\tfrac12$, so the exponent $\tfrac12$ is optimal among them (the power of $\hbar$ is not
claimed optimal). For irrational $b$, and for cusped quotients other than $\mathrm{PSL}(2,\mathbb Z)\backslash\Hb$, the necessity of the factor is not decided
here; a bound depending on $b$ without that factor is not excluded. The mechanism below is not special to $b=\infty$: in the frame where
$b=\infty$, the top of every $\Gamma$-image of a cusp horocycle is a stationary point of $-\log y$.
\end{remark}

The mechanism is the following. The error contains the explicit first Dyson term $\ue^{-it/(2\hbar)}(S(t)\,\tfrac{i\hbar}2P_1(t)a)\ue^{i\varphi/\hbar}$,
periodized. High in the cusp the sum over $\Gamma$ concentrates, by stationary phase, at the tops $a/c+i/(c^2Y)$ of the horocycles $\gamma(\mathbb R+iY)$,
whose images under $D_t:(x,y)\mapsto(x,\ue^ty)$ are the points $a/c+i/c^2$,
and its size is $\ue^{t/2}\sqrt{\pi\hbar}\cdot\hbar t$; the remainder is controlled by the $L^2$ theorem.

\begin{lemma}[semiclassical local Sobolev inequality]\label{cl:ssob}
For every Fuchsian $\Gamma$ there is an absolute $C$ with $|u(z)|\le C\,m(z)^{1/2}\bigl(r^{-1}\norm u_{L^2(M)}+r^3\norm{\Delta^2u}_{L^2(M)}\bigr)$ for
$0<r\le1$, $z\in M$ and $u\in H^4(M)$. For $\Gamma=\mathrm{PSL}(2,\mathbb Z)$ and $z=x+iY$ with $Y\ge4$, $m(z)\le5Y$.
\end{lemma}
\begin{proof}
By homogeneity take $\tilde z=i$. In normal coordinates $\xi$ at $i$ put $v(\xi)=\tilde u(\exp_i(r\xi))$ for $|\xi|<1$. Then
$r^4(\Delta^2\tilde u)(\exp_i(r\xi))=L_r^2v(\xi)$, where $L_r$ is $r^2\Delta$ in these coordinates: a uniformly elliptic operator whose coefficients are
bounded in $C^\infty$ uniformly for $0<r\le1$. Interior regularity for $L_r^2$ and $H^4\subset L^\infty$ give
$|v(0)|\le C(\norm v_{L^2(B_1)}+\norm{L_r^2v}_{L^2(B_1)})$. Since $d\nu\asymp r^2d\xi$ on the ball,
$|\tilde u(i)|\le C(r^{-1}\norm{\tilde u}_{L^2(B(i,r))}+r^3\norm{\Delta^2\tilde u}_{L^2(B(i,r))})$, and $B(\tilde z,r)\subset B(\tilde z,1)$, so the unfolding step of
Lemma~\ref{pc:g43} applies. For the count: if $\gamma\notin\Gamma_\infty$ then $\operatorname{Im}\gamma z\le1/Y$ and $d(z,\gamma z)\ge2\log Y\ge2$; and
$d(z,z+n)=2\operatorname{arcsinh}(|n|/(2Y))<2$ if and only if $|n|<2Y\sinh1$. So $m(z)\le4Y\sinh1+1\le5Y$ for $Y\ge4$.
\end{proof}

\begin{lemma}[cusp asymptotics]\label{cl:asym}
Let $\vartheta$ be either $X_0^{-2}$ or $\ue^{-d(i,\cdot)}(1+d(i,\cdot))^{-3/2}$, and for $A\in C^4(\Hb)$ put
$B(A):=\sup_w\vartheta(w)^{-1}\sum_{|S|\le4}|Z_SA(w)|$. There is a $C$ (depending only on the choice of $\vartheta$) such that for $0<\hbar\le\tfrac12$,
$t\ge2\log(1/\hbar)$, $Y=\ue^t$ and $z=x+iY$,
\begin{multline*}
\Bigl|\bigl[(S(t)A)\ue^{i\varphi/\hbar}\bigr]_\Gamma(z)-\ue^{t/2}\sqrt{\pi\hbar}\,\ue^{i\pi/4}\ue^{it/\hbar}\sum_{c\ge1}c^{2i/\hbar}\sum_{a\in\mathbb Z,(a,c)=1}A\Bigl(\frac ac+\frac i{c^2}\Bigr)\Bigr|\\
\le C\,\ue^{t/2}\hbar^{3/2}B(A),
\end{multline*}
and the double series converges absolutely.
\end{lemma}
\begin{proof}
\emph{Cosets.} An element of $\mathrm{PSL}(2,\mathbb Z)$ is a translation $T^k$, or has lower-left entry $c\ge1$ and acts by $w\mapsto a/c-1/(c(cw+d))$ with
$ad\equiv1\pmod c$. For fixed $c$ and $a$ coprime to $c$, $d$ runs over $d_a+c\mathbb Z$, and $d=d_a+cn$ corresponds to $w=z+n$. Hence
$[F]_\Gamma(z)=\sum_kF(z+k)+\sum_{c\ge1}\sum_{(a,c)=1}\sum_{n\in\mathbb Z}G_{a,c}(n)$ with $G_{a,c}(u)=F(\gamma_{a,c}(z+u))$. The translation part is
bounded by $\ue^{-t/2}B(A)\sum_k\vartheta(x+k+i\ue^{2t})\le C\ue^{-t/2}B(A)$, which is $\le C\ue^{t/2}\hbar^{3/2}B(A)$.

\emph{Change of variable.} With $\tau=(c(x+u)+d_a)/Y$, the point $D_t\gamma_{a,c}(z+u)$, where $D_t(x,y)=(x,\ue^ty)$ and $S(t)A=\ue^{-t/2}A\circ D_t$, is
\[
Q(\tau)=\frac ac-\frac{\tau}{cY(\tau^2+c^2)}+\frac{i}{\tau^2+c^2},\qquad \varphi(\gamma_{a,c}(z+u))=t+\log(\tau^2+c^2),
\]
so $G_{a,c}(u)=g(\tau):=\ue^{-t/2}A(Q(\tau))\,\ue^{i(t+\log(\tau^2+c^2))/\hbar}$ and $du=(Y/c)\,d\tau$. Along $Q$,
$\partial_\tau A(Q)=-(YA)(Q)\,\frac{2\tau}{\tau^2+c^2}-(XA)(Q)\,\frac{c^2-\tau^2}{cY(\tau^2+c^2)}$, so $|\partial_\tau^jA(Q)|\le C_jc^{-j}B(A)\vartheta(Q)$ for
$j\le4$. From $X_0(Q)=(q^2+1)/(2\operatorname{Im}Q)+\operatorname{Im}Q/2$ with $|q-a/c|\le1/(2c^2Y)$ one gets
$\vartheta(Q(\tau))\le C\,\vartheta_{a,c}\,(1+\tau^2/c^2)^{-s}$, where $s=2$, $\vartheta_{a,c}=(a^2+c^2)^{-2}$ for $\vartheta=X_0^{-2}$, and $s=1$,
$\vartheta_{a,c}=(a^2+c^2)^{-1}(1+\log(a^2+c^2))^{-3/2}$ for the second choice. In both cases $\sum_{(a,c)}\vartheta_{a,c}<\infty$ (for the second, by
partial summation, since $\#\{(a,c):a^2+c^2\le N\}\le CN$).

\emph{Poisson summation.} $\sum_nG(n)=\sum_m\hat G(m)$ with
\[
\hat G(m)=\frac Yc\,\ue^{2\pi im(x+d_a/c)}\int g(\tau)\,\ue^{-2\pi imY\tau/c}\,d\tau .
\]
For $m\ne0$ the total phase $\Phi$ has $|\Phi'|\ge2\pi|m|Y/c-1/(\hbar c)\ge\pi|m|Y/c$ (as $\hbar Y\ge\hbar^{-1}\ge2$), $|\Phi''|\le2/(\hbar c^2)$ and
$|\Phi'''|\le C/(\hbar c^3)$. Two integrations by parts give
\[
\Bigl|\int g\,\ue^{-2\pi imY\tau/c}d\tau\Bigr|\le C\Bigl(\frac c{|m|Y}\Bigr)^2c^{-2}\ue^{-t/2}B(A)\int\vartheta(Q(\tau))\,d\tau\le C\ue^{-t/2}B(A)\,\vartheta_{a,c}\,\frac c{m^2Y^2},
\]
using $\int(1+\tau^2/c^2)^{-s}d\tau\le Cc$. Hence $\sum_{m\ne0}|\hat G(m)|\le C\ue^{-t/2}B(A)\vartheta_{a,c}/Y$, and the sum over $(a,c)$ is $\le C\ue^{-3t/2}B(A)$.

\emph{Stationary phase.} With $\tau=c\sigma$, $\hat G(0)=Y\ue^{it/\hbar}c^{2i/\hbar}\int f(\sigma)\ue^{i\theta(\sigma)/\hbar}d\sigma$, where
$f(\sigma)=\ue^{-t/2}A(Q(c\sigma))$ and $\theta(\sigma)=\log(1+\sigma^2)$. Let $\chi$ be a cutoff, $\chi=1$ on $[-\tfrac12,\tfrac12]$, $\operatorname{supp}\chi\subset[-1,1]$.
Near part: $\theta=s(\sigma)^2$ with $s(\sigma)=\operatorname{sign}(\sigma)(\log(1+\sigma^2))^{1/2}$ smooth and $s'(0)=1$; for $h(s)=(\chi f)(\sigma(s))\sigma'(s)$,
\[
\int h\,\ue^{is^2/\hbar}ds=\sqrt{\pi\hbar}\,\ue^{i\pi/4}\,\frac1{2\pi}\int\hat h(\xi)\,\ue^{-i\hbar\xi^2/4}d\xi,
\]
and $|\ue^{-i\hbar\xi^2/4}-1|\le\hbar\xi^2/4$ gives the error $C\hbar^{3/2}\norm h_{C^4}\le C\hbar^{3/2}\ue^{-t/2}B(A)\vartheta_{a,c}$. Far part, on $|\sigma|\ge\tfrac12$, where $|\theta'|\ge(1+|\sigma|)^{-1}/2$ and
$|\theta^{(j)}|\le C(1+|\sigma|)^{-j}$. First truncate: with a cutoff $\chi_Y$ equal to $1$ on $|\sigma|\le Y$ and $0$ on $|\sigma|\ge2Y$, the integral over
the support of $1-\chi_Y$ is bounded absolutely by $C\ue^{-t/2}B(A)\vartheta_{a,c}/Y\le C\hbar^2\ue^{-t/2}B(A)\vartheta_{a,c}$. On $\tfrac12\le|\sigma|\le2Y$
integrate by parts once with $D^Tu=-(u/\theta')'$; the derivatives of $\chi_Y$ add terms of order $\hbar/Y$. In $\partial_\sigma f$ the $Y$-terms carry
$2\sigma/(1+\sigma^2)$ and gain $(1+|\sigma|)^{-1}$ per derivative, and the $X$-terms carry $1/Y$. The $X$-part of $D^T[(1-\chi)\chi_Yf]$ is integrated directly,
giving at most $C\hbar\,\ue^{-t/2}B(A)\vartheta_{a,c}Y^{-1}\log Y\le C\hbar^{3/2}\ue^{-t/2}B(A)\vartheta_{a,c}$ (as $Y^{-1}\log Y=t\ue^{-t}\le\hbar^{1/2}$); the
remaining symbol part is integrated by parts once more, giving at most $C\hbar^2\ue^{-t/2}B(A)\vartheta_{a,c}$. Since $f(0)=\ue^{-t/2}A(a/c+i/c^2)$,
\[
\hat G(0)=\ue^{t/2}\sqrt{\pi\hbar}\,\ue^{i\pi/4}\ue^{it/\hbar}c^{2i/\hbar}A\Bigl(\frac ac+\frac i{c^2}\Bigr)+O\bigl(\ue^{t/2}\hbar^{3/2}B(A)\vartheta_{a,c}\bigr).
\]
Summing over $(a,c)$ proves the lemma.
\end{proof}

\begin{lemma}[the Poisson kernel off the diagonal]\label{cl:kernel}
Fix $r_0\in(0,1]$. For $0<\alpha\le\alpha_0<\pi/2$ the kernel $k_\alpha(r)$ of $\ue^{-\alpha\Lambda}$ extends holomorphically to
$\{a+i\theta:a>3r_0/4,\ |\theta|<r_0/4\}$ with $|\tilde k_\alpha(a+i\theta)|\le C\alpha\,\ue^{-a}(1+a)^{-3/2}$. Hence
$|\partial_r^jk_\alpha(r)|\le C_j\alpha\,\ue^{-r}(1+r)^{-3/2}$ for $r\ge r_0$, uniformly in $\alpha\le\alpha_0$.
\end{lemma}
\begin{proof}
By Lemma~\ref{os:hk} with $\theta_0=r_0/4$,
\[
|\tilde k_\alpha(a+i\theta)|\le C\int_0^\infty\alpha(4\pi)^{-1/2}s^{-3/2}\ue^{-(\alpha^2-\theta^2)/(4s)}p_s(a)\,ds .
\]
The computation in the proof of Lemma~\ref{os:poisson}, with $\alpha^2-\theta^2$ in place of $\alpha^2$, turns this into
\[
C\alpha\int_a^\infty\frac{\rho\,K_2\bigl(\sqrt{\rho^2+\alpha^2-\theta^2}/2\bigr)}{\rho^2+\alpha^2-\theta^2}\,\frac{d\rho}{(\cosh\rho-\cosh a)^{1/2}},
\]
which converges because $\rho^2+\alpha^2-\theta^2\ge r_0^2/2$, and is bounded as in \textup{(\ref{os:ub})}. Cauchy estimates on discs of radius $r_0/8$ give the derivative bounds.
\end{proof}

\begin{lemma}[the amplitudes]\label{cl:amp}
(i) $a=\ue^{-\kappa X_0}$ has finite swapped norm for every $\alpha\le r_1<\tfrac12$ and $\beta$, and $|Z_Sa|\le C_{S,\kappa}X_0^{|S|}\ue^{-\kappa X_0}$.
(ii) $a=\ue^{-\alpha\Lambda}\ue^{-\kappa X_0}$ has $\norm a_{\alpha,\beta}=\norm{\ue^{\beta\langle d\rangle}\ue^{-\kappa X_0}}<\infty$, and
$|Z_Sa(w)|\le C_{S,\kappa}\ue^{-d(i,w)}(1+d(i,w))^{-3/2}$ uniformly in $\alpha\le\alpha_0$.
(iii) Put $A_0=(Y^2-Y)a$ and $S_0(\hbar)=\sum_{c\ge1}c^{2i/\hbar}\sum_{(a,c)=1}A_0(a/c+i/c^2)$. In case (i) with $\kappa=20$, $|S_0(\hbar)|\ge0.98\,\kappa\ue^{-\kappa}$
for all $\hbar$; in case (ii), $|S_0(\hbar)|\ge\tfrac12\kappa\ue^{-\kappa}$ for $\kappa\ge\kappa_*$, $\alpha\le c\kappa^{-6}$ and all $\hbar$.
\end{lemma}
\begin{proof}
(i) $f=\ue^{\beta\langle d\rangle}a$ continues to the crown with $\operatorname{Re}\tilde X_0=\cos(2r)X_0$ and $\operatorname{Re}\widetilde{\langle d\rangle}\le\langle d\rangle+C_Hr^2$ (proof of
Lemma~\ref{lem:ex}), so its tube norm is finite and $\norm{\ue^{\alpha\Lambda}f}\le C_0N_\alpha(f)$ as there. The vector fields act by $XX_0=x$, $YX_0=y-X_0$,
$Xx=y$, $Yy=y$, $Xy=Yx=0$, and $|x|,y\le2X_0$; this gives the word bound.
(ii) The norm is the definition. For the words, transport $i$ to $w'$ by the affine group, under which $X$ and $Y$ are invariant; then
$|Z_S^w\cosh d(w,w')|\le C_S\cosh d(w,w')$, and Lemma~\ref{cl:kernel} gives $|Z_S^wk_\alpha(d(w,w'))|\le C_S\alpha\ue^{-d(w,w')}(1+d(w,w'))^{-3/2}$ for
$d(w,w')\ge\tfrac14$. Split $h=\ue^{-\kappa X_0}$ near each $w$ as $h\psi_w+h(1-\psi_w)$, $\psi_w=1$ on $B(w,\tfrac18)$, $\operatorname{supp}\psi_w\subset B(w,\tfrac14)$:
the far part is bounded with the kernel estimate, and the near part (with $\psi_w$ frozen at the evaluation point) by local Sobolev and
$\norm{\ue^{-\alpha\Lambda}g}_{H^k}\le\norm g_{H^k}$; the words of $h$ carry powers of $\kappa$ and $X_0$, so the near part is at most
$C_{S,\kappa}\exp(-\tfrac\kappa2\cosh(\max(0,d(i,w)-\tfrac14)))$, which is dominated by $\ue^{-d}(1+d)^{-3/2}$.
(iii) $Y^2X_0=X_0$ gives $A_0=a[\kappa^2(y-X_0)^2-\kappa(2X_0-y)]$ in case (i), so $A_0(i)=-\kappa\ue^{-\kappa}$. At the other points $a/c+i/c^2$,
$X_0=(a^2+c^2+c^{-2})/2$, $X_0-1\ge\max(\tfrac12,(n-2)/2)$ with $n=a^2+c^2$, and $|A_0|\le\kappa\ue^{-\kappa}\ue^{-\kappa(X_0-1)}(\kappa X_0^2+2X_0)$. For
$\kappa=20$ and $q=\ue^{-10}<1/20000$: the two points $\pm1+i$ give at most $96q<0.005$; there is no coprime pair with $c\ge1$ and $n\in\{3,4\}$; the rest is
at most $\sum_{n\ge5}(4n+4)[5(n+1)^2+n+1]q^{n-2}\le5184q^3/(1-8q)<0.01$. In case (ii), $|A_0(i)+\kappa\ue^{-\kappa}|\le C\norm{(\ue^{-\alpha\Lambda}-1)h}_{H^4}\le
C\alpha\norm{(1-\Delta)^3h}\le C\alpha\kappa^6\ue^{-\kappa}$. The other points are $0.96$-separated ($\cosh d\ge\tfrac32$ for distinct ones). Split
$h=h_1+h_2$ with $h_1=h\psi$, $\psi=1$ on $B(i,0.44)$, $\operatorname{supp}\psi\subset B(i,0.45)$, so that $h_2$ vanishes on $B(i,0.44)$: $h_1$ contributes at most $C\alpha\kappa^{-1}\ue^{-\kappa}\vartheta(w)$ at each point by Lemma~\ref{cl:kernel}; for $h_2$, whose
derivatives up to order $4$ are at most $C(1+\kappa)^4(1+X_0)^4\ue^{-\kappa X_0}$, the split of (ii) at each point and the packing bound
$\#\{w:d(i,w)\le R\}=O(\ue^R)$ give at most $C\kappa^8\ue^{-\kappa\cosh0.44}$ in total, with $C$ independent of $\kappa\ge1$ and $\alpha\le\alpha_0$. Hence the sum
over the other points is at most $C(\alpha/\kappa^2+\kappa^8\ue^{-0.098\kappa})\kappa\ue^{-\kappa}$.
\end{proof}

\begin{proof}[Proof of Theorem~\ref{cl:main}]
With $c_\hbar=i\hbar/2$, the definition of $u^{(K)}$ gives $E_0=M_1+M_2+E_2$, where
$M_k=\ue^{-it/(2\hbar)}[(S(t)A_k)\ue^{i\varphi/\hbar}]_\Gamma$, $A_1=c_\hbar P_1(t)a$ and $A_2=c_\hbar^2P_2(t)a$. Since
$P_1(t)a=\int_0^t\Delta(s)a\,ds=\tfrac{1-\ue^{-2t}}2X^2a+t(Y^2-Y)a$, Lemma~\ref{cl:asym} gives
\[
|M_1(z)|\ge\ue^{t/2}\sqrt{\pi\hbar}\,\tfrac\hbar2\bigl(t|S_0(\hbar)|-S_X\bigr)-C\ue^{t/2}\hbar^{3/2}\hbar t,
\]
with $S_X=\sum_{(a,c)}|X^2a(a/c+i/c^2)|<\infty$. $P_2(t)a$ is a combination of
words of length at most $4$ with coefficients bounded by $t^2/2$, so $|M_2(z)|\le C\ue^{t/2}\hbar^{1/2}\hbar^2t^2$. For $E_2$, Theorem~\ref{thm:dd-appx} in case (i)
and Theorem~\ref{pc:main} in case (ii), at $K=2$ and $N=0,2$, give $\norm{E_2}\le C(\hbar t/\alpha^2)^3\norm a$ and
$\norm{\Delta^2E_2}\le C(\alpha\hbar)^{-4}(\hbar t/\alpha^2)^3\norm a$; Lemma~\ref{cl:ssob} with $r=\hbar$ gives $|E_2(z)|\le C\ue^{t/2}\hbar^{-1}\alpha^{-10}(\hbar t)^3\norm a$.
Hence
\[
|E_0(z)|\ge\ue^{t/2}\hbar^{3/2}t\Bigl[\tfrac{\sqrt\pi}2|S_0(\hbar)|-\tfrac{\sqrt\pi}2\frac{S_X}t-C\hbar-C\hbar t-C\hbar^{1/2}t^2\Bigr],
\]
and the bracket is at least $\tfrac{\sqrt\pi}4|S_0(\hbar)|$ for $t\ge t_*$, $t\le\hbar^{-1/5}$ and $\hbar\le\hbar_0$. Lemma~\ref{cl:amp}(iii) bounds $|S_0|$ below.
\end{proof}



\section{Sharpness and limitations}\label{sec:g5}
\subsection{The order \texorpdfstring{$\hbar t$}{ht} of the leading error is attained}
\begin{proposition}\label{prop:g5}
Let $\Gamma=\{\mathrm{id}\}$, allowed in \cite{S} since a Fuchsian group there is any discrete subgroup.
Let $a\ne0$ be such that every word in $X,Y$ applied to $a$ lies in $L^2(\Hb)$. This holds for every $a$ of finite swapped
norm (Lemmas~\ref{lem:m2w} and~\ref{lem:sc}) and for every $a$ of finite printed norm
(Lemma~\ref{lem:lift}, since then $a\in\operatorname{Dom}(\ue^{\alpha\Lambda})$), in both cases with a smaller radius in the
cited lemmas if necessary; nonzero amplitudes of the two kinds are
$\ue^{-\beta\langle d\rangle}\ue^{-\alpha\Lambda}g$ and $\ue^{-\alpha\Lambda}\ue^{-\beta\langle d\rangle}g$ with $0\ne g\in L^2(\Hb)$, whose swapped and printed
norm, respectively, is $\norm g$. Then for $t=d\,\hbar^{-1/2}$, $d>0$ fixed,
\[
\norm{u^{(0)}(t)-\cU(t)(a\ue^{i\varphi/\hbar})}_{L^2(\Hb)}=\tfrac d2\sqrt\hbar\,\norm{Aa}+O_{a,d}(\hbar),\qquad A=Y^2-Y,
\]
and $\norm{Aa}\ge\norm a/4>0$. Hence the order $\hbar t$ of Theorem~\ref{thm:main} is attained by every nonzero amplitude of the
swapped class, and the bound
$C\hbar\norm a_{\alpha,\beta}$ of \cite[(1.18)]{S} cannot hold on the window \cite[(1.19)]{S} for any nonzero amplitude of the printed class.
\end{proposition}
\begin{proof}
By \eqref{eq:factor} and unitarity of $\cU$ and $V(t)$, the left side equals $\norm{V(t)a-a}=\norm{a-V(t)^*a}$. The Dyson
identity (D1) of Proposition~\ref{prop:dd-id} at $K=1$, whose hypothesis is the one made on $a$, reads, with $c=i\hbar/2$,
\[a-V(t)^*a=c\,V(t)^*P_1(t)a+R,\qquad R=c^2\int_{t>s_1>s_2>0}V(s_1)^*\Delta(s_1)\Delta(s_2)a\,ds_1\,ds_2,\]
and, since $\Delta(s)=A+\ue^{-2s}X^2$,
\[P_1(t)=\int_0^t\Delta(s)\,ds=tA+\tfrac{1-\ue^{-2t}}2X^2\]
and
\[\norm{\Delta(r)\Delta(s)a}\le M_a:=\norm{A^2a}+\norm{AX^2a}+\norm{X^2Aa}+\norm{X^4a}<\infty\qquad(r,s\ge0).\]
So $\norm R\le\tfrac{\hbar^2t^2}{8}M_a$, and since $V(t)^*$ is unitary,
\[\bigl|\norm{V(t)a-a}-\tfrac{\hbar t}2\norm{Aa}\bigr|\le\tfrac\hbar4\norm{X^2a}+\tfrac{\hbar^2t^2}{8}M_a ,\]
which is $\hbar\,(\tfrac14\norm{X^2a}+\tfrac{d^2}8M_a)$ at $t=d\,\hbar^{-1/2}$. In the variable $r=\log y$ and after the unitary
gauge $a=\ue^{r/2}g$, $A$ becomes $\partial_r^2-1/4$, so $-A\ge1/4$.
\end{proof}

\begin{remark}[The formalized statement]\label{rem:g5-lean}
In the formalization (Appendix~\ref{app:lean}) the statement of this section is the negation of the printed theorem \cite[Theorem~1.3]{S}; the formal theorem does not prove the asymptotic formula. The printed statement is formalized as: for every Fuchsian $\Gamma$ and all $\alpha>0$, $\beta>1$ there are $C,\delta,\hbar_0>0$ such that
$\norm{u^{(0)}(t)_\Gamma-\cU(t)(a\ue^{i\varphi_b/\hbar})_\Gamma}_{L^2(M)}\le C\hbar\norm a_{\alpha,\beta}$ for all $0<\hbar\le\hbar_0$, all $0\le t\le\delta\hbar^{-1/2}$, all $b$ and all $C^\infty$ amplitudes of finite printed norm. The restrictions to $\hbar\le\hbar_0$ and to $C^\infty$ amplitudes weaken the printed claim, and the formal theorem states that this weaker claim is false. Its proof uses $\Gamma=\{\mathrm{id}\}$, $\alpha=\tfrac12$, $\beta=2$, $b=\infty$ and the amplitude $a=\ue^{-\alpha\Lambda}\ue^{-3X_0}$, $X_0=\cosh d(i,\cdot)$. It takes the Dyson identity (D4) of Proposition~\ref{prop:dd-id} at $K=0$ and $K=1$, which gives the expansion above with $R$ the second-order term, and uses that periodization over the trivial group is an isometry. The bound $\norm R\le\tfrac{\hbar^2t^2}8\sup_{r,s}\norm{\Delta(r)\Delta(s)a}$ comes from Lemma~\ref{lem:wb} at $N=0$, and $\norm{Aa}\ge\norm a/4$ from $Y^*=1-Y$ on $\operatorname{Dom}(\Delta)$: $\langle a,Aa\rangle=-\norm{Ya}^2$ and $2\operatorname{Re}\langle a,Ya\rangle=\norm a^2$.
\end{remark}

\subsection{The mean of the approximation}\label{subsec:mean}
The integral of $u^{(K)}(t)$ over $\Hb$ can be computed exactly. On a surface of finite area this bounds the error of the
approximation from below at large times. For $z\in\mathbb C$ let $T_K(z):=\sum_{k=0}^{K}z^k/k!$.

\begin{proposition}[the mean of the approximation]\label{prop:mean}
Let $\beta>\tfrac12$, $0<\alpha\le1$ and let $a$ have finite swapped norm $\norm{\ue^{\alpha\Lambda}\ue^{\beta\langle d\rangle}a}<\infty$.
Let $b\in\partial\Hb$, $\hbar>0$, $K\in\mathbb N_0$ and $t\ge0$. Then $a\,\ue^{i\varphi_b/\hbar}$ and $u^{(K)}(t)$ are integrable
over $\Hb$, and
\[
  \int_{\Hb}u^{(K)}(t)\,d\nu=\ue^{-z}\,T_K(z)\,J,\qquad z:=-\frac t2-\frac{it}{2\hbar},\qquad
  J:=\int_{\Hb}a\,\ue^{i\varphi_b/\hbar}\,d\nu .
\]
Consequently, for every Fuchsian group $\Gamma$ with $\operatorname{area}(M)<\infty$, with $\psi_0:=(a\,\ue^{i\varphi_b/\hbar})_\Gamma$,
\[
  \ue^{t/2}\,|T_K(z)|\,|J|\le\operatorname{area}(M)^{1/2}\Bigl(\norm{\cU(t)\psi_0-u^{(K)}(t)_\Gamma}_{L^2(M)}+\norm{\psi_0}_{L^2(M)}\Bigr).
\]
\end{proposition}
\begin{proof}
By Lemma~\ref{lem:base} we may take $x_0=i$ and $b=\infty$: with $q$ as there, $q$ is an isometry, and both
$a\,\ue^{i\varphi_b/\hbar}$ and $u^{(K)}_b(t)[a]$ are carried to the corresponding functions for $\infty$ and $a\circ q$, times the same
constant of modulus $1$; the norm of $a\circ q$ with base point $i$ is that of $a$ with base point $x_0$. Write $e:=\ue^{i\varphi/\hbar}=y^{-i/\hbar}$ and $J(g):=\int_{\Hb}g\,e\,d\nu$ for $g\in L^1(\Hb)$.

\emph{Step 1 (the words of $a$ are integrable).} The amplitude $a$ is $C^\infty$ (Lemma~\ref{lem:lift}). By Lemma~\ref{lem:m2w} and
Lemma~\ref{lem:sc}, $X_0^\beta Z_Sa\in L^2(\Hb)$ for every word $Z_S$ in $X,Y$. Since $2\beta>1$, $X_0^{-\beta}\in L^2(\Hb)$. So
$Z_Sa\in L^1(\Hb)$ by the Cauchy--Schwarz inequality.

\emph{Step 2 (two integrations by parts).} Let $g$ be $C^\infty$ with $g$, $Xg$, $Yg$ in $L^1(\Hb)$. Then
\[
  J(Xg)=0,\qquad J(Yg)=\Bigl(1+\frac i\hbar\Bigr)J(g).
\]
For the first identity let $\sigma_s(x+iy):=x+sy+iy$. This shear preserves $\nu$ and $y$, and
$\tfrac{d}{ds}(g\circ\sigma_s)=(Xg)\circ\sigma_s$. So $g\circ\sigma_1-g=\int_0^1(Xg)\circ\sigma_s\,ds$ pointwise. Multiply by $e$ and
integrate over $\Hb$; the order of integration may be exchanged, because $\int_{\Hb}|(Xg)\circ\sigma_s|\,d\nu=\norm{Xg}_{L^1}$ for
every $s$. Since $J(h\circ\sigma_s)=J(h)$ for $h\in L^1$, this gives $0=J(g\circ\sigma_1)-J(g)=J(Xg)$. For the second identity let
$\delta_s(x+iy):=x+i\ue^sy$. Then $(\delta_s)_*\nu=\ue^s\nu$, $e\circ\delta_s=\ue^{-is/\hbar}e$ and
$\tfrac{d}{ds}(g\circ\delta_s)=(Yg)\circ\delta_s$, so that $J(h\circ\delta_s)=\ue^{cs}J(h)$ for $h\in L^1$, with $c:=1+i/\hbar$. In
the same way
\[
  (\ue^{c}-1)J(g)=J(g\circ\delta_1)-J(g)=\int_0^1J\bigl((Yg)\circ\delta_s\bigr)ds=J(Yg)\,\frac{\ue^{c}-1}{c},
\]
and $\ue^c\ne1$ because $|\ue^c|=\ue$.

\emph{Step 3 (the Dyson terms).} By Steps~1 and~2, $J(Z_Sa)=\chi(S)J(a)$ for every word, where $\chi$ is multiplicative on
words, $\chi(X)=0$ and $\chi(Y)=c$. So for every finite linear combination $g$ of words of $a$,
\[
  J(\Delta(s)g)=\lambda J(g),\qquad\lambda:=c^2-c=\frac i\hbar-\frac1{\hbar^2},
\]
because $\Delta(s)=\ue^{-2s}X^2+Y^2-Y$. The function $P_k(t)a$ is such a combination, with coefficients that are continuous in
$t$, and $P_{k+1}(t)a=\int_0^t\Delta(s)P_k(s)a\,ds$. By induction on $k$, $J(P_k(t)a)=(\lambda t)^kJ(a)/k!$. Since
$\tfrac{i\hbar}2\lambda t=z$,
\[
  J\bigl(a^{(K)}(t)\bigr)=\sum_{k=0}^{K}\Bigl(\frac{i\hbar}2\Bigr)^k\frac{(\lambda t)^k}{k!}J(a)=T_K(z)\,J(a).
\]

\emph{Step 4 (the transport).} By definition $u^{(K)}(t)=\ue^{-it/(2\hbar)}\ue^{-t/2}\bigl(a^{(K)}(t)\circ\delta_t\bigr)e$. So
$u^{(K)}(t)$ is integrable, and by Step~2
\[
  \int_{\Hb}u^{(K)}(t)\,d\nu=\ue^{-it/(2\hbar)-t/2}\,\ue^{ct}\,J\bigl(a^{(K)}(t)\bigr)=\ue^{-z}\,T_K(z)\,J(a).
\]

\emph{Step 5 (the surface).} Let $u\in L^1(\Hb)$ be such that $\sum_\gamma|u(\gamma z)|<\infty$ for almost every $z$ and
$u_\Gamma\in L^2(M)$. Unfolding gives $\int_{\Hb}u\,d\nu=\int_Mu_\Gamma\,d\nu$, hence
$|\int_{\Hb}u\,d\nu|\le\operatorname{area}(M)^{1/2}\norm{u_\Gamma}_{L^2(M)}$. For $u=u^{(K)}(t)$ the hypotheses hold by Step~4 and by
Lemma~\ref{lem:t2} with Lemma~\ref{lem:sc}, as in the proof of Theorem~\ref{thm:dd-appx}. Finally
$\norm{u^{(K)}(t)_\Gamma}\le\norm{\cU(t)\psi_0-u^{(K)}(t)_\Gamma}+\norm{\psi_0}$, because $\cU(t)$ is unitary on $L^2(M)$, and
$|\ue^{-z}|=\ue^{t/2}$.
\end{proof}

\begin{remark}[the time of validity on a surface of finite area]\label{rem:mean}
On a surface of finite area the solution $\cU(t)\psi_0$ has the norm $\norm{\psi_0}$ at all times, and by
Proposition~\ref{prop:mean} the error of the approximation of order $K$ is at least
$\operatorname{area}(M)^{-1/2}\,\ue^{t/2}|T_K(z)||J|-\norm{\psi_0}$. For $K=0$, where $T_0=1$, and $J\ne0$ the error is at least
$\norm{\psi_0}$ for all $t\ge2\log\bigl(2\operatorname{area}(M)^{1/2}\norm{\psi_0}/|J|\bigr)$. The number $J$ depends on $\hbar$ and
on $a$; it is not estimated here. On $\Hb$ itself, and on surfaces of infinite area, the proposition gives no such bound.
\end{remark}

% Subsection for sec_printed.tex: the cusp factor of Corollary pc:cor16 is necessary.
% Notes source: notes/Cor16_cusps.md, Propositions G4.4 and G4.4P (reviews R22, R22c, R23, R23b). Labels prefix cl.
\subsection{The cusp factor is necessary}\label{subsec:cl}

Throughout this subsection $\Gamma=\mathrm{PSL}(2,\mathbb Z)$, $M=\Gamma\backslash\Hb$ (cusp at $\infty$ of width $1$), $x_0=i$, $X_0=\cosh d(i,\cdot)$,
$b=\infty$ and $\varphi=-\log y$. For a function $F$ on $\Hb$ we write $[F]_\Gamma(z)=\sum_{\gamma\in\Gamma}F(\gamma z)$, and
$E_K(t):=\cU(t)(a\,\ue^{i\varphi/\hbar})_\Gamma-u^{(K)}(t)_\Gamma$. Words $Z_S$ are products of $X$ and $Y$.

\begin{theorem}[the cusp factor is necessary]\label{cl:main}
Let $\beta>\tfrac12$, let $\alpha_0<\tfrac12$ be small enough for Theorems~\ref{thm:dd-appx} and~\ref{pc:main} and for Lemma~\ref{cl:amp}(i), and fix
$\alpha\in(0,\alpha_0]$; the constants below depend on the fixed $\alpha$, $\beta$ and $\kappa$.
\begin{enumerate}
\item[(i)] (Swapped class.) For $a=\ue^{-20X_0}$ there are $\hbar_0,c_*,t_*>0$ such that for $0<\hbar\le\hbar_0$, every
$t$ with $\max(2\log(1/\hbar),t_*)\le t\le\min(\hbar^{-1/5},\delta(\alpha/\hbar)^{1/2})$, and every $z=x+i\ue^t$,
\[
|E_0(t)(z)|\ge c_*\,\ue^{t/2}\hbar^{3/2}t,\qquad m(z)\le5\ue^t .
\]
\item[(ii)] (Printed class.) The same holds for $a=\ue^{-\alpha\Lambda}\ue^{-\kappa X_0}$, with $\kappa\ge\kappa_*$ and $\alpha\le\alpha_*(\kappa):=\min(\alpha_0,c\kappa^{-6})$; here the
proof uses Theorem~\ref{pc:main}.
\end{enumerate}
Consequently, on $\mathrm{PSL}(2,\mathbb Z)\backslash\Hb$ and for this $b$, the factor $m(z)^{1/2}$ in Corollary~\ref{pc:cor16} (and in its swapped-class
version) cannot be removed. For every $p\ge0$ no bound $\sup_M|E_0(t)|\le C(\alpha\hbar)^{-p}\norm a$ with $C$ independent of $\hbar$ holds on the window (take
$t=(2p+4)\log(1/\hbar)$), and for $\gamma<\tfrac12$ no bound $C\,m(z)^{\gamma}\hbar^{-p}\norm a$ with $C$ independent of $\hbar$ holds (take $t=\hbar^{-1/5}$). In particular the printed Corollary~1.6, \cite[(1.34)]{S} at
$N=K=0$, which is claimed for every Fuchsian $\Gamma$ and every $b$, is false when its constants $C,\delta$ do not depend on $\hbar$
(Remark~\ref{cl:reading}): at $t=20\log(1/\hbar)$ the lower bound (ii) exceeds $C(\alpha\hbar)^{-8}\norm a_{\alpha,\beta}$ for small $\hbar$.
\end{theorem}

\begin{remark}[the reading of the printed constants]\label{cl:reading}
\cite[Theorem~1.4 and Corollary~1.6]{S} choose $C,\delta$ after ``for all $\alpha,\beta$ with $\beta>1/2$ and $\alpha^3\ge\beta^2\hbar$''
and do not say whether they may depend on $\hbar$. We read them as independent of $\hbar$, as in a semiclassical estimate and as in all
our theorems. Under that reading the printed Corollary~1.6 is false. The contradiction is asymptotic: it needs one $C$ valid as $\hbar$ tends
to $0$. It does not rule out a bound with $\hbar$-dependent constants, and we do not claim that such a bound holds.
\end{remark}

\begin{remark}[the approximation is still accurate there]\label{cl:relative}
At the points of Theorem~\ref{cl:main}, in case (i) (the swapped-class amplitude $a=\ue^{-20X_0}$), the solution itself is large and the error remains
small relative to it: for
$\hbar\le\hbar_0$ and $t$ in the range of Theorem~\ref{cl:main},
\[
c\,\hbar t\le\frac{|E_0(t)(z)|}{|u^{(0)}(t)_\Gamma(z)|}\le C\,\hbar t,\qquad z=x+i\ue^t .
\]
So what fails is the form of the printed bound (a fixed multiple of $\norm a$, independent of $z$), not the accuracy of the approximation: the
periodized approximation and the solution both grow like $m(z)^{1/2}\hbar^{1/2}$ at these points, and their difference is a factor $O(\hbar t)$
smaller. (This is a statement about this amplitude on the window of Theorem~\ref{cl:main}; it is not a general relative-error theorem.)
\end{remark}
\begin{proof}
Lemma~\ref{cl:asym} with $A=a$ gives
\[
u^{(0)}(t)_\Gamma(z)=\ue^{-it/(2\hbar)}\bigl(\ue^{t/2}\sqrt{\pi\hbar}\,\ue^{i\pi/4}\ue^{it/\hbar}S_a(\hbar)+O(\ue^{t/2}\hbar^{3/2})\bigr),
\]
where $S_a(\hbar)=\sum_{c\ge1}c^{2i/\hbar}\sum_{(a',c)=1}a(a'/c+i/c^2)$.
With $a=\ue^{-\kappa X_0}$ and $\kappa=20$, $a(i)=\ue^{-\kappa}$ and the other terms are at most $\ue^{-\kappa}\ue^{-\kappa(X_0-1)}$; the count in the proof of
Lemma~\ref{cl:amp}(iii) gives $0.99\,\ue^{-\kappa}\le|S_a(\hbar)|\le1.01\,\ue^{-\kappa}$. The proof of Theorem~\ref{cl:main} gives
\[
|E_0(z)|=\ue^{t/2}\sqrt{\pi\hbar}\,\tfrac{\hbar t}2\,|S_0(\hbar)|\,\bigl(1+O(1/t)+O(\hbar t)+O(\hbar^{1/2}t^2)\bigr),
\]
and $0.98\,\kappa\ue^{-\kappa}\le|S_0(\hbar)|\le1.02\,\kappa\ue^{-\kappa}$. Divide.
\end{proof}

\begin{remark}[scope]\label{cl:scope}
Theorem~\ref{cl:main} extends to every $b\in\mathbb Q$, the other cusps of $\mathrm{PSL}(2,\mathbb Z)$: if $b=\gamma\infty$ with $\gamma\in\Gamma$, then
$\varphi_b=\varphi_\infty\circ\gamma^{-1}$ up to a constant, the approximation for $b$ is the $\gamma$-transport of the one for $\infty$, and the error for
$(b,a)$ is a constant phase times the error for $(\infty,a\circ\gamma)$. So take $a_b:=a\circ\gamma^{-1}$, that is $\ue^{-\kappa\cosh d(\gamma i,\cdot)}$ in case
(i) and $\ue^{-\alpha\Lambda}\ue^{-\kappa\cosh d(\gamma i,\cdot)}$ in case (ii) ($\Lambda$ commutes with isometries); these have finite norms about the base
point $i$, since $\langle d(i,\cdot)\rangle\le\langle d(\gamma i,\cdot)\rangle+d(i,\gamma i)$ and the tube argument of Lemma~\ref{cl:amp}(i) is isometry
invariant. Like the printed Corollary~1.6, Corollary~\ref{pc:cor16} is stated uniformly for every Fuchsian $\Gamma$ and every $b$; by Theorem~\ref{cl:main} no
statement of that form can have $m(z)^\gamma$ with $\gamma<\tfrac12$, so the exponent $\tfrac12$ is optimal among them (the power of $\hbar$ is not
claimed optimal). For irrational $b$, and for cusped quotients other than $\mathrm{PSL}(2,\mathbb Z)\backslash\Hb$, the necessity of the factor is not decided
here; a bound depending on $b$ without that factor is not excluded. The mechanism below is not special to $b=\infty$: in the frame where
$b=\infty$, the top of every $\Gamma$-image of a cusp horocycle is a stationary point of $-\log y$.
\end{remark}

The mechanism is the following. The error contains the explicit first Dyson term $\ue^{-it/(2\hbar)}(S(t)\,\tfrac{i\hbar}2P_1(t)a)\ue^{i\varphi/\hbar}$,
periodized. High in the cusp the sum over $\Gamma$ concentrates, by stationary phase, at the tops $a/c+i/(c^2Y)$ of the horocycles $\gamma(\mathbb R+iY)$,
whose images under $D_t:(x,y)\mapsto(x,\ue^ty)$ are the points $a/c+i/c^2$,
and its size is $\ue^{t/2}\sqrt{\pi\hbar}\cdot\hbar t$; the remainder is controlled by the $L^2$ theorem.

\begin{lemma}[semiclassical local Sobolev inequality]\label{cl:ssob}
For every Fuchsian $\Gamma$ there is an absolute $C$ with $|u(z)|\le C\,m(z)^{1/2}\bigl(r^{-1}\norm u_{L^2(M)}+r^3\norm{\Delta^2u}_{L^2(M)}\bigr)$ for
$0<r\le1$, $z\in M$ and $u\in H^4(M)$. For $\Gamma=\mathrm{PSL}(2,\mathbb Z)$ and $z=x+iY$ with $Y\ge4$, $m(z)\le5Y$.
\end{lemma}
\begin{proof}
By homogeneity take $\tilde z=i$. In normal coordinates $\xi$ at $i$ put $v(\xi)=\tilde u(\exp_i(r\xi))$ for $|\xi|<1$. Then
$r^4(\Delta^2\tilde u)(\exp_i(r\xi))=L_r^2v(\xi)$, where $L_r$ is $r^2\Delta$ in these coordinates: a uniformly elliptic operator whose coefficients are
bounded in $C^\infty$ uniformly for $0<r\le1$. Interior regularity for $L_r^2$ and $H^4\subset L^\infty$ give
$|v(0)|\le C(\norm v_{L^2(B_1)}+\norm{L_r^2v}_{L^2(B_1)})$. Since $d\nu\asymp r^2d\xi$ on the ball,
$|\tilde u(i)|\le C(r^{-1}\norm{\tilde u}_{L^2(B(i,r))}+r^3\norm{\Delta^2\tilde u}_{L^2(B(i,r))})$, and $B(\tilde z,r)\subset B(\tilde z,1)$, so the unfolding step of
Lemma~\ref{pc:g43} applies. For the count: if $\gamma\notin\Gamma_\infty$ then $\operatorname{Im}\gamma z\le1/Y$ and $d(z,\gamma z)\ge2\log Y\ge2$; and
$d(z,z+n)=2\operatorname{arcsinh}(|n|/(2Y))<2$ if and only if $|n|<2Y\sinh1$. So $m(z)\le4Y\sinh1+1\le5Y$ for $Y\ge4$.
\end{proof}

\begin{lemma}[cusp asymptotics]\label{cl:asym}
Let $\vartheta$ be either $X_0^{-2}$ or $\ue^{-d(i,\cdot)}(1+d(i,\cdot))^{-3/2}$, and for $A\in C^4(\Hb)$ put
$B(A):=\sup_w\vartheta(w)^{-1}\sum_{|S|\le4}|Z_SA(w)|$. There is a $C$ (depending only on the choice of $\vartheta$) such that for $0<\hbar\le\tfrac12$,
$t\ge2\log(1/\hbar)$, $Y=\ue^t$ and $z=x+iY$,
\begin{multline*}
\Bigl|\bigl[(S(t)A)\ue^{i\varphi/\hbar}\bigr]_\Gamma(z)-\ue^{t/2}\sqrt{\pi\hbar}\,\ue^{i\pi/4}\ue^{it/\hbar}\sum_{c\ge1}c^{2i/\hbar}\sum_{a\in\mathbb Z,(a,c)=1}A\Bigl(\frac ac+\frac i{c^2}\Bigr)\Bigr|\\
\le C\,\ue^{t/2}\hbar^{3/2}B(A),
\end{multline*}
and the double series converges absolutely.
\end{lemma}
\begin{proof}
\emph{Cosets.} An element of $\mathrm{PSL}(2,\mathbb Z)$ is a translation $T^k$, or has lower-left entry $c\ge1$ and acts by $w\mapsto a/c-1/(c(cw+d))$ with
$ad\equiv1\pmod c$. For fixed $c$ and $a$ coprime to $c$, $d$ runs over $d_a+c\mathbb Z$, and $d=d_a+cn$ corresponds to $w=z+n$. Hence
$[F]_\Gamma(z)=\sum_kF(z+k)+\sum_{c\ge1}\sum_{(a,c)=1}\sum_{n\in\mathbb Z}G_{a,c}(n)$ with $G_{a,c}(u)=F(\gamma_{a,c}(z+u))$. The translation part is
bounded by $\ue^{-t/2}B(A)\sum_k\vartheta(x+k+i\ue^{2t})\le C\ue^{-t/2}B(A)$, which is $\le C\ue^{t/2}\hbar^{3/2}B(A)$.

\emph{Change of variable.} With $\tau=(c(x+u)+d_a)/Y$, the point $D_t\gamma_{a,c}(z+u)$, where $D_t(x,y)=(x,\ue^ty)$ and $S(t)A=\ue^{-t/2}A\circ D_t$, is
\[
Q(\tau)=\frac ac-\frac{\tau}{cY(\tau^2+c^2)}+\frac{i}{\tau^2+c^2},\qquad \varphi(\gamma_{a,c}(z+u))=t+\log(\tau^2+c^2),
\]
so $G_{a,c}(u)=g(\tau):=\ue^{-t/2}A(Q(\tau))\,\ue^{i(t+\log(\tau^2+c^2))/\hbar}$ and $du=(Y/c)\,d\tau$. Along $Q$,
$\partial_\tau A(Q)=-(YA)(Q)\,\frac{2\tau}{\tau^2+c^2}-(XA)(Q)\,\frac{c^2-\tau^2}{cY(\tau^2+c^2)}$, so $|\partial_\tau^jA(Q)|\le C_jc^{-j}B(A)\vartheta(Q)$ for
$j\le4$. From $X_0(Q)=(q^2+1)/(2\operatorname{Im}Q)+\operatorname{Im}Q/2$ with $|q-a/c|\le1/(2c^2Y)$ one gets
$\vartheta(Q(\tau))\le C\,\vartheta_{a,c}\,(1+\tau^2/c^2)^{-s}$, where $s=2$, $\vartheta_{a,c}=(a^2+c^2)^{-2}$ for $\vartheta=X_0^{-2}$, and $s=1$,
$\vartheta_{a,c}=(a^2+c^2)^{-1}(1+\log(a^2+c^2))^{-3/2}$ for the second choice. In both cases $\sum_{(a,c)}\vartheta_{a,c}<\infty$ (for the second, by
partial summation, since $\#\{(a,c):a^2+c^2\le N\}\le CN$).

\emph{Poisson summation.} $\sum_nG(n)=\sum_m\hat G(m)$ with
\[
\hat G(m)=\frac Yc\,\ue^{2\pi im(x+d_a/c)}\int g(\tau)\,\ue^{-2\pi imY\tau/c}\,d\tau .
\]
For $m\ne0$ the total phase $\Phi$ has $|\Phi'|\ge2\pi|m|Y/c-1/(\hbar c)\ge\pi|m|Y/c$ (as $\hbar Y\ge\hbar^{-1}\ge2$), $|\Phi''|\le2/(\hbar c^2)$ and
$|\Phi'''|\le C/(\hbar c^3)$. Two integrations by parts give
\[
\Bigl|\int g\,\ue^{-2\pi imY\tau/c}d\tau\Bigr|\le C\Bigl(\frac c{|m|Y}\Bigr)^2c^{-2}\ue^{-t/2}B(A)\int\vartheta(Q(\tau))\,d\tau\le C\ue^{-t/2}B(A)\,\vartheta_{a,c}\,\frac c{m^2Y^2},
\]
using $\int(1+\tau^2/c^2)^{-s}d\tau\le Cc$. Hence $\sum_{m\ne0}|\hat G(m)|\le C\ue^{-t/2}B(A)\vartheta_{a,c}/Y$, and the sum over $(a,c)$ is $\le C\ue^{-3t/2}B(A)$.

\emph{Stationary phase.} With $\tau=c\sigma$, $\hat G(0)=Y\ue^{it/\hbar}c^{2i/\hbar}\int f(\sigma)\ue^{i\theta(\sigma)/\hbar}d\sigma$, where
$f(\sigma)=\ue^{-t/2}A(Q(c\sigma))$ and $\theta(\sigma)=\log(1+\sigma^2)$. Let $\chi$ be a cutoff, $\chi=1$ on $[-\tfrac12,\tfrac12]$, $\operatorname{supp}\chi\subset[-1,1]$.
Near part: $\theta=s(\sigma)^2$ with $s(\sigma)=\operatorname{sign}(\sigma)(\log(1+\sigma^2))^{1/2}$ smooth and $s'(0)=1$; for $h(s)=(\chi f)(\sigma(s))\sigma'(s)$,
\[
\int h\,\ue^{is^2/\hbar}ds=\sqrt{\pi\hbar}\,\ue^{i\pi/4}\,\frac1{2\pi}\int\hat h(\xi)\,\ue^{-i\hbar\xi^2/4}d\xi,
\]
and $|\ue^{-i\hbar\xi^2/4}-1|\le\hbar\xi^2/4$ gives the error $C\hbar^{3/2}\norm h_{C^4}\le C\hbar^{3/2}\ue^{-t/2}B(A)\vartheta_{a,c}$. Far part, on $|\sigma|\ge\tfrac12$, where $|\theta'|\ge(1+|\sigma|)^{-1}/2$ and
$|\theta^{(j)}|\le C(1+|\sigma|)^{-j}$. First truncate: with a cutoff $\chi_Y$ equal to $1$ on $|\sigma|\le Y$ and $0$ on $|\sigma|\ge2Y$, the integral over
the support of $1-\chi_Y$ is bounded absolutely by $C\ue^{-t/2}B(A)\vartheta_{a,c}/Y\le C\hbar^2\ue^{-t/2}B(A)\vartheta_{a,c}$. On $\tfrac12\le|\sigma|\le2Y$
integrate by parts once with $D^Tu=-(u/\theta')'$; the derivatives of $\chi_Y$ add terms of order $\hbar/Y$. In $\partial_\sigma f$ the $Y$-terms carry
$2\sigma/(1+\sigma^2)$ and gain $(1+|\sigma|)^{-1}$ per derivative, and the $X$-terms carry $1/Y$. The $X$-part of $D^T[(1-\chi)\chi_Yf]$ is integrated directly,
giving at most $C\hbar\,\ue^{-t/2}B(A)\vartheta_{a,c}Y^{-1}\log Y\le C\hbar^{3/2}\ue^{-t/2}B(A)\vartheta_{a,c}$ (as $Y^{-1}\log Y=t\ue^{-t}\le\hbar^{1/2}$); the
remaining symbol part is integrated by parts once more, giving at most $C\hbar^2\ue^{-t/2}B(A)\vartheta_{a,c}$. Since $f(0)=\ue^{-t/2}A(a/c+i/c^2)$,
\[
\hat G(0)=\ue^{t/2}\sqrt{\pi\hbar}\,\ue^{i\pi/4}\ue^{it/\hbar}c^{2i/\hbar}A\Bigl(\frac ac+\frac i{c^2}\Bigr)+O\bigl(\ue^{t/2}\hbar^{3/2}B(A)\vartheta_{a,c}\bigr).
\]
Summing over $(a,c)$ proves the lemma.
\end{proof}

\begin{lemma}[the Poisson kernel off the diagonal]\label{cl:kernel}
Fix $r_0\in(0,1]$. For $0<\alpha\le\alpha_0<\pi/2$ the kernel $k_\alpha(r)$ of $\ue^{-\alpha\Lambda}$ extends holomorphically to
$\{a+i\theta:a>3r_0/4,\ |\theta|<r_0/4\}$ with $|\tilde k_\alpha(a+i\theta)|\le C\alpha\,\ue^{-a}(1+a)^{-3/2}$. Hence
$|\partial_r^jk_\alpha(r)|\le C_j\alpha\,\ue^{-r}(1+r)^{-3/2}$ for $r\ge r_0$, uniformly in $\alpha\le\alpha_0$.
\end{lemma}
\begin{proof}
By Lemma~\ref{os:hk} with $\theta_0=r_0/4$,
\[
|\tilde k_\alpha(a+i\theta)|\le C\int_0^\infty\alpha(4\pi)^{-1/2}s^{-3/2}\ue^{-(\alpha^2-\theta^2)/(4s)}p_s(a)\,ds .
\]
The computation in the proof of Lemma~\ref{os:poisson}, with $\alpha^2-\theta^2$ in place of $\alpha^2$, turns this into
\[
C\alpha\int_a^\infty\frac{\rho\,K_2\bigl(\sqrt{\rho^2+\alpha^2-\theta^2}/2\bigr)}{\rho^2+\alpha^2-\theta^2}\,\frac{d\rho}{(\cosh\rho-\cosh a)^{1/2}},
\]
which converges because $\rho^2+\alpha^2-\theta^2\ge r_0^2/2$, and is bounded as in \textup{(\ref{os:ub})}. Cauchy estimates on discs of radius $r_0/8$ give the derivative bounds.
\end{proof}

\begin{lemma}[the amplitudes]\label{cl:amp}
(i) $a=\ue^{-\kappa X_0}$ has finite swapped norm for every $\alpha\le r_1<\tfrac12$ and $\beta$, and $|Z_Sa|\le C_{S,\kappa}X_0^{|S|}\ue^{-\kappa X_0}$.
(ii) $a=\ue^{-\alpha\Lambda}\ue^{-\kappa X_0}$ has $\norm a_{\alpha,\beta}=\norm{\ue^{\beta\langle d\rangle}\ue^{-\kappa X_0}}<\infty$, and
$|Z_Sa(w)|\le C_{S,\kappa}\ue^{-d(i,w)}(1+d(i,w))^{-3/2}$ uniformly in $\alpha\le\alpha_0$.
(iii) Put $A_0=(Y^2-Y)a$ and $S_0(\hbar)=\sum_{c\ge1}c^{2i/\hbar}\sum_{(a,c)=1}A_0(a/c+i/c^2)$. In case (i) with $\kappa=20$, $|S_0(\hbar)|\ge0.98\,\kappa\ue^{-\kappa}$
for all $\hbar$; in case (ii), $|S_0(\hbar)|\ge\tfrac12\kappa\ue^{-\kappa}$ for $\kappa\ge\kappa_*$, $\alpha\le c\kappa^{-6}$ and all $\hbar$.
\end{lemma}
\begin{proof}
(i) $f=\ue^{\beta\langle d\rangle}a$ continues to the crown with $\operatorname{Re}\tilde X_0=\cos(2r)X_0$ and $\operatorname{Re}\widetilde{\langle d\rangle}\le\langle d\rangle+C_Hr^2$ (proof of
Lemma~\ref{lem:ex}), so its tube norm is finite and $\norm{\ue^{\alpha\Lambda}f}\le C_0N_\alpha(f)$ as there. The vector fields act by $XX_0=x$, $YX_0=y-X_0$,
$Xx=y$, $Yy=y$, $Xy=Yx=0$, and $|x|,y\le2X_0$; this gives the word bound.
(ii) The norm is the definition. For the words, transport $i$ to $w'$ by the affine group, under which $X$ and $Y$ are invariant; then
$|Z_S^w\cosh d(w,w')|\le C_S\cosh d(w,w')$, and Lemma~\ref{cl:kernel} gives $|Z_S^wk_\alpha(d(w,w'))|\le C_S\alpha\ue^{-d(w,w')}(1+d(w,w'))^{-3/2}$ for
$d(w,w')\ge\tfrac14$. Split $h=\ue^{-\kappa X_0}$ near each $w$ as $h\psi_w+h(1-\psi_w)$, $\psi_w=1$ on $B(w,\tfrac18)$, $\operatorname{supp}\psi_w\subset B(w,\tfrac14)$:
the far part is bounded with the kernel estimate, and the near part (with $\psi_w$ frozen at the evaluation point) by local Sobolev and
$\norm{\ue^{-\alpha\Lambda}g}_{H^k}\le\norm g_{H^k}$; the words of $h$ carry powers of $\kappa$ and $X_0$, so the near part is at most
$C_{S,\kappa}\exp(-\tfrac\kappa2\cosh(\max(0,d(i,w)-\tfrac14)))$, which is dominated by $\ue^{-d}(1+d)^{-3/2}$.
(iii) $Y^2X_0=X_0$ gives $A_0=a[\kappa^2(y-X_0)^2-\kappa(2X_0-y)]$ in case (i), so $A_0(i)=-\kappa\ue^{-\kappa}$. At the other points $a/c+i/c^2$,
$X_0=(a^2+c^2+c^{-2})/2$, $X_0-1\ge\max(\tfrac12,(n-2)/2)$ with $n=a^2+c^2$, and $|A_0|\le\kappa\ue^{-\kappa}\ue^{-\kappa(X_0-1)}(\kappa X_0^2+2X_0)$. For
$\kappa=20$ and $q=\ue^{-10}<1/20000$: the two points $\pm1+i$ give at most $96q<0.005$; there is no coprime pair with $c\ge1$ and $n\in\{3,4\}$; the rest is
at most $\sum_{n\ge5}(4n+4)[5(n+1)^2+n+1]q^{n-2}\le5184q^3/(1-8q)<0.01$. In case (ii), $|A_0(i)+\kappa\ue^{-\kappa}|\le C\norm{(\ue^{-\alpha\Lambda}-1)h}_{H^4}\le
C\alpha\norm{(1-\Delta)^3h}\le C\alpha\kappa^6\ue^{-\kappa}$. The other points are $0.96$-separated ($\cosh d\ge\tfrac32$ for distinct ones). Split
$h=h_1+h_2$ with $h_1=h\psi$, $\psi=1$ on $B(i,0.44)$, $\operatorname{supp}\psi\subset B(i,0.45)$, so that $h_2$ vanishes on $B(i,0.44)$: $h_1$ contributes at most $C\alpha\kappa^{-1}\ue^{-\kappa}\vartheta(w)$ at each point by Lemma~\ref{cl:kernel}; for $h_2$, whose
derivatives up to order $4$ are at most $C(1+\kappa)^4(1+X_0)^4\ue^{-\kappa X_0}$, the split of (ii) at each point and the packing bound
$\#\{w:d(i,w)\le R\}=O(\ue^R)$ give at most $C\kappa^8\ue^{-\kappa\cosh0.44}$ in total, with $C$ independent of $\kappa\ge1$ and $\alpha\le\alpha_0$. Hence the sum
over the other points is at most $C(\alpha/\kappa^2+\kappa^8\ue^{-0.098\kappa})\kappa\ue^{-\kappa}$.
\end{proof}

\begin{proof}[Proof of Theorem~\ref{cl:main}]
With $c_\hbar=i\hbar/2$, the definition of $u^{(K)}$ gives $E_0=M_1+M_2+E_2$, where
$M_k=\ue^{-it/(2\hbar)}[(S(t)A_k)\ue^{i\varphi/\hbar}]_\Gamma$, $A_1=c_\hbar P_1(t)a$ and $A_2=c_\hbar^2P_2(t)a$. Since
$P_1(t)a=\int_0^t\Delta(s)a\,ds=\tfrac{1-\ue^{-2t}}2X^2a+t(Y^2-Y)a$, Lemma~\ref{cl:asym} gives
\[
|M_1(z)|\ge\ue^{t/2}\sqrt{\pi\hbar}\,\tfrac\hbar2\bigl(t|S_0(\hbar)|-S_X\bigr)-C\ue^{t/2}\hbar^{3/2}\hbar t,
\]
with $S_X=\sum_{(a,c)}|X^2a(a/c+i/c^2)|<\infty$. $P_2(t)a$ is a combination of
words of length at most $4$ with coefficients bounded by $t^2/2$, so $|M_2(z)|\le C\ue^{t/2}\hbar^{1/2}\hbar^2t^2$. For $E_2$, Theorem~\ref{thm:dd-appx} in case (i)
and Theorem~\ref{pc:main} in case (ii), at $K=2$ and $N=0,2$, give $\norm{E_2}\le C(\hbar t/\alpha^2)^3\norm a$ and
$\norm{\Delta^2E_2}\le C(\alpha\hbar)^{-4}(\hbar t/\alpha^2)^3\norm a$; Lemma~\ref{cl:ssob} with $r=\hbar$ gives $|E_2(z)|\le C\ue^{t/2}\hbar^{-1}\alpha^{-10}(\hbar t)^3\norm a$.
Hence
\[
|E_0(z)|\ge\ue^{t/2}\hbar^{3/2}t\Bigl[\tfrac{\sqrt\pi}2|S_0(\hbar)|-\tfrac{\sqrt\pi}2\frac{S_X}t-C\hbar-C\hbar t-C\hbar^{1/2}t^2\Bigr],
\]
and the bracket is at least $\tfrac{\sqrt\pi}4|S_0(\hbar)|$ for $t\ge t_*$, $t\le\hbar^{-1/5}$ and $\hbar\le\hbar_0$. Lemma~\ref{cl:amp}(iii) bounds $|S_0|$ below.
\end{proof}


% Main-text section: expectation values.  Input file; no preamble.  Labels prefix ev.
\section{Expectation values}\label{sec:expect}

The theorems compare the solution with the periodized approximation in $L^2(M)$. This section records what follows for expectation
values, and shows on the modular surface that accuracy of the approximation and classical behaviour of expectation values are
different things: the approximation stays accurate up to times of order $\hbar^{-1/2}$, while the classical transport of the
initial density fails to describe the expectation values after a time of order $\log(1/\hbar)\log\log(1/\hbar)$.

\subsection{Transfer and cancellation}\label{subsec:ev-transfer}
\begin{corollary}\label{ev:transfer}
Assume the hypotheses of Theorem~\ref{thm:dd14}, or those of Theorem~\ref{pc:main} with $\pcNorm a$ in place of $\norm a$. Let
$K\in\mathbb N_0$ and $0\le t\le\delta(\alpha/\hbar)^{1/2}$, and put
\[
  \psi(t):=\cU(t)\big[a\,\ue^{i\varphi_b/\hbar}\big]_\Gamma,\qquad
  \epsilon_K:=C^{K+1}K!\,\Big(\frac{\hbar t}{\alpha^2}\Big)^{K+1}\norm a .
\]
\begin{itemize}
\item[(i)] For every bounded operator $B$ on $L^2(M)$,
\[
  \big|\langle B\psi(t),\psi(t)\rangle-\langle Bu^{(K)}(t)_\Gamma,u^{(K)}(t)_\Gamma\rangle\big|
  \le\norm B\,\epsilon_K\big(2\norm{\psi(0)}+\epsilon_K\big).
\]
\item[(ii)] $\big|\,\norm{u^{(K)}(t)_\Gamma}_{L^2(M)}-\norm{[a\,\ue^{i\varphi_b/\hbar}]_\Gamma}_{L^2(M)}\big|\le\epsilon_K$.
\end{itemize}
\end{corollary}
\begin{proof}
The theorem at $N=0$ gives $\norm{\psi(t)-u^{(K)}(t)_\Gamma}\le\epsilon_K$, and $\norm{\psi(t)}=\norm{\psi(0)}$ since $\cU(t)$ is unitary.
This is (ii). For (i) write $u=u^{(K)}(t)_\Gamma$ and
$\langle B\psi,\psi\rangle-\langle Bu,u\rangle=\langle B(\psi-u),\psi\rangle+\langle Bu,\psi-u\rangle$, and use
$\norm u\le\norm{\psi(0)}+\epsilon_K$.
\end{proof}

\begin{remark}[cancellation]\label{ev:cancel}
At $K=0$, (ii) says that $\norm{[(S(t)a)\,\ue^{i\varphi_b/\hbar}]_\Gamma}_{L^2(M)}$ stays within $C\hbar t\alpha^{-2}\norm a$ of its value at $t=0$,
for all $t\le\delta(\alpha/\hbar)^{1/2}$. Without the phase this fails: if $M$ has finite area and $a\in L^1(\Hb)$, then by the
Cauchy--Schwarz inequality and unfolding
\[
  \norm{[S(t)|a|]_\Gamma}_{L^2(M)}\ge\operatorname{area}(M)^{-1/2}\int_\Hb S(t)|a|\,d\nu=\operatorname{area}(M)^{-1/2}\,\ue^{t/2}\norm a_{L^1(\Hb)} .
\]
So the norm of the periodization with the phase is smaller than the norm of the periodization of the moduli by a factor of order
at most $\ue^{-t/2}$: the terms cancel. The following picture is geometric intuition and is not proved here. The terms that
contribute at a point come from about $\ue^{t}$ sheets of the cover. Up to the Ehrenfest time $\log(1/\hbar)$ their phase
gradients differ by more than $\hbar$, so that they are almost orthogonal one by one; beyond it the phase gradients of many pairs
differ by less than $\hbar$, and the cancellation is a property of the sum.
\end{remark}

By (i), for bounded observables the expectation values of $\psi(t)$ up to times $\delta(\alpha/\hbar)^{1/2}$ are those of the
explicit sum $u^{(K)}(t)_\Gamma$, up to an error that tends to $0$ with $\hbar$ for fixed $\alpha$ and $K$. The next subsection answers one such question in a case where the sum can be evaluated.

\subsection{The modular surface}\label{subsec:ev-mod}
In this subsection $\Gamma=\mathrm{PSL}(2,\mathbb Z)$, $M=\Gamma\backslash\Hb$, $x_0=i$, $b=\infty$ and $\varphi=-\log y$. Let
$F=\{|x|\le\tfrac12,\ |z|\ge1\}$ be the standard fundamental domain; for $Y\ge1$ let $M_Y\subset M$ be the image of $\{z\in F:y\le Y\}$,
a compact set, and $C_Y$ the image of $\{z\in F:y>Y\}=\{|x|\le\tfrac12,\ y>Y\}$, the part of the cusp above height $Y$. Recall
$\operatorname{area}(M)=\pi/3$.

Fix $\alpha\in(0,\tfrac12]$ and $h\in C_c^\infty(\Hb)$, $h\ge0$, such that
\[
  h_0(y):=\sum_{n\in\mathbb Z}h(x+n,y)
\]
does not depend on $x$ and is not identically zero; for instance $h(x,y)=\int_0^1h_1(x+s,y)\,ds$ with $h_1\in C_c^\infty(\Hb)$,
$h_1\ge0$, $h_1\not\equiv0$. Let
\[
  a:=\ue^{-\alpha\Lambda}h .
\]
Then $\pcNorm a=\norm{\ue^{\beta\langle d\rangle}h}<\infty$ for every $\beta$, so $a$ is an amplitude of the printed class, and it does
not depend on $\hbar$. Let $k_\alpha\ge0$ be the radial kernel of $\ue^{-\alpha\Lambda}$ (Lemma~\ref{os:poisson}). It is integrable on
$\Hb$, by \textup{(\ref{os:ub})} and \textup{(\ref{os:small})}. Since $k_\alpha(d(z,w))$ is invariant under $z\mapsto z+n$, $w\mapsto w+n$,
\begin{equation}\label{ev:a0}
  a_0(y):=\sum_{n\in\mathbb Z}a(x+n,y)=\int_\Hb k_\alpha(d(z,w))\,h_0(\operatorname{Im}w)\,d\nu(w)
\end{equation}
does not depend on $x$, and $a_0>0$. Put
\[
  \mu_0:=\int_0^\infty|a_0(y)|^2\,\frac{dy}{y^2}.
\]

\begin{proposition}[accurate approximation, non-classical expectation values]\label{ev:mod}
Let $\Gamma$, $\alpha$, $h$, $a$, $a_0$, $\mu_0$ be as above, $\psi(t)=\cU(t)[a\,\ue^{i\varphi/\hbar}]_\Gamma$, and let $C,\delta$ be the
constants of Theorem~\ref{pc:main} for $\beta=1$ and $\kappa=1$. Then $0<\mu_0<\infty$, and there is an absolute constant $C_1$ such
that, with $T(\hbar):=C_1\log(1/\hbar)\log\log(1/\hbar)$, the following hold.
\begin{itemize}
\item[(i)] $\norm{\psi(t)}^2_{L^2(M)}=\mu_0+O(\hbar^{n})$ for every $n$, and for all $K$ and $0\le t\le\delta(\alpha/\hbar)^{1/2}$,
$\norm{\psi(t)-u^{(K)}(t)_\Gamma}_{L^2(M)}\le C^{K+1}K!\,(\hbar t/\alpha^2)^{K+1}\pcNorm a$.
\item[(ii)] For every $Y\ge1$ and every $K$, as $\hbar\to0$,
\[
  \int_{M_Y}|\psi(t)|^2\,d\nu+\int_{M_Y}|u^{(K)}(t)_\Gamma|^2\,d\nu\longrightarrow0
  \qquad\text{uniformly for }T(\hbar)+\log Y\le t\le\delta(\alpha/\hbar)^{1/2}.
\]
\item[(iii)] For every continuous $q$ with compact support in $M$, lifted to a $\Gamma$-invariant function on $\Hb$,
\[
  \int_0^\infty\!\!\int_0^1q(z)\,\big|(S(t)a_0)(y)\big|^2\,dx\,\frac{dy}{y^2}\longrightarrow\frac3\pi\int_Mq\,d\nu\cdot\mu_0\qquad(t\to\infty).
\]
\end{itemize}
\end{proposition}
The left side of (iii) is the classical prediction for $\langle q\,\psi(t),\psi(t)\rangle$: the density $|a_0|^2$ on
$\Gamma_\infty\backslash\Hb$, transported by the geodesic flow for the time $t$ and projected to $M$; it is the sum of the squared
moduli of the terms of $u^{(0)}(t)_\Gamma$, summed over $\Gamma_\infty\backslash\Gamma$. By (iii) this prediction tends to the
equidistributed value. By (ii), for $q\ge0$ with support in $M_Y$,
$\langle q\,\psi(t),\psi(t)\rangle\le\sup q\int_{M_Y}|\psi(t)|^2\to0$ uniformly on the window of (ii): the $L^2$ mass of the state in every
compact set tends to $0$; the state has moved into the cusp. So on this window the expectation value of a compactly supported
$q\ge0$, $q\not\equiv0$, is not given by the classical transport, whose limit $\tfrac3\pi\mu_0\int_Mq\,d\nu$ is positive, although by (i) and
(ii) the periodized approximation remains accurate and has the same behaviour. The time $T(\hbar)$ is a sufficient time; it is not
claimed to be the time at which the departure begins. For the wave
group and Lagrangian states on unstable horocycles of the modular surface, the departure from classical transport at logarithmic
times is due to Brooks \cite{Bro17}.

\begin{remark}[the time scale]\label{ev:time}
The proof uses the classical zero-free region of the Riemann zeta function, which gives the sufficient time $T(\hbar)$. For $r>0$
and a continuous choice of the argument, the scattering coefficient $\Phi$ of \eqref{ev:Phi} satisfies
$-\tfrac{d}{dr}\arg\Phi(\tfrac12+ir)=2\log(r/\pi)+4\operatorname{Re}\tfrac{\zeta'}{\zeta}(1+2ir)+O(r^{-2})$. Heuristically this is the time delay of a wave
packet at $r=1/\hbar$, and its first term is twice the Ehrenfest time $\log(1/\hbar)$, up to a constant. The bound for $\zeta'/\zeta$ used here is
of the same order $\log r$ as that term, so this is an indication only: we do not prove a statement with the constant $2$, and we
do not prove that the classical prediction holds before that time.
\end{remark}

\subsubsection*{Proof of Proposition~\ref{ev:mod}}
Throughout, $n$ denotes an arbitrary positive integer and constants may depend on $n$, $\alpha$ and $h$.

\emph{Step 1: the function $a_0$.} Let $w$ be the bounded function $w(z)=h_0(\operatorname{Im}z)$ on $\Hb$. Since $k_\alpha\ge0$ is
integrable, \eqref{ev:a0} gives $0<a_0\le\norm{k_\alpha}_{L^1(\Hb)}\sup h_0$, and by Tonelli's theorem
\[
  \int_0^\infty a_0(y)\frac{dy}{y^2}=\int_0^\infty\!\!\int_0^1\sum_na(x+n,y)\,dx\frac{dy}{y^2}=\int_\Hb a\,d\nu\le\norm{k_\alpha}_{L^1}\norm h_{L^1}<\infty .
\]
Hence $\mu_0\le\sup a_0\int_0^\infty a_0\,dy/y^2<\infty$, and $\mu_0>0$. Convolution with a radial kernel commutes with the isometries
$z\mapsto\lambda z$, $\lambda>0$, whose generator acts on functions of $y$ alone as $y\partial_y$. So
$(y\partial_y)^ja_0$ is given by \eqref{ev:a0} with $(y\partial_y)^jh_0$ in place of $h_0$, and since
$|(y\partial_y)^jh_0|\le C_j\tilde h_0$ for a fixed $\tilde h_0\in C_c^\infty((0,\infty))$, $\tilde h_0\ge0$, that equals $1$ on the support of $h_0$, the same
bounds hold: with $\tilde a_0$ defined by \eqref{ev:a0} from $\tilde h_0$,
\begin{equation}\label{ev:a0bd}
  |(y\partial_y)^ja_0|\le C_j\tilde a_0,\qquad \sup\tilde a_0<\infty,\qquad\int_0^\infty\tilde a_0\,\frac{dy}{y^2}<\infty\qquad(j\ge0).
\end{equation}
For a function $f$ on $(0,\infty)$ let $\hat f(s):=\int_0^\infty f(y)y^{-s}\,dy/y$ be its Mellin transform. By \eqref{ev:a0bd},
$\hat a_0(s)$ converges absolutely for $0<\operatorname{Re}s\le1$, is holomorphic in $0<\operatorname{Re}s<1$ and continuous up to
$\operatorname{Re}s=1$. In the variable $u=\log y$ the functions $y^{-\sigma}(y\partial_y)^ja_0$, $\tfrac14\le\sigma\le1$, are integrable with integrable
derivatives, so they vanish at $\pm\infty$, and integration by parts gives $s^j\hat a_0(s)=\widehat{(y\partial_y)^ja_0}(s)$. Hence
\begin{equation}\label{ev:mellin}
  |\hat a_0(\sigma+i\rho)|\le C_n(1+|\rho|)^{-n}\qquad(\tfrac14\le\sigma\le1,\ \rho\in\mathbb R).
\end{equation}
The Mellin--Plancherel formula reads $\int_0^\infty|f|^2\,dy/y^2=\tfrac1{2\pi}\int_{\mathbb R}|\hat f(\tfrac12+ir)|^2\,dr$.

\emph{Step 2: incomplete Eisenstein series.} For a function $f$ on $(0,\infty)$ put $E(z|f):=\sum_{\gamma\in\Gamma_\infty\backslash\Gamma}f(\operatorname{Im}\gamma z)$,
where $\Gamma_\infty$ is the group of the translations $z\mapsto z+n$. The cosets are the identity and the classes of the
matrices with bottom row $(c,d)$, $c\ge1$, $\gcd(c,d)=1$, and $\operatorname{Im}\gamma z=y/|cz+d|^2$. Let $\mathcal F$ be the set of smooth $f$ with
$\int_0^\infty|f|\,dy/y^2<\infty$ and $\int_0^\infty|f|^2\,dy/y^2<\infty$. For $f\in\mathcal F$, unfolding gives
$\int_ME(z||f|)\,d\nu=\int_0^\infty|f|\,dy/y^2<\infty$; so the series converges absolutely almost everywhere and $E(\cdot|f)\in L^1(M)$. For $y_1>0$ the
set $\{0\le x\le1,\ y>y_1\}$ is covered by finitely many, say $\mathcal N(y_1)$, $\Gamma$-translates of $F$, with $\mathcal N(1)=2$. So
$\int_{y_1}^\infty\int_0^1E(z||f|)\,d\nu\le\mathcal N(y_1)\int_0^\infty|f|\,dy/y^2$, and for almost every $y>0$ the zero Fourier mode of $E(\cdot|f)$ is, by
Fubini's theorem and the substitution $d=d_0+cn$,
\begin{equation}\label{ev:zero}
  \int_0^1E(x+iy|f)\,dx=f(y)+R_f(y),\qquad R_f(y):=\sum_{c\ge1}N_c\int_{\mathbb R}f\Big(\frac{y}{c^2(x^2+y^2)}\Big)dx ,
\end{equation}
where $N_c$ is the number of residues $d_0$ modulo $c$ that are prime to $c$. The same computation with $|f|$ shows
$\int_{y_1}^\infty|R_f|\,dy/y^2\le\mathcal N(y_1)\int_0^\infty|f|\,dy/y^2$ for every $y_1>0$. Let
\begin{equation}\label{ev:Phi}
  \Phi(s):=\sqrt\pi\,\frac{\Gamma(s-\frac12)}{\Gamma(s)}\,\frac{\zeta(2s-1)}{\zeta(2s)}
  =\pi^{2s-1}\,\frac{\Gamma(1-s)}{\Gamma(s)}\,\frac{\zeta(2-2s)}{\zeta(2s)} ,
\end{equation}
the two forms being equal by the functional equation of $\zeta$. $\Phi$ is holomorphic in $\operatorname{Re}s\ge\tfrac12$ except for a simple
pole at $s=1$ with residue $3/\pi$, $|\Phi(\tfrac12+ir)|=1$ (because $\Phi(s)\Phi(1-s)=1$ and $\overline{\Phi(s)}=\Phi(\bar s)$), and $\Phi$ is bounded by a power of $1+|s|$ in
$\tfrac12\le\operatorname{Re}s\le2$, $|s-1|\ge\tfrac14$: the Gamma quotient by Stirling's formula \cite[5.11.9]{DLMF}, $\zeta(2s-1)$ by its polynomial
growth in vertical strips, and $1/\zeta(2s)$ by \cite[Theorem~6.7]{MV07}.
We claim that for $f\in\mathcal F$, for almost every $y>0$,
\begin{equation}\label{ev:R}
  R_f(y)=\frac3\pi\hat f(1)+\mathcal O_f(y),\qquad \mathcal O_f(y):=\frac1{2\pi}\int_{\mathbb R}\Phi(\tfrac12+ir)\hat f(\tfrac12+ir)\,y^{1/2-ir}\,dr ,
\end{equation}
where the integral is the $L^2$ inverse Mellin transform, so that
\begin{equation}\label{ev:Onorm}
  \int_0^\infty|\mathcal O_f|^2\,\frac{dy}{y^2}=\int_0^\infty|f|^2\,\frac{dy}{y^2}.
\end{equation}
For $f\in C_c^\infty((0,\infty))$, $\hat f$ is entire and rapidly decreasing on vertical lines, and Mellin inversion on $\operatorname{Re}s=2$ gives
\[
  R_f(y)=\frac1{2\pi i}\int_{(2)}\hat f(s)\Big(\sum_{c\ge1}\frac{N_c}{c^{2s}}\Big)\Big(\int_{\mathbb R}\frac{y^s\,dx}{(x^2+y^2)^s}\Big)ds
  =\frac1{2\pi i}\int_{(2)}\hat f(s)\,\Phi(s)\,y^{1-s}\,ds ,
\]
because $\sum_cN_cc^{-2s}=\zeta(2s-1)/\zeta(2s)$ and $\int_{\mathbb R}(1+v^2)^{-s}dv=\sqrt\pi\,\Gamma(s-\tfrac12)/\Gamma(s)$. Moving the line to
$\operatorname{Re}s=\tfrac12$ gives \eqref{ev:R}. For general $f\in\mathcal F$ take $f_j=f\chi(j^{-1}\log y)$ with $\chi\in C_c^\infty(\mathbb R)$, $0\le\chi\le1$, $\chi=1$ near $0$:
then $f_j\to f$ in $L^1$ and in $L^2$ of $dy/y^2$, so $R_{f_j}\to R_f$ in $L^1((y_1,\infty),dy/y^2)$ for every $y_1>0$, $\hat f_j(1)\to\hat f(1)$ and
$\mathcal O_{f_j}\to\mathcal O_f$ in $L^2(dy/y^2)$, and \eqref{ev:R} follows.

\emph{Step 3: the approximation as an incomplete Eisenstein series.} Let
\[
  f_t(y):=(S(t)a_0)(y)\,y^{-i/\hbar}=\ue^{-t/2}a_0(\ue^ty)\,y^{-i/\hbar}.
\]
By \eqref{ev:a0bd}, $f_t\in\mathcal F$, with $\int_0^\infty|f_t|^2dy/y^2=\mu_0$ and $\int_0^\infty|f_t|\,dy/y^2=\ue^{t/2}\int_0^\infty a_0\,dy/y^2$. By Corollary~\ref{os:cwp}, $\omega a=\omega\,\ue^{-\alpha\Lambda}h\in L^2(\Hb)$ for the periodization weight $\omega$ of
Section~\ref{sec:printed}. So by Proposition~\ref{os:p1} and (W1), (W2) of Section~\ref{subsec:dd-surf}, the orbit sums of $S(t)a$ and of
$(S(t)a)\ue^{i\varphi/\hbar}$ converge absolutely almost everywhere and lie in $L^2(M)$. They may be summed over $\Gamma_\infty$ first, so
$E(\cdot||f_t|)=[S(t)a]_\Gamma\in L^2(M)$ and
\[
  \big[(S(t)a)\ue^{i\varphi/\hbar}\big]_\Gamma=E(\cdot|f_t),\qquad u^{(0)}(t)_\Gamma=\ue^{-it/(2\hbar)}E(\cdot|f_t),\qquad\psi(0)=E(\cdot|f_0).
\]
The substitution $\eta=\ue^ty$ gives
\begin{equation}\label{ev:ft}
  \hat f_t(s)=\ue^{t(s-1/2)}\,\ue^{it/\hbar}\,\hat a_0(s+i/\hbar)\qquad(0<\operatorname{Re}s\le1).
\end{equation}
Theorem~\ref{pc:main} with $\beta=\kappa=1$ applies for $\hbar\le\alpha^3$; write $\epsilon_0=C\hbar t\alpha^{-2}\pcNorm a$. For
$t\le\delta(\alpha/\hbar)^{1/2}$ we have $\epsilon_0\le C\delta\alpha^{-3/2}\hbar^{1/2}\pcNorm a\to0$, and the second statement of (i) is the theorem.

\emph{Step 4: the norm and the mean.} Since $E(\cdot||f_0|)\in L^2(M)$, unfolding one factor of $|E(\cdot|f_0)|^2$ is justified, and with
\eqref{ev:zero} and \eqref{ev:R} it gives
\[
  \norm{\psi(0)}^2=\int_0^\infty f_0\,\overline{\Big(f_0+\frac3\pi\hat f_0(1)+\mathcal O_{f_0}\Big)}\,\frac{dy}{y^2}
  =\mu_0+\frac3\pi|\hat f_0(1)|^2+\int_0^\infty f_0\,\overline{\mathcal O_{f_0}}\,\frac{dy}{y^2},
\]
where the three integrals converge absolutely because $f_0\in\mathcal F$ and $\mathcal O_{f_0}\in L^2(dy/y^2)$, and $\int_0^\infty f_0\,dy/y^2=\hat f_0(1)$.
By \eqref{ev:ft} and \eqref{ev:mellin}, $|\hat f_0(1)|=|\hat a_0(1+i/\hbar)|\le C_n\hbar^n$. The Mellin transform of $\mathcal O_{f_0}$ at
$\tfrac12+ir$ is $\Phi(\tfrac12-ir)\hat f_0(\tfrac12-ir)$, so by the Parseval formula and \eqref{ev:ft} the last integral is at most
$\tfrac1{2\pi}\int|\hat a_0(\tfrac12+i(r+\tfrac1\hbar))|\,|\hat a_0(\tfrac12+i(-r+\tfrac1\hbar))|\,dr\le C_n\hbar^n$: at every $r$ one of the two arguments has
imaginary part at least $1/\hbar$ in modulus, and \eqref{ev:mellin} is used with the power $n+2$ for that factor and with the power $2$ for the
other, which is integrable in $r$. Hence $\norm{\psi(t)}^2=\norm{\psi(0)}^2=\mu_0+O(\hbar^n)$, which is the first statement of (i).

Since $M$ has finite area, the constant function $1$ lies in $L^2(M)$ and $\cU(t)1=1$. So $\langle\psi(t),1\rangle=\langle\psi(0),1\rangle=\hat f_0(1)$ and
$\langle E(\cdot|f_t),1\rangle=\hat f_t(1)$ (unfolding), and by the theorem
\begin{equation}\label{ev:mean}
  |\hat f_t(1)|\le|\hat f_0(1)|+\operatorname{area}(M)^{1/2}\epsilon_0\le C_n\hbar^n+2\epsilon_0\qquad(0\le t\le\delta(\alpha/\hbar)^{1/2}).
\end{equation}

\emph{Step 5: the outgoing term below height $Y$.} We bound $\int_0^Y|\mathcal O_{f_t}|^2dy/y^2$. Let $c_0\in(0,1]$ be strictly smaller than the constant $c$ of \cite[Theorem~6.6]{MV07}. Then the closed region
$\operatorname{Re}w\ge1-c_0/(2\log(|\operatorname{Im}w|+4))$ has a neighbourhood in which $\zeta(w)\ne0$, and in this region
$|\log\zeta(w)|\le\log\log(|\operatorname{Im}w|+4)+C$ where $|\operatorname{Im}w|\ge\tfrac78$, and $\zeta(w)(w-1)$ is bounded above and away from $0$ where
$|\operatorname{Im}w|\le\tfrac78$ \cite[Theorems~6.6 and~6.7]{MV07}. Put
\[
  \vartheta(r):=\frac{c_0}{4\log(2|r|+4)}\in(0,\tfrac14),\qquad D:=\{s=\sigma+ir:\ \tfrac12-\vartheta(r)\le\sigma\le\tfrac12\}.
\]
For $s\in D$ the points $2s$ and $2-2s$ lie in that region. So for $|r|\ge1$, $|\zeta(2-2s)|$ and $|1/\zeta(2s)|$ are at most
$C\log(|r|+3)$; and $|\Gamma(1-s)/\Gamma(s)|\le C(1+|r|)^{1-2\sigma}\le C(1+|r|)^{2\vartheta(r)}\le C'$ by Stirling's formula \cite[5.11.9]{DLMF}. For
$|r|\le1$, $\Phi$ is holomorphic on a neighbourhood of the compact set $D\cap\{|r|\le1\}$: $\zeta(2s)$ has no zero there, and in the
quotient $\zeta(2-2s)/\zeta(2s)$ the two poles at $s=\tfrac12$ cancel. So by the second form in \eqref{ev:Phi}, $\Phi$ is holomorphic on a neighbourhood of $D$ and
\begin{equation}\label{ev:PhiD}
  |\Phi(s)|\le C\log^2(|r|+3)\qquad(s\in D).
\end{equation}
By \eqref{ev:ft} and \eqref{ev:mellin}, on $D$ the function $\hat f_t$ is holomorphic and
$|\hat f_t(\sigma+ir)|\le\ue^{-t(1/2-\sigma)}C_n(1+|r+\tfrac1\hbar|)^{-n}$. So in the integral that defines $\mathcal O_{f_t}(y)$, written as
$\tfrac1{2\pi i}\int_{(1/2)}\Phi(s)\hat f_t(s)y^{1-s}ds$, the line may be moved to the left edge $\sigma=\tfrac12-\vartheta(r)$ of $D$, a piecewise smooth curve: for $f_t$ the integral converges
absolutely and defines the $L^2$ inverse Mellin transform pointwise, and the integrals over the horizontal segments $\{\sigma+ir:\tfrac12-\vartheta(r)\le\sigma\le\tfrac12\}$
tend to $0$ as $|r|\to\infty$, by \eqref{ev:PhiD} and the decay of $\hat f_t$. On that edge $|y^{1-s}|=y^{1/2+\vartheta(r)}$, the
norm of $y^{1/2+\vartheta}$ in $L^2((0,Y),dy/y^2)$ is $Y^{\vartheta}(2\vartheta)^{-1/2}$, and $|ds|\le2\,dr$. Minkowski's inequality gives, for $t\ge\log Y$,
\[
  \Big(\int_0^Y|\mathcal O_{f_t}|^2\frac{dy}{y^2}\Big)^{1/2}
  \le C_n\int_{\mathbb R}\log^{5/2}(|r|+4)\,\ue^{-(t-\log Y)\vartheta(r)}\,(1+|r+\tfrac1\hbar|)^{-n}\,dr .
\]
On $|r+\tfrac1\hbar|\le\tfrac1{2\hbar}$ we have $\vartheta(r)\ge\vartheta_\hbar:=c_0/(4\log(3/\hbar+4))$; on the rest the exponential is at most $1$ and
the integral is $O(\hbar^{n-2})$. So
\begin{equation}\label{ev:out}
  \Big(\int_0^Y|\mathcal O_{f_t}|^2\frac{dy}{y^2}\Big)^{1/2}\le C\log^{5/2}(4/\hbar)\,\ue^{-(t-\log Y)\vartheta_\hbar}+C_n\hbar^{n-2}.
\end{equation}
With $C_1:=13/c_0$ and $t-\log Y\ge T(\hbar)$, the exponent is at least $\tfrac{13}{4}\log\log(1/\hbar)\cdot\log(1/\hbar)/\log(3/\hbar+4)$, and the right side of
\eqref{ev:out} tends to $0$ as $\hbar\to0$, uniformly in such $t$ and in $Y\ge1$.

\emph{Step 6: the mass in the cusp.} Let $Y\ge1$ and $T(\hbar)+\log Y\le t\le\delta(\alpha/\hbar)^{1/2}$, and write $v=E(\cdot|f_t)$. By the Cauchy--Schwarz
inequality in $x$, \eqref{ev:zero} and \eqref{ev:R},
\begin{align*}
  \Big(\int_{C_Y}|v|^2d\nu\Big)^{1/2}&\ge\Big(\int_Y^\infty\Big|\int_0^1v\,dx\Big|^2\frac{dy}{y^2}\Big)^{1/2}\\
  &\ge\norm{\mathcal O_{f_t}}_{L^2((Y,\infty))}-\norm{f_t}_{L^2((Y,\infty))}-\frac3\pi|\hat f_t(1)|\,Y^{-1/2},
\end{align*}
with the norms of $L^2(dy/y^2)$. By \eqref{ev:Onorm} and \eqref{ev:out}, $\norm{\mathcal O_{f_t}}^2_{L^2((Y,\infty))}=\mu_0-o(1)$. Next
$\norm{f_t}^2_{L^2((Y,\infty))}=\int_{\ue^tY}^\infty|a_0|^2d\eta/\eta^2\le\sup|a_0|^2\,\ue^{-t}\to0$, and $|\hat f_t(1)|\to0$ by \eqref{ev:mean}. So
$\int_{C_Y}|v|^2d\nu\ge\mu_0-o(1)$. By Corollary~\ref{ev:transfer}(ii) and Step 4, $\norm v^2_{L^2(M)}\le(\norm{\psi(0)}+\epsilon_0)^2=\mu_0+o(1)$. Hence
\[
  \int_{M_Y}|u^{(0)}(t)_\Gamma|^2d\nu=\norm v^2-\int_{C_Y}|v|^2d\nu=o(1),
\]
uniformly on the window. Then $\norm{\mathbf 1_{M_Y}\psi(t)}\le\norm{\mathbf 1_{M_Y}u^{(0)}(t)_\Gamma}+\epsilon_0$ and
$\norm{\mathbf 1_{M_Y}u^{(K)}(t)_\Gamma}\le\norm{\mathbf 1_{M_Y}\psi(t)}+\epsilon_K$, with $\epsilon_K=C^{K+1}K!(\hbar t/\alpha^2)^{K+1}\pcNorm a\to0$ on the window for
fixed $K$. This proves (ii).

\emph{Step 7: the classical transport.} For $q$ continuous with compact support in $M$, the closed horocycles equidistribute:
$\int_0^1q(x+iy)\,dx\to\tfrac3\pi\int_Mq\,d\nu$ as $y\to0$. For smooth $q$ with compact support this is in \cite[Section~1, pp.~279--280]{Zag81}
(see also \cite{Sar81}); for continuous $q$ it follows by uniform approximation, since the measures $dx$ on the horocycles and
$\tfrac3\pi d\nu$ have mass $1$. The left side of (iii) equals
$\int_0^\infty|a_0(\eta)|^2\big(\int_0^1q(x+i\ue^{-t}\eta)\,dx\big)\,d\eta/\eta^2$, and (iii) follows by dominated convergence. \qed


\section{Open problems and outlook}\label{sec:outlook}
\begin{enumerate}
\item \emph{Longer times.} The window $t\le\delta(\alpha/\hbar)^{1/2}$ comes from the mollification at scale
$\varepsilon=(\alpha\hbar)^{1/2}/2$: the mollified flow is controlled up to times of order $\varepsilon/\hbar$
(Propositions~\ref{prop:a0} and~\ref{prop:ws}), and the defect of the mollifier is small only when $\varepsilon$ is small
compared with $\alpha$ (Lemma~\ref{lem:unw}). \cite[Conjecture~1]{S} is a dispersive estimate for the unmollified flow on the
time scale $1/\hbar$; according to \cite{S} it would allow one to extend the time scale from $\hbar^{-1/2}$ to $\hbar^{-1}$. We have
not examined which of the other steps of the present proofs would then extend. Whether the periodized approximation stays accurate beyond
$\hbar^{-1/2}$ is open. When $J\ne0$, on a surface of finite area Proposition~\ref{prop:mean} limits the time of validity of each
approximation $u^{(K)}(t)_\Gamma$ of fixed order $K$ in terms of $J$ (Remark~\ref{rem:mean}); the size of $J$ as $\hbar\to0$ is not
determined here.
\item \emph{Expectation values on compact surfaces.} By Corollary~\ref{ev:transfer} the expectation values of
$\cU(t)(a\,\ue^{i\varphi_b/\hbar})_\Gamma$ up to times of order $\hbar^{-1/2}$ are those of the explicit sum $u^{(K)}(t)_\Gamma$. Beyond the Ehrenfest time
the sum is not evaluated here, except in the case of Proposition~\ref{ev:mod}. On a compact surface there is no cusp to
escape into, and neither equidistribution in the sense of \cite{Sch05} nor the random wave behaviour of \cite{Ber77,BB79,IR22} is
known for these states between the Ehrenfest time and $\hbar^{-1/2}$.
\item \emph{Expectation values on the modular surface.} Proposition~\ref{ev:mod} concerns amplitudes whose periodization along
the cusp does not depend on $x$, and it gives a sufficient time only. We conjecture, on the heuristic of Remark~\ref{ev:time},
that the state leaves after about $2\log(1/\hbar)$ and that the classical prediction holds before; neither is proved here. For general amplitudes of the class the part with
nonzero Fourier modes in $x$ is not treated. The states of Brooks \cite{Bro22} are not in the class
(Section~\ref{subsec:context}); a transfer to them by approximation in $L^2$ would need a bound for the periodized
approximation of the difference, which Proposition~\ref{os:p1} gives only with the factor $\ue^{t/2}(1+t/\log2)^{c'}$ and a
weighted norm.
\item \emph{The cusp factor.} Whether the factor $m(z)^{1/2}$ in Corollary~\ref{pc:cor16} is needed for irrational $b$,
and on cusped surfaces other than the modular surface. The proof of
Theorem~\ref{cl:main} would have to be redone with the cusp and $b$ at different
points (Remark~\ref{cl:scope}); a bound depending on $b$ may be smaller there.
\item \emph{Compact surfaces.} Proposition~\ref{prop:g5} shows that the order $\hbar t$ of the leading error is attained on
$\Hb$, that is, for the trivial group. Whether it is attained on compact surfaces is not decided here.
\item \emph{Other geometries.} The exact factorization \eqref{eq:factor} uses constant curvature and a phase attached to a
boundary point. On real hyperbolic space of dimension $n$, in the upper half-space model with $Y=y\partial_y$, the
analogous factorization holds, with $S(t)a(x,y)=\ue^{-(n-1)t/2}a(x,\ue^ty)$ and
$\Delta(t)=\ue^{-2t}y^2\Delta_x+Y^2-(n-1)Y$; an extension to quotients of dimension $n$ would need the analytic, weighted and
periodization estimates of this paper in that dimension. In variable negative curvature a factorization of the same form
is described in \cite[Remark~1.11]{S}, for Busemann functions $\varphi_b$ and the transport along their gradient flow; what
is lost there is the explicit operator $\Delta(t)=\ue^{-2t}X^2+Y^2-Y$, on which all estimates of this paper rest.
\end{enumerate}

\section*{AI-use disclosure}

AI systems were used to develop arguments, write the manuscript, formalize proofs and check the work.
Anthropic's Claude Opus~5.5 (\texttt{claude-opus-5-5}), running in Claude Code, developed the analytic estimates in
Appendices~\ref{app:kl}--\ref{app:ws}, the critical-weight multiplier theorem (Theorem~\ref{os:cw}), and their Lean formalizations.
Claude Fable~5.2 (\texttt{claude-fable-5-2}) developed the transfer proof of Theorem~\ref{pc:main}, the extension to arbitrary
radii (Corollary~\ref{pc:allalpha}), the mean formula (Proposition~\ref{prop:mean}), and Section~\ref{sec:expect}; it wrote the
corresponding formalizations and revised the manuscript to follow the formal proofs. Claude Haiku~5.5
(\texttt{claude-haiku-5-5}) checked citations, transcriptions and numerical consistency; Fable checked and applied its corrections.

OpenAI GPT-6.1 Sol (\texttt{gpt-6.1-sol}) supplied the counterexample on $\Hb$ (Proposition~\ref{prop:g5}),
the Mellin product computation in Appendix~\ref{sec:dd}, the Mehler comparison in Lemma~\ref{lem:gutzmer}, and contributions to
the tail bounds and truncation in Section~\ref{subsec:cl}; Fable completed the cusp proofs and the treatment of the example amplitudes. Sol, Google Gemini~4 Argon (\texttt{gemini-4-argon}),
OpenAI GPT-6~Pro (\texttt{gpt-6-pro}), and Claude Fable~5.2 performed AI reviews of working notes or manuscript versions.
Gemini also contributed to an earlier proof route that was subsequently replaced.
OpenAI Codex and ChatGPT sessions checked revisions and compared manuscript statements with their formal counterparts;
Codex also carried out the final presentation review and editorial reorganization.
These were AI checks and do not constitute external peer review. The chronological contribution record is retained in earlier
repository revisions; the account preceding this editorial revision is at commit \texttt{a223ee5}.

The main approximation results and the supporting results specified in Appendix~\ref{app:lean} have been formalized in Lean~4;
that appendix states the precise scope and the exclusions. Numerical computations served only as consistency checks, never as
proof inputs. The author directed the research, selected and checked the final arguments, and takes responsibility for all
mathematical claims, citations, computations and wording.

\appendix

% Appendix tags contain a section letter, so they use lining characters;
% \oldstylenums is retained for the numeric equation tags in the main text.
\makeatletter
\def\tagform@#1{%
  \maketag@@@{(\ignorespaces#1\unskip\@@italiccorr)}}
\renewcommand{\eqref}[1]{\textup{(\ref{#1})}}
\makeatother



\section[The analyticity weight and the mollified flow]{The analyticity weight and the mollified flow\\ (Kontorovich--Lebedev)}\label{app:kl}
Let $J_\varepsilon=\chi_\varepsilon(\varepsilon\Lambda)$, $\chi_\varepsilon=g_\varepsilon*\chi$, where, as in
\cite[(4.20)--(4.21)]{S}, $\chi\in C_c^\infty(\mathbb R)$ has $\operatorname{supp}\chi\subset[-2,2]$ and $\chi=1$ on $[-1,1]$, and in
addition $\chi$ is real, even and $0\le\chi\le1$; here $g_\varepsilon=\varepsilon^{-1}g(\cdot/\varepsilon)$ with $g$ the unit
Gaussian, so $J_\varepsilon=j_\varepsilon(\Lambda)$ with $j_\varepsilon=g*\chi(\varepsilon\,\cdot)$ and $0\le j_\varepsilon\le1$. Put $\Delta_\varepsilon(t)=J_\varepsilon\Delta(t)J_\varepsilon$, with $\Delta(t)=\ue^{-2t}X^2+Y^2-Y$.

\begin{lemma}[elliptic estimate]\label{lem:ell}
For $f,g\in C_c^\infty(\Hb)$,
$\langle Xf,g\rangle=-\langle f,Xg\rangle$ and $\langle Yf,g\rangle=-\langle f,Yg\rangle+\langle f,g\rangle$, and
\[
\norm{X^2f}^2+\norm{(Y^2-Y)f}^2+2\norm{XYf}^2=\norm{\Delta f}^2+2\norm{Xf}^2\le\norm{\Delta f}^2+2\norm f\norm{\Delta f},
\]
and $\norm{Xf}^2+\norm{Yf}^2=\langle f,-\Delta f\rangle$. By closure, these three relations hold for every
$f\in\operatorname{Dom}(\Delta)$, and then $\norm{X^2f}$, $\norm{(Y^2-Y)f}\le\norm{\Delta f}+\norm f$,
$\norm{Y^2f}\le2(\norm{\Delta f}+\norm f)$, and
$\norm{\Delta(t)f}\le2(\norm{\Delta f}+\norm f)$ for $t\ge0$. Moreover
$\Delta_\varepsilon(t)$ is a bounded self-adjoint operator on $L^2(\Hb)$, norm continuous in $t\ge0$.
\end{lemma}
\begin{proof}
Let $f,g\in C_c^\infty(\Hb)$. Integration by parts with $d\nu=y^{-2}dx\,dy$ gives the two formulas for $X=y\partial_x$ and $Y=y\partial_y$;
that is, $X^*=-X$ and $Y^*=1-Y$ on test functions. So $-\Delta=X^*X+Y^*Y$, which gives
$\norm{Xf}^2+\norm{Yf}^2=\langle f,-\Delta f\rangle$, and
$Y^2-Y=-Y^*Y$. Since $[Y,X]=X$, $YX^2=X^2Y+2X^2$, and therefore
$\langle X^2f,Y^*Yf\rangle=\langle YX^2f,Yf\rangle=-\norm{XYf}^2-2\langle Xf,XYf\rangle$. Moreover $XY=YX-X$ and $Y+Y^*=1$ give
$\mathrm{Re}\langle Xf,XYf\rangle=\mathrm{Re}\langle Xf,YXf\rangle-\norm{Xf}^2=-\tfrac12\norm{Xf}^2$. With $\Delta=X^2-Y^*Y$,
\begin{align*}
\norm{\Delta f}^2&=\norm{X^2f}^2+\norm{Y^*Yf}^2+2\norm{XYf}^2+4\,\mathrm{Re}\langle Xf,XYf\rangle\\
&=\norm{X^2f}^2+\norm{(Y^2-Y)f}^2+2\norm{XYf}^2-2\norm{Xf}^2,
\end{align*}
which is the identity; the inequality holds because $\norm{Xf}^2\le\langle f,-\Delta f\rangle\le\norm f\norm{\Delta f}$. Now let
$f\in\operatorname{Dom}(\Delta)$. Since $\Delta$ is essentially self-adjoint on $C_c^\infty(\Hb)$, there are test functions $f_n\to f$ with
$\Delta f_n\to\Delta f$; by the inequality applied to $f_n-f_m$, the sequences $Xf_n$, $Yf_n$, $X^2f_n$, $(Y^2-Y)f_n$, $XYf_n$ converge in
$L^2$, to the distributional derivatives of $f$, and the identities and the inequality pass to the limit. The right side is at most
$(\norm{\Delta f}+\norm f)^2$, which gives the bounds for $X^2f$ and $(Y^2-Y)f$. Since $\norm{Yf}^2\le\norm f\norm{\Delta f}$,
$\norm{Yf}\le\tfrac12(\norm{\Delta f}+\norm f)$, which gives the bound for $Y^2f=(Y^2-Y)f+Yf$; and $\Delta(t)=\ue^{-2t}X^2+(Y^2-Y)$.
Finally $J_\varepsilon$ maps $L^2$ into $\operatorname{Dom}(\Delta)$ with
$\norm{\Delta J_\varepsilon}\le\sup_\lambda\lambda^2j_\varepsilon(\lambda)<\infty$, the operators $X^2$ and $Y^2-Y=-Y^*Y$ are symmetric on
$\operatorname{Dom}(\Delta)$, and $\Delta_\varepsilon(t)-\Delta_\varepsilon(t')=(\ue^{-2t}-\ue^{-2t'})J_\varepsilon X^2J_\varepsilon$.
\end{proof}

By Lemma~\ref{lem:ell} the equation $i\partial_sv=-\tfrac\hbar2\Delta_\varepsilon(s)v$ has a unitary propagator
$V_\varepsilon(s_1,s_0)$ on $L^2(\Hb)$, $s_0,s_1\ge0$, norm differentiable in both variables. The mollified flow of
\cite[(4.32)--(4.33)]{S} is $V_\varepsilon(s)=V_\varepsilon(s,0)$, and $V_\varepsilon(s_1,s_0)=V_\varepsilon(s_1)V_\varepsilon(s_0)^*$.
All estimates below are stated for $V_\varepsilon(s_1,s_0)$; this contains $V_\varepsilon(t)$ and $V_\varepsilon(t)^*$.

\medskip\noindent\textbf{The Kontorovich--Lebedev picture.}
Fourier transform in $x$; on the mode $\ue^{ikx}$, $k\ne0$, put $z=|k|y$ and write the function of $y$ as $z^{1/2}f(z)$. Then
$L^2(dx\,dy/y^2)$ becomes $L^2(dz/z)$ in $f$, $-\Delta$ becomes $-(z\partial_z)^2+z^2+\tfrac14$, and $X$ becomes multiplication by
$i\operatorname{sgn}(k)\,z$; on the mode $k=0$, $X=0$. The Kontorovich--Lebedev transform $F(\nu)=\int_0^\infty K_{i\nu}(z)f(z)\,dz/z$, with Parseval
$\int|f|^2dz/z=\int|F|^2w(\nu)\,d\nu$, $w(\nu)=\tfrac2{\pi^2}\nu\sinh(\pi\nu)$ (\cite[eqs.~(1.1)--(1.3)]{Y}, applied to $f(z)/z$), diagonalizes
$-\Delta=\nu^2+\tfrac14$. So $\Lambda$ acts as multiplication by $\Lambda(\nu)=(\nu^2+\tfrac14)^{1/2}$, and
$|\Lambda(\nu)-\Lambda(\nu')|\le|\nu-\nu'|$. In the unitary picture $G=w^{1/2}F$, multiplication by $z$ has the kernel
$Zt_1(\nu,\nu')$ computed in step (X) of the proof of Lemma~\ref{lem:fc}, and
\begin{equation}\label{eq:zt1}
0\le Zt_1(\nu,\nu')\le\frac{(\nu\nu')^{1/2}}{2\cosh(\pi(\nu-\nu')/2)} .
\end{equation}
The operators are the same for every $k$, so all bounds are uniform in $k$.

\begin{lemma}[clipped conjugation of $XJ_\varepsilon$]\label{lem:clip}
Let $\tau_0<\pi/2$, $L\ge\tfrac12$ and $\ell:=\min(\Lambda,L)$. For real $a$ let $R_a:=\ue^{a\ell}XJ_\varepsilon\ue^{-a\ell}$, a bounded
operator. Put
\begin{gather*}
g_K(x):=\frac{(1+|x|)\,\ue^{\tau_0|x|}}{2\cosh(\pi x/2)},\qquad g_Q(x):=\frac{(1+3|x|)\,|x|\,\ue^{\tau_0|x|}}{2\cosh(\pi x/2)},\\
A_1:=\int_{\mathbb R}g_K(x)\,dx,\qquad A_2:=\int_{\mathbb R}g_Q(x)\,dx .
\end{gather*}
For $0\le\tau\le\tau_0$, $\varepsilon>0$ and $u\in L^2(\Hb)$,
\[
\bigl\|\tfrac12(R_\tau+R_{-\tau})u\bigr\|\le A_1\norm{\Lambda J_\varepsilon u},\qquad
\bigl\|\tfrac12(R_\tau-R_{-\tau})u\bigr\|\le A_2\,\tau\,\norm{\ell J_\varepsilon u},
\]
and consequently
\[
\bigl|\mathrm{Im}\langle R_{-\tau}u,R_\tau u\rangle\bigr|\le2A_1A_2\,\tau\,\norm{\Lambda J_\varepsilon u}\,\norm{\ell J_\varepsilon u}.
\]
\end{lemma}
\begin{proof}
In the Kontorovich--Lebedev picture $\ell$ and $J_\varepsilon$ are multiplications by $\ell(\nu)=\min(\Lambda(\nu),L)$ and
$j_\varepsilon(\Lambda(\nu))$. Put $x=\nu-\nu'$ and $\delta=\ell(\nu)-\ell(\nu')$, so $|\delta|\le|x|$. Up to the factor
$i\operatorname{sgn}(k)$ of modulus one, the operators $\tfrac12(R_\tau+R_{-\tau})$ and $\tfrac12(R_\tau-R_{-\tau})$ have the kernels
\[
Zt_1(\nu,\nu')\cosh(\tau\delta)\,j_\varepsilon(\Lambda(\nu'))\qquad\text{and}\qquad Zt_1(\nu,\nu')\sinh(\tau\delta)\,j_\varepsilon(\Lambda(\nu')).
\]

\emph{Even part.} Since $(\nu\nu')^{1/2}\le\nu'+|x|/2$ and $\Lambda(\nu')\ge\max(\nu',\tfrac12)$, we have
$(\nu\nu')^{1/2}\le(1+|x|)\Lambda(\nu')$. By \eqref{eq:zt1} and $\cosh(\tau\delta)\le\ue^{\tau_0|x|}$, the even part is
$K_A\circ\Lambda J_\varepsilon$ with a kernel $|K_A(\nu,\nu')|\le g_K(\nu-\nu')$, and the Schur test gives $\norm{K_A}\le A_1$.

\emph{Odd part.} Here $|\sinh(\tau\delta)|\le\tau|\delta|\,\ue^{\tau_0|x|}$, and $\delta=0$ when $\Lambda(\nu)\ge L$ and
$\Lambda(\nu')\ge L$. In the remaining cases $(\nu\nu')^{1/2}\le(1+|x|)\,\ell(\nu')$: if $\Lambda(\nu')\le L$ this is the inequality
of the even part; if $\Lambda(\nu')>L\ge\Lambda(\nu)$, then $\ell(\nu')=L$ and
$(\nu\nu')^{1/2}\le\nu+|x|/2\le L+|x|L$, because $\nu\le\Lambda(\nu)\le L$ and $L\ge\tfrac12$. So the odd part is $K_B\circ\ell J_\varepsilon$ with
$|K_B(\nu,\nu')|\le\tau\,|x|\,g_K(x)\le\tau\,g_Q(x)$, and the Schur test gives $\norm{K_B}\le A_2\tau$.

\emph{The imaginary part.} With $P=\tfrac12(R_\tau+R_{-\tau})u$ and $Q=\tfrac12(R_\tau-R_{-\tau})u$,
$\langle R_{-\tau}u,R_\tau u\rangle=\langle P-Q,P+Q\rangle=\norm P^2-\norm Q^2+2i\,\mathrm{Im}\langle P,Q\rangle$, so
$|\mathrm{Im}\langle R_{-\tau}u,R_\tau u\rangle|\le2\norm P\norm Q$.
\end{proof}

\begin{proposition}[A0]\label{prop:a0}
Let $\tau_0<\pi/2$ and $c_A:=(1+32A_1A_2)^{-1}$, with $A_1,A_2$ as in Lemma~\ref{lem:clip}. Let $0<\varepsilon\le1$, $\hbar>0$,
$0\le\alpha'\le\alpha\le\tau_0$ with $\alpha>0$, and $s_0,s_1\ge0$ with
\[
|s_1-s_0|\le c_A\,\frac{\alpha-\alpha'}{\alpha}\,\frac{\varepsilon}{\hbar}.
\]
If $a\in\operatorname{Dom}(\ue^{\alpha\Lambda})$, then $V_\varepsilon(s_1,s_0)a\in\operatorname{Dom}(\ue^{\alpha'\Lambda})$ and
\[
\norm{\ue^{\alpha'\Lambda}V_\varepsilon(s_1,s_0)a}\le\norm{\ue^{\alpha\Lambda}a}.
\]
\end{proposition}
\begin{proof}
\emph{A frequency bound.} For $L\ge4/\varepsilon$ and $\lambda\ge0$, $\lambda^2j_\varepsilon(\lambda)^2\le(32/\varepsilon)\min(\lambda,L)$.
Indeed $j_\varepsilon(\lambda)\le\int_{|\mu|\le2/\varepsilon}g(\lambda-\mu)\,d\mu\le1$. If $\lambda\le4/\varepsilon$ the left side is at most
$(4/\varepsilon)\lambda$. If $\lambda>4/\varepsilon$, then $\lambda-\mu>\lambda/2$ on the range of integration, so
$j_\varepsilon(\lambda)\le\ue^{-\lambda^2/8}$ and $\lambda^2j_\varepsilon(\lambda)^2\le4/\ue$, while $\min(\lambda,L)\ge4/\varepsilon\ge4$.

\emph{The clipped energy.} Fix $L\ge4/\varepsilon$, put $\ell=\min(\Lambda,L)$ and $\rho:=32A_1A_2\hbar/\varepsilon$, and let
$\tau(s):=\alpha\,\ue^{-\rho|s-s_0|}$, $v(s):=V_\varepsilon(s,s_0)a$, $u(s):=\ue^{\tau(s)\ell}v(s)$. All operators involved are bounded.
Since $v'=\tfrac{i\hbar}2\Delta_\varepsilon(s)v$,
\[
\tfrac{d}{ds}\norm{u}^2=2\tau'\langle u,\ell u\rangle-\hbar\,\mathrm{Im}\bigl\langle u,\ue^{\tau\ell}\Delta_\varepsilon(s)\ue^{-\tau\ell}u\bigr\rangle .
\]
Now $\Delta_\varepsilon(s)=J_\varepsilon\Delta J_\varepsilon-(1-\ue^{-2s})J_\varepsilon X^2J_\varepsilon$. The first term commutes with $\ell$ and is
self-adjoint, so it does not contribute. For the second, $\ue^{\tau\ell}J_\varepsilon X^2J_\varepsilon\ue^{-\tau\ell}=-R_{-\tau}^*R_\tau$,
because $X^*=-X$. Hence, by Lemma~\ref{lem:clip}, $\norm{\ell J_\varepsilon u}\le\norm{\Lambda J_\varepsilon u}$ and the frequency bound,
for $s\ge0$,
\[
\Bigl|\tfrac{d}{ds}\norm u^2-2\tau'\langle u,\ell u\rangle\Bigr|
=\hbar(1-\ue^{-2s})\bigl|\mathrm{Im}\langle R_{-\tau}u,R_\tau u\rangle\bigr|
\le2A_1A_2\hbar\,\tau\,\langle u,\Lambda^2J_\varepsilon^2u\rangle\le2\rho\,\tau\,\langle u,\ell u\rangle .
\]
Since $\tau'=-\rho\tau$ for $s>s_0$ and $\tau'=\rho\tau$ for $s<s_0$, $\norm{u(s)}$ does not increase as $s$ moves away from $s_0$, in
either direction; all times between $s_0$ and $s_1$ are nonnegative. Therefore
$\norm{\ue^{\tau(s_1)\ell}v(s_1)}\le\norm{\ue^{\alpha\ell}a}\le\norm{\ue^{\alpha\Lambda}a}$.

\emph{The radius.} By the hypothesis on $|s_1-s_0|$ and $32A_1A_2c_A\le1$,
$\tau(s_1)\ge\alpha(1-\rho|s_1-s_0|)\ge\alpha'$. So $\norm{\ue^{\alpha'\min(\Lambda,L)}v(s_1)}\le\norm{\ue^{\alpha\Lambda}a}$ for every
$L\ge4/\varepsilon$, and monotone convergence in the spectral representation of $\Lambda$, as $L\to\infty$, gives the claim.
\end{proof}

\section{Multiplication by the weight and analyticity}\label{app:m}
Here $X_0=\cosh d(x_0,\cdot)$ and $\tilde\psi=\log X_0$. The theorems use the two lemmas of the first subsection, in which the radius
is reduced by a fixed factor. A fixed factor changes only constants, since $\alpha$ enters the conclusions through
$(\alpha\hbar)^{-2N}(\hbar t/\alpha^2)^{K+1}$, the condition $\alpha^3\ge\kappa\hbar$ and the time window. The second subsection contains
a comparison by the Gutzmer identity \cite[\S7.2]{K}, which loses no fixed factor of the radius. It is not used by the theorems; it is
used only for Lemma~\ref{lem:ex}, on the example amplitudes of \cite{S}.

\subsection{The forms used by the theorems}
\begin{lemma}[analytic vectors from word bounds]\label{lem:expw}
Let $0<\alpha\le1$, $M\ge1$ and $B,N_0\ge0$. Let $G\in L^2(\Hb)$ be such that every word $Z_1\cdots Z_m$, $Z_i\in\{X,Y\}$, $m\ge0$,
applied to $G$ in the sense of distributions lies in $L^2$, with
\[
\norm{Z_1\cdots Z_mG}\le B\,(m+1)\,m!\,(M/\alpha)^m\,N_0 .
\]
Then $G\in\operatorname{Dom}(\ue^{b\Lambda})$ and $\norm{\ue^{b\Lambda}G}\le C_*BN_0$ for $b=\alpha/(4M)$, with an absolute constant $C_*$.
\end{lemma}
\begin{proof}
$\Delta$ is essentially self-adjoint on $C_c^\infty(\Hb)$, so $G\in\operatorname{Dom}(\Delta^n)$ for every $n$, and
$(-\Delta)^n=(-X^2-Y^2+Y)^n$ is a sum of $3^n$ words of length at most $2n$ with coefficients $\pm1$. Since
$m\mapsto(m+1)m!(M/\alpha)^m$ is increasing,
$\norm{\Lambda^{2n}G}=\norm{(-\Delta)^nG}\le3^nB(2n+1)(2n)!(M/\alpha)^{2n}N_0$, and so
\[
\sum_{n\ge0}\frac{b^{2n}}{(2n)!}\norm{\Lambda^{2n}G}\le BN_0\sum_{n\ge0}(2n+1)\Bigl(\frac{3}{16}\Bigr)^n\le2BN_0 .
\]
The partial sums of $\sum_n(b\lambda)^{2n}/(2n)!$ increase to $\cosh(b\lambda)$, so monotone convergence in the spectral representation
of $\Lambda$ gives $\norm{\cosh(b\Lambda)G}\le2BN_0$, and $\ue^{b\lambda}\le2\cosh(b\lambda)$. So $C_*=4$ works.
\end{proof}

\begin{lemma}[M, reduced radius]\label{lem:mw}
Let $\beta\ge0$ and $\tau_0<\pi/2$, let $B_\beta$, $C_3$ be the constants of Lemma~\ref{lem:sc}, and put $c_M:=1/(4\max(C_3,1))$. For
$0<\alpha\le\min(1,\tau_0)$ and every $a$ with $\norm{\ue^{\alpha\Lambda}X_0^\beta a}<\infty$,
\[
\norm{\ue^{c_M\alpha\Lambda}a}\le C_*B_\beta\,\norm{\ue^{\alpha\Lambda}X_0^\beta a},
\]
with $C_*$ the constant of Lemma~\ref{lem:expw}.
\end{lemma}
\begin{proof}
By Lemma~\ref{lem:sc} and $X_0\ge1$, every word satisfies
$\norm{Z_1\cdots Z_ma}\le\norm{X_0^\beta Z_1\cdots Z_ma}\le B_\beta(m+1)m!(C_3/\alpha)^m\norm{\ue^{\alpha\Lambda}X_0^\beta a}$.
Apply Lemma~\ref{lem:expw} with $M=\max(C_3,1)$.
\end{proof}

\begin{lemma}[\texorpdfstring{M2, reduced radius: the weights $\langle d\rangle$ and $\tilde\psi$}{M2, reduced radius}]\label{lem:m2w}
For $\beta\ge0$ there are $c_0\in(0,1]$ and $C$, depending only on $\beta$, such that for every base point $x_0$, every
$0<\alpha\le1$ and every $a$ with $\norm{\ue^{\alpha\Lambda}\ue^{\beta\langle d\rangle}a}<\infty$,
\[
\norm{\ue^{c_0\alpha\Lambda}X_0^\beta a}\le C\,\norm{\ue^{\alpha\Lambda}\ue^{\beta\langle d\rangle}a}.
\]
\end{lemma}
\begin{proof}
The fields $X,Y$ and the operator $\Lambda$ commute with the affine isometries $z\mapsto pz+q$, $p>0$, which act transitively on
$\Hb$; so we may take $x_0=i$. Put $F:=\ue^{\beta\langle d\rangle}a\in\operatorname{Dom}(\ue^{\alpha\Lambda})$. Then $X_0^\beta a=mF$ with
\[
m:=X_0^\beta\ue^{-\beta\langle d\rangle}=\Phi(\tilde\psi),\qquad\Phi(\psi):=\exp\beta\bigl(\psi-(1+d^2)^{1/2}+1\bigr),
\]
where $d\ge0$ is the function of $\psi\ge0$ given by $\psi=\log\cosh d$.

\emph{Step 1: $|\Phi^{(k)}|\le A\,k!\,\rho^{-k}$ on $[0,\infty)$, with $A,\rho>0$ depending only on $\beta$.} The function $\Phi$
extends holomorphically to a neighbourhood of every point of $[0,\infty)$. Near $0$ this holds because $d^2$ is the analytic inverse at
$0$ of $s\mapsto\log\cosh\sqrt s$, which has derivative $\tfrac12$ there. Near a point of $(0,\infty)$ it holds because
$d=\operatorname{arcosh}(\ue^\psi)$ is analytic there. On the strip $\{\operatorname{Re}\psi>\tfrac{17}{10},\ |\operatorname{Im}\psi|<\tfrac14\}$,
\begin{gather*}
\Phi=\exp\beta\bigl(1-\ell-(D+R)^{-1}\bigr),\\
\ell=\log\bigl(1+(1-\ue^{-2\psi})^{1/2}\bigr),\qquad D=\psi+\ell,\qquad R=(1+D^2)^{1/2},
\end{gather*}
with principal branches: on the real axis $D=d$ and $\psi-R+1=1-\ell-(D+R)^{-1}$, because $R^2-D^2=1$. On the strip
$\operatorname{Re}(1-\ue^{-2\psi})>0$, $\operatorname{Re}\ell\ge0$, $\operatorname{Re}(1+D^2)>0$ and $\operatorname{Re}(D+R)>\tfrac{17}{10}$, so $\Phi$
is holomorphic there with $|\Phi|\le\ue^{2\beta}$. Compactness of $[0,2]$, and Cauchy's estimate on discs of radius $\tfrac15$ around
the points of $[2,\infty)$, give Step~1.

\emph{Step 2: $|Z_Sm|\le A\,c^j\,j!$ for every word $Z_S$ of length $j$, with $c=4+3/\rho$.} Put $\xi:=x/X_0$ and $\upsilon:=y/X_0$, so
$|\xi|\le1$, $0<\upsilon\le2$, $X\tilde\psi=\xi$ and $Y\tilde\psi=\upsilon-1$. The polynomials in $(\xi,\upsilon)$ are stable under $X$ and $Y$:
$X\xi=\upsilon-\xi^2$, $X\upsilon=-\xi\upsilon$, $Y\xi=\xi-\xi\upsilon$, $Y\upsilon=2\upsilon-\upsilon^2$. For
$Q=\sum c_{ab}\xi^a\upsilon^b$ let $\norm Q_*:=\sum|c_{ab}|2^b$, so $|Q|\le\norm Q_*$ pointwise. The four formulas give
$\norm{ZQ}_*\le4(\deg Q)\norm Q_*$, and $\norm{(Z\tilde\psi)Q}_*\le3\norm Q_*$. By induction on $j$,
$Z_S\Phi(\tilde\psi)=\sum_{k=0}^j\Phi^{(k)}(\tilde\psi)\,Q_{S,k}(\xi,\upsilon)$ with polynomials of degree at most $j$ satisfying
$Q_{ZS,k}=(Z\tilde\psi)Q_{S,k-1}+ZQ_{S,k}$. With $q_{j,k}:=\max_{|S|=j}\norm{Q_{S,k}}_*$ this gives
$q_{j+1,k}\le3q_{j,k-1}+4j\,q_{j,k}$, and then $T_j:=\sum_kk!\,\rho^{-k}q_{j,k}$ satisfies
$T_{j+1}\le(3(j+1)/\rho+4j)T_j\le c\,(j+1)\,T_j$, so $T_j\le c^jj!$. By Step~1, $|Z_Sm|\le A\,T_j$.

\emph{Step 3.} By the Leibniz rule for words, Step~2 and Lemma~\ref{lem:wc} at $\alpha''=0$ (with $\tau_0=1$), for a word of length $j$,
\[
\norm{Z_S(mF)}\le\sum_{S'\subseteq S}A\,c^{|S'|}|S'|!\,\Bigl(\frac{\ue C_1}{\alpha}\Bigr)^{j-|S'|}(j-|S'|)!\,\norm{\ue^{\alpha\Lambda}F}
\le A\,(j+1)\,j!\,(M/\alpha)^j\norm{\ue^{\alpha\Lambda}F},
\]
with $M:=\max(c,\ue C_1,1)$. Lemma~\ref{lem:expw} gives the claim with $c_0=1/(4M)$ and $C=C_*A$.
\end{proof}

\subsection{A comparison by the Gutzmer identity}\label{subsec:gutzmer}
This subsection is used only for Lemma~\ref{lem:ex}. Let $G=SL(2,\mathbb R)$, $h_0=\mathrm{diag}(1,-1)$ and $x_0=i$, so that
$d(x_0,\exp(th_0)x_0)=2|t|$. For $0<R\le\pi/4$ let $\Xi_R=G\exp(i(-R,R)h_0)x_0$, a $G$-invariant domain in the complexification
$G_{\mathbb C}/K_{\mathbb C}$ of $\Hb=G/K$; $\Xi_{\pi/4}$ is the crown \cite[\S7.2]{K}. For a holomorphic function $f$ on $\Xi_R$ and $0\le r<R$ put
\[
N_r(f)^2:=\int_G|f(g\exp(irh_0)x_0)|^2\,dg ,
\]
with the Haar measure normalized so that $N_0(f)=\norm f_{L^2(\Hb)}$.

\begin{lemma}[comparison with the orbit norm]\label{lem:gutzmer}
Let $0<r<R\le\pi/4$ with $r\le\tfrac12$.
\begin{itemize}
\item[(a)] Let $f\in L^2(\Hb)$ extend holomorphically to $\Xi_R$ with $\sup_{0\le\rho<R}N_\rho(f)<\infty$. Then
$f\in\operatorname{Dom}(\ue^{r\Lambda})$ and $\norm{\ue^{r\Lambda}f}\le(\pi\ue)^{1/2}N_r(f)$.
\item[(b)] The function $X_0=\cosh d(x_0,\cdot)$ extends holomorphically to $\Xi_{\pi/4}$, and for $g\in G$ its value at
$g\exp(irh_0)x_0$ is $\cos(2r)\,A+i\sin(2r)\,B$, where $A=X_0(gx_0)\ge1$ and $B$ is real with $B^2\le A^2-1$.
\end{itemize}
\end{lemma}
\begin{proof}
(b) Realize $\Hb$ as the sheet $p_0>0$ of the hyperboloid $\{p\in\mathbb R^3:-p_0^2+p_1^2+p_2^2=-1\}$, with $x_0=e_0=(1,0,0)$, the group $G$
acting through $SO_0(2,1)$, and $G_{\mathbb C}/K_{\mathbb C}$ as the complex quadric given by the same equation. Then $X_0(p)=p_0$ is a linear function,
$\exp(th_0)e_0=(\cosh2t,\sinh2t,0)$, and so $\exp(irh_0)e_0=\cos(2r)\,e_0+i\sin(2r)\,e_1$. Hence
$X_0(g\exp(irh_0)x_0)=\cos(2r)\,(ge_0)_0+i\sin(2r)\,(ge_1)_0$. Here $(ge_0)_0=X_0(gx_0)=A$, and $q:=ge_1$ is a real vector with
$-q_0^2+q_1^2+q_2^2=1$ that is orthogonal to $p:=ge_0$ for the form: $Aq_0=p_1q_1+p_2q_2$. By the Cauchy--Schwarz inequality
$A^2q_0^2\le(A^2-1)(1+q_0^2)$, that is $q_0^2\le A^2-1$.

(a) In the spectral decomposition of $L^2(\Hb)$ by the Helgason--Fourier transform, $-\Delta$ acts as $\lambda^2+\tfrac14$, $\lambda\ge0$, and
$\Lambda$ as $(\lambda^2+\tfrac14)^{1/2}\le\lambda+\tfrac12$. The Gutzmer identity (\cite[\S7.2, the identity before Theorem~7.2]{K}, where the
spectral parameter is $2\lambda$) states that for $f$ as in (a) and $r<R$ the number $N_r(f)^2$ is the integral of $|\hat f|^2$ against the
Plancherel measure with the additional weight $\phi_\lambda(\exp(2irh_0))$, the value at $\exp(2irh_0)x_0$ of the holomorphic extension of the
spherical function $\phi_\lambda$. A radial eigenfunction of $\Delta$ with eigenvalue $-(\lambda^2+\tfrac14)$ that equals $1$ at $x_0$ is the regular
solution of Legendre's equation in the variable $X_0$, so $\phi_\lambda=P_{-1/2+i\lambda}(X_0)$, a conical function \cite[\S14.20]{DLMF}; by (b),
$X_0(\exp(2irh_0)x_0)=\cos4r$. With $\theta:=4r\in(0,2]$, the Mehler--Dirichlet integral \cite[14.20.9]{DLMF} and the substitution $\phi=\theta v$ give
\[
P_{-1/2+i\lambda}(\cos\theta)=\frac2\pi\int_0^1\frac{\cosh(\lambda\theta v)\,\theta}{\sqrt{2(\cos\theta v-\cos\theta)}}\,dv
\ge\frac2\pi\int_0^1\frac{\cosh(\lambda\theta v)}{\sqrt{1-v^2}}\,dv=I_0(\lambda\theta),
\]
because $2(\cos\theta v-\cos\theta)=4\sin\tfrac{\theta(1+v)}2\sin\tfrac{\theta(1-v)}2\le\theta^2(1-v^2)$; the last equality is
\cite[10.32.1]{DLMF}. From the same formula, $I_0(x)=\tfrac1\pi\int_0^\pi\ue^{x\cos t}dt\ge\tfrac1\pi\int_0^{(1+x)^{-1/2}}\ue^{x(1-t^2/2)}dt
\ge(\pi\sqrt{\ue})^{-1}(1+x)^{-1/2}\ue^{x}$ for $x\ge0$. Since $\ue^{4r\lambda}\ge(1+4r\lambda)^{1/2}\ue^{2r\lambda}$ and
$\ue^{2r\Lambda}\le\ue^{r}\ue^{2r\lambda}\le\sqrt{\ue}\,\ue^{2r\lambda}$,
\[
\phi_\lambda(\exp(2irh_0))\ge(\pi\sqrt{\ue})^{-1}\ue^{2r\lambda}\ge(\pi\ue)^{-1}\ue^{2r(\lambda^2+1/4)^{1/2}} ,
\]
and integrating against $|\hat f|^2$ gives (a).
\end{proof}

\section{Weighted bounds for the mollifier and the mollified flow}\label{app:ws}
In this appendix the base point is $x_0=i$ (Lemma~\ref{lem:base} reduces the general case to it), $X_0=\cosh d(i,\cdot)$, and
$D:=(-\Delta-\tfrac14)^{1/2}$, so that $\Lambda=(D^2+\tfrac14)^{1/2}$. Explicitly $X_0=(x^2+y^2+1)/(2y)$, $XX_0=x$, $YX_0=y-X_0$, and
$(XX_0)^2+(YX_0)^2=X_0^2-1$.

\medskip\noindent\textbf{Bounded weights.} A \emph{bounded weight} is a smooth function $W>0$ on $\Hb$ such that $Z_SW$ is bounded
for every word $Z_S$ in $X,Y$, the empty word included. For $\rho,K\ge0$ let $\mathcal W(\rho,K)$ be the set of bounded weights
with, for $Z\in\{X,Y\}$,
\[
|ZW|\le\rho\,W,\qquad|Z^2W|\le K\,W,\qquad|\Delta W|\le K\,W,\qquad|Z^2(W^2)|\le K\,W^2 .
\]
For $\gamma\ge0$ and an integer $n\ge1$ put $W_{\gamma,n}:=(X_0^{-\gamma}+1/n)^{-1}$. Then $W_{\gamma,n}\le n$,
$W_{\gamma,n}\uparrow X_0^\gamma$ as $n\to\infty$,
\[
(X\log W_{\gamma,n})^2+(Y\log W_{\gamma,n})^2\le\gamma^2,
\]
and $W_{\gamma,n}\in\mathcal W(\gamma,K_\gamma)$ with a constant $K_\gamma$ depending only on $\gamma$ ($K_\gamma=8\gamma^2+7\gamma$ works).
Indeed, with $\phi=X_0^{-\gamma}$ one has $(X\phi)^2+(Y\phi)^2\le\gamma^2\phi^2$ and, since $X^2X_0=y\le2X_0$ and $Y^2X_0=X_0$,
$|Z^2\phi|\le(\gamma^2+3\gamma)\phi$; now differentiate $W_{\gamma,n}=(\phi+1/n)^{-1}$ and $W_{\gamma,n}^2$, and use $\phi\le\phi+1/n$.
All words applied to $\phi$ are bounded (Lemma~\ref{lem:fdb}), hence so are all words applied to $W_{\gamma,n}$.
A bound for all $W_{\gamma,n}$ with a constant independent of $n$ passes to the weight $X_0^\gamma$ by monotone convergence. The
bounds below are used in this way.

\begin{lemma}[wave energy at a bounded weight]\label{lem:wave}
Let $W$ be a bounded weight, $u\in L^2(\Hb)$ and $s\in\mathbb R$.
\begin{itemize}
\item[(a)] If $|XW|\le\rho W$ and $|YW|\le\rho W$, then $\norm{W\cos(s\Lambda)u}\le\ue^{2\rho|s|}\norm{Wu}$. If moreover $u$ lies in the
range of a spectral projection $1_{[0,R]}(\Lambda)$, $R<\infty$, then, with $q=\Lambda^{-1}\sin(s\Lambda)u$,
\[
\norm{WXq}\le\ue^{2\rho|s|}\norm{Wu},\qquad\norm{WYq}\le\ue^{2\rho|s|}\norm{Wu}.
\]
\item[(b)] If $l\ge0$ and $(X\log W)^2+(Y\log W)^2\le l^2$, then $\norm{W\cos(sD)u}\le\ue^{(l+1/2)|s|}\norm{Wu}$.
\end{itemize}
\end{lemma}
\begin{proof}
First let $u$ lie in the range of a spectral projection $1_{[0,L]}(\Lambda)$. Then all functions below are smooth, lie in the domain of
every power of $\Delta$, and are differentiable in $s$. For $f\in\operatorname{Dom}(\Delta)$ and $g$ in the form domain,
$\langle g,\Delta f\rangle=-\langle Xg,Xf\rangle-\langle Yg,Yf\rangle$ (Lemma~\ref{lem:ell}), and $W^2g$ is in the form domain when $g$ is,
since $W$ and $XW$, $YW$ are bounded.

(a) Let $p=\cos(s\Lambda)u$, so $q'=p$ and $p'=\Delta q$, and $E:=\norm{Wp}^2+\norm{WXq}^2+\norm{WYq}^2$. Then
\begin{align*}
E'&=2\,\mathrm{Re}\langle W^2p,\Delta q\rangle+2\,\mathrm{Re}\langle W^2Xq,Xp\rangle+2\,\mathrm{Re}\langle W^2Yq,Yp\rangle\\
&=-2\,\mathrm{Re}\langle X(W^2)p,Xq\rangle-2\,\mathrm{Re}\langle Y(W^2)p,Yq\rangle .
\end{align*}
Since $|X(W^2)|\le2\rho W^2$ and $|Y(W^2)|\le2\rho W^2$, $|E'|\le4\rho\norm{Wp}(\norm{WXq}+\norm{WYq})\le4\rho E$, so
$E(s)\le\ue^{4\rho|s|}E(0)=\ue^{4\rho|s|}\norm{Wu}^2$.

(b) Let $p=\cos(sD)u$ and $q=\int_0^s\cos(rD)u\,dr$, so $q'=p$ and $p'=(\Delta+\tfrac14)q$, and
$E:=\norm{Wp}^2+\norm{WXq}^2+\norm{WYq}^2+\tfrac14\norm{Wq}^2$. The same computation gives
$E'=-2\,\mathrm{Re}\langle X(W^2)p,Xq\rangle-2\,\mathrm{Re}\langle Y(W^2)p,Yq\rangle+\mathrm{Re}\langle W^2p,q\rangle$. Pointwise
$|X(W^2)\,Xq+Y(W^2)\,Yq|\le2l\,W^2(|Xq|^2+|Yq|^2)^{1/2}$, so the first two terms are at most
$4l\norm{Wp}(\norm{WXq}^2+\norm{WYq}^2)^{1/2}\le2lE$ together, and $|\langle W^2p,q\rangle|\le\norm{Wp}\norm{Wq}\le E$. So $|E'|\le(2l+1)E$.

For general $u$ the bounds for $\cos(s\Lambda)u$ and $\cos(sD)u$ follow by applying this to $u_L=1_{[0,L]}(\Lambda)u$ and letting
$L\to\infty$, since $W$ is bounded.
\end{proof}

\begin{lemma}[commutator with the wave group]\label{lem:wcomm}
Let $W\in\mathcal W(\rho,K)$. For $y,g$ in the range of a spectral projection $1_{[0,R]}(\Lambda)$, $R<\infty$, and $s\in\mathbb R$,
\[
\bigl|\langle y,[\cos(s\Lambda),W^2]g\rangle\bigr|\le\bigl(8\rho|s|+(2K+2\rho)s^2\bigr)\ue^{4\rho|s|}\norm{Wy}\norm{Wg}.
\]
\end{lemma}
\begin{proof}
Put $h=W^2$, $c(s)=\cos(s\Lambda)$ and $S(s)=\Lambda^{-1}\sin(s\Lambda)$. The function
$v(s)=[c(s),h]g$ satisfies $v''=\Delta v+[\Delta,h]c(s)g$, $v(0)=v'(0)=0$, so
\[
v(s)=\int_0^sS(s-r)[\Delta,h]c(r)g\,dr.
\]
Here
$[\Delta,h]f=(\Delta h)f+2(Xh)Xf+2(Yh)Yf$, and $[\Delta,h]^*=-[\Delta,h]$. With $q(r):=S(r)y$,
\[
\langle y,v(s)\rangle=-\int_0^s\Bigl(\bigl\langle(\Delta h)\,q(s-r),c(r)g\bigr\rangle+2\bigl\langle(Xh)Xq(s-r)+(Yh)Yq(s-r),c(r)g\bigr\rangle\Bigr)dr,
\]
in which no derivative falls on $c(r)g$. Now
$|Xh|,|Yh|\le2\rho h$ and $|\Delta h|\le|X^2h|+|Y^2h|+|Yh|\le(2K+2\rho)h$. By Lemma~\ref{lem:wave}(a),
$\norm{Wc(r)g}\le\ue^{2\rho|r|}\norm{Wg}$, $\norm{WXq(r)}+\norm{WYq(r)}\le2\ue^{2\rho|r|}\norm{Wy}$, and
$\norm{Wq(r)}\le|r|\ue^{2\rho|r|}\norm{Wy}$ because $q(r)=\int_0^rc(r')y\,dr'$. So the integrand is at most
$\bigl((2K+2\rho)|s-r|+8\rho\bigr)\ue^{2\rho|s|}\norm{Wy}\norm{Wg}$, and integrating gives the claim (with $\ue^{2\rho|s|}$ and
$(K+\rho)s^2$).
\end{proof}

By Fourier inversion, since $j_\varepsilon$ is even,
\[
J_\varepsilon=\frac1\pi\int_0^\infty h_\varepsilon(s)\cos(s\Lambda)\,ds,\qquad h_\varepsilon(s)=\varepsilon^{-1}\hat\chi(s/\varepsilon)\,\ue^{-s^2/2},
\]
where $h_\varepsilon$ is real, even and Schwartz. For $0<\varepsilon\le1$, $k,r\ge0$ and $C\ge0$,
\begin{equation}\label{eq:hmom}
\int_0^\infty|h_\varepsilon^{(r)}(s)|\,\ue^{Cs}\,ds\le C_{r}\,\varepsilon^{-r},\qquad\int_0^\infty|h_\varepsilon(s)|\,s^k\,\ue^{Cs}\,ds\le C_{k}\,\varepsilon^{k},
\end{equation}
with constants depending on $C$ and $\chi$ (substitute $s=\varepsilon\sigma$ and use $\ue^{Cs-s^2/2}\le\ue^{C^2/2}$).

\begin{lemma}[weighted mollifier]\label{lem:wmoll}
Let $W\in\mathcal W(\rho,K)$ and $0<\varepsilon\le1$. There is $M=M(\rho,K,\chi)$ such that for all $u,y,g\in L^2(\Hb)$ and
$Z\in\{X,Y\}$:
\begin{itemize}
\item[(i)] $\norm{WJ_\varepsilon u}\le M\norm{Wu}$;
\item[(ii)] $\norm{WZJ_\varepsilon u}\le M\varepsilon^{-1}\norm{Wu}$;
\item[(iii)] $\norm{W\Delta J_\varepsilon u}\le M\varepsilon^{-2}\norm{Wu}$ and $\norm{WZ^2J_\varepsilon u}\le M\varepsilon^{-2}\norm{Wu}$;
\item[(iv)] $|\langle y,[J_\varepsilon,W^2]g\rangle|\le M\varepsilon\norm{Wy}\norm{Wg}$.
\end{itemize}
In particular $\norm{W\Delta(s)J_\varepsilon u}\le3M\varepsilon^{-2}\norm{Wu}$ for $s\ge0$.
\end{lemma}
\begin{proof}
In (ii) and (iv) it suffices to take $u$, $y$, $g$ in the range of a spectral projection $1_{[0,R]}(\Lambda)$, with $M$ independent of
$R$: the projection commutes with $J_\varepsilon$, the operators $ZJ_\varepsilon$ and $[J_\varepsilon,W^2]$ are bounded on $L^2$, and $W$ is bounded.
(i) follows from Lemma~\ref{lem:wave}(a) and \eqref{eq:hmom}. (ii) With $q(s)=\Lambda^{-1}\sin(s\Lambda)u$, $q'=\cos(s\Lambda)u$ and
$q(0)=0$, an integration by parts gives $ZJ_\varepsilon u=-\tfrac1\pi\int_0^\infty h_\varepsilon'(s)\,Zq(s)\,ds$; use
Lemma~\ref{lem:wave}(a) and \eqref{eq:hmom} with $r=1$. (iii) Since $\Delta\cos(s\Lambda)u=\partial_s^2\cos(s\Lambda)u$ and
$h_\varepsilon'(0)=0$, two integrations by parts give $\Delta J_\varepsilon u=\tfrac1\pi\int_0^\infty h_\varepsilon''(s)\cos(s\Lambda)u\,ds$, and
\eqref{eq:hmom} with $r=2$ gives the first bound. For the second put $f=J_\varepsilon u$. The function $Wf$ lies in
$\operatorname{Dom}(\Delta)$, and
\begin{align*}
Z^2(Wf)&=WZ^2f+2(ZW)Zf+(Z^2W)f,\\
\Delta(Wf)&=W\Delta f+2(XW)Xf+2(YW)Yf+(\Delta W)f,
\end{align*}
with $|ZW|\le\rho W$, $|Z^2W|\le KW$ and $|\Delta W|\le KW$. So Lemma~\ref{lem:ell}, applied to $Wf$, bounds
$\norm{WZ^2f}$ by a constant times $\norm{W\Delta f}+\max_Z\norm{WZf}+\norm{Wf}$, and (i), (ii) and the first bound of (iii) apply. (iv) follows from
Lemma~\ref{lem:wcomm} and \eqref{eq:hmom} with $k=1,2$. The last claim holds because $\Delta(s)=\ue^{-2s}X^2+Y^2-Y$.
\end{proof}

\begin{proposition}[WS: the mollified flow at a weight]\label{prop:ws}
Let $W\in\mathcal W(\rho,K)$. There is $G=G(\rho,K,\chi)$ such that for $0<\varepsilon\le1$, $\hbar>0$, $s_0,s_1\ge0$ and $a\in L^2(\Hb)$,
\[
\norm{W\,V_\varepsilon(s_1,s_0)a}\le\ue^{G\hbar|s_1-s_0|/(2\varepsilon)}\norm{Wa}.
\]
Consequently, for $\gamma\ge0$ and $X_0^\gamma a\in L^2$,
$\norm{X_0^\gamma V_\varepsilon(s_1,s_0)a}\le\ue^{G_\gamma\hbar|s_1-s_0|/(2\varepsilon)}\norm{X_0^\gamma a}$ with $G_\gamma=G(\gamma,K_\gamma,\chi)$.
\end{proposition}
\begin{proof}
Let $v(s)=V_\varepsilon(s,s_0)a$ and $h=W^2$. Since $v'=\tfrac{i\hbar}2\Delta_\varepsilon(s)v$ in $L^2$ and $W$ is bounded,
$\tfrac d{ds}\norm{Wv}^2=-\hbar\,\mathrm{Im}\langle hv,\Delta_\varepsilon(s)v\rangle$. Write
\[
\langle hv,J_\varepsilon\Delta(s)J_\varepsilon v\rangle=\langle hJ_\varepsilon v,\Delta(s)J_\varepsilon v\rangle+\langle[J_\varepsilon,h]v,\Delta(s)J_\varepsilon v\rangle .
\]
For $f=J_\varepsilon v$, integrating by parts as in Lemma~\ref{lem:ell},
$\mathrm{Im}\langle hf,\Delta(s)f\rangle=-\ue^{-2s}\mathrm{Im}\langle(Xh)f,Xf\rangle-\mathrm{Im}\langle(Yh)f,Yf\rangle$, which for $s\ge0$ is at
most $4\rho\norm{Wf}\max_Z\norm{WZf}\le4\rho M^2\varepsilon^{-1}\norm{Wv}^2$ in absolute value, by Lemma~\ref{lem:wmoll}(i),(ii). By
Lemma~\ref{lem:wmoll}(iv) and its last claim, the second term is at most
$M\varepsilon\,\norm{W\Delta(s)J_\varepsilon v}\norm{Wv}\le3M^2\varepsilon^{-1}\norm{Wv}^2$. So
$|\tfrac d{ds}\norm{Wv}^2|\le G\hbar\varepsilon^{-1}\norm{Wv}^2$ with $G:=(4\rho+3)M^2$, at all times $s\ge0$, and Gr\"onwall's inequality gives
the first claim. The second follows with $W=W_{\gamma,n}$ and $n\to\infty$.
\end{proof}

\begin{lemma}[the path of the mollified flow]\label{lem:path}
Let $\gamma\ge0$, $0<\varepsilon\le1$, $\hbar>0$ and $s_1\ge0$ with $\hbar s_1\le\varepsilon$. Let $\mathcal S:=\{1,X,Y,X^2,Y^2\}$, let $b'$ be such that
$X_0^\gamma Zb'\in L^2(\Hb)$ for every $Z\in\mathcal S$, and put
$g(s):=V_\varepsilon(s,s_1)b'$ for $0\le s\le s_1$. There is $C=C(\gamma,\chi)$ such that:
\begin{itemize}
\item[(a)] $g(s)=b'+J_\varepsilon h(s)$ with $h(s):=\tfrac{i\hbar}2\int_{s_1}^s\Delta(r)J_\varepsilon g(r)\,dr$, where $h$ is Lipschitz from
$[0,s_1]$ to $L^2(X_0^{2\gamma}d\nu)$ and $\norm{X_0^\gamma h(s)}\le C\varepsilon^{-1}\norm{X_0^\gamma b'}$;
\item[(b)] the defect $d(s):=(\Delta(s)-\Delta_\varepsilon(s))g(s)$ satisfies
$\norm{X_0^\gamma d(s)}\le C\varepsilon^{-3}\sum_{Z\in\mathcal S}\norm{X_0^\gamma Zb'}$;
\item[(c)] $g$ is $C^1$ from $[0,s_1]$ to $L^2(X_0^{2\gamma}d\nu)$ with $g'=\tfrac{i\hbar}2\Delta_\varepsilon(s)g$, and $Zg$ is
continuous from $[0,s_1]$ to $L^2(X_0^{2\gamma}d\nu)$ for every $Z\in\mathcal S$.
\end{itemize}
\end{lemma}
\begin{proof}
The bounds of Lemma~\ref{lem:wmoll}(i)--(iii) and of Proposition~\ref{prop:ws} hold at the weight $X_0^\gamma$, with $M=M(\gamma,K_\gamma,\chi)$
and $G_\gamma$.
(a) Since $\partial_sg=\tfrac{i\hbar}2J_\varepsilon\Delta(s)J_\varepsilon g$ and $g(s_1)=b'$, the formula for $g$ holds. For $0\le r\le s_1$,
$\hbar|r-s_1|\le\varepsilon$, so $\norm{X_0^\gamma g(r)}\le\ue^{G_\gamma/2}\norm{X_0^\gamma b'}$ and
$\norm{X_0^\gamma\Delta(r)J_\varepsilon g(r)}\le3M\varepsilon^{-2}\ue^{G_\gamma/2}\norm{X_0^\gamma b'}$. Hence $h$ is Lipschitz, and
$\norm{X_0^\gamma h(s)}\le\tfrac12\hbar|s-s_1|\cdot3M\varepsilon^{-2}\ue^{G_\gamma/2}\norm{X_0^\gamma b'}\le\tfrac32M\ue^{G_\gamma/2}\varepsilon^{-1}\norm{X_0^\gamma b'}$.
(b) By (a),
$d(s)=\Delta(s)b'-J_\varepsilon\Delta(s)J_\varepsilon b'+\Delta(s)J_\varepsilon h(s)-J_\varepsilon\Delta(s)J_\varepsilon J_\varepsilon h(s)$. The first term is
bounded by the words $X^2b'$, $Y^2b'$, $Yb'$; the others by $3M^2\varepsilon^{-2}\norm{X_0^\gamma b'}$ and
$(3M+3M^3)\varepsilon^{-2}\norm{X_0^\gamma h(s)}$. (c) follows from (a) and Lemma~\ref{lem:wmoll}: $ZJ_\varepsilon$ is bounded on
$L^2(X_0^{2\gamma}d\nu)$ for $Z\in\mathcal S$, and so is $J_\varepsilon\Delta(s)J_\varepsilon$, continuously in $s$.
\end{proof}

% Appendix: Theorem 1.4 with explicit constants, the periodization lemmas (T2, G4.2) and Lemma EX.

\section{Companion results and analytic estimates}\label{sec:dd}

\subsection{Companion results for the reversed weight order}\label{sec:swapped}
The definition of the amplitude norm in \cite{S} and the proofs there take the two weights in different orders, and
\cite[Proposition~7.1]{S} states that the two orders give equivalent norms. For
$\beta>\tfrac12$ they do not: the two orders define different classes (Remark~\ref{pc:classes}). The principal result for the printed class is Theorem~\ref{pc:main}. In this subsection
\[
\norm a:=\norm{\ue^{\alpha\Lambda}\ue^{\beta\langle d\rangle}a}_{L^2(\Hb)} ,
\]
the order used in the proofs of \cite{S}. We call the amplitudes with $\norm a<\infty$ the \emph{swapped class}, and those with
$\norm a_{\alpha,\beta}<\infty$ the \emph{printed class}.
By Lemma~\ref{lem:m2w} below, a finite norm $\norm a$ at radius $\alpha$ gives a finite norm with the weight
$\ue^{\beta\tilde\psi}$ in place of $\ue^{\beta\langle d\rangle}$, at the radius $c_0\alpha$ for a fixed
$c_0=c_0(\beta)\in(0,1]$; the proofs use $\tilde\psi$.
\begin{theoremR}\label{thm:main}
Let $\Gamma$ be any Fuchsian group, $\beta>1$ and $\alpha\in(0,\alpha_0]$ constants ($\alpha_0<\pi/4$ fixed).
There are $C,\delta>0$ such that for all $0<\hbar\le1$, every $b\in\partial\Hb$ and every $a$ with $\norm a<\infty$,
\[
\norm{u^{(0)}(t)_\Gamma-\cU(t)[a\ue^{i\varphi_b/\hbar}]_\Gamma}_{L^2(\Hb/\Gamma)}\le C\hbar t\norm a,\qquad 0\le t\le\delta\hbar^{-1/2}.
\]
\end{theoremR}
For $\alpha>\alpha_0$ the class with radius $\alpha$ is contained in the class with radius $\alpha_0$
(in the swapped order this nesting is trivial), so the theorem holds for every constant $\alpha>0$,
as in \cite{S}. By Proposition~\ref{prop:g5} the order $\hbar t$ cannot be improved.

\begin{proof}
This is the case $K=N=0$ of Theorem~\ref{thm:dd14} below, with $\kappa:=\alpha^3$ and with $\alpha_0$ as here: then
$\alpha^3\ge\kappa\hbar$ for all $0<\hbar\le1$, $u^{(0)}$ is the approximation of order $0$, and the bound of Theorem~\ref{thm:dd14} reads
$C\alpha^{-2}\hbar t\norm a$ for $0\le t\le\delta\alpha^{1/2}\hbar^{-1/2}$. The constants of the present theorem are
$C\alpha^{-2}$ and $\delta\alpha^{1/2}$, with $C$ and $\delta$ the constants of Theorem~\ref{thm:dd14} for $\kappa=\alpha^3$; they depend on the fixed $\alpha$.
\end{proof}

\begin{remark}
No Volterra resummation and no Sobolev inequality on $M$ is used, so the proof holds for every Fuchsian
group; the cusp issue of Remark~\ref{rem:g4} does not arise. Since Theorem~\ref{thm:dd14} holds for every $\beta>1/2$, so does
the present theorem; we state it with $\beta>1$ as in \cite{S}.
\end{remark}

\medskip\noindent\textbf{The estimate of order $K$ and its corollaries.}
Let $u^{(K)}$ be the approximation of order $K$ of Section~\ref{sec:setup}.
\begin{theoremR}\label{thm:dd14}
Fix $\beta>1/2$, $\kappa>0$ and $\alpha_0\in(0,\min(1,\pi/4))$. There is $\delta>0$, and for every $P<\infty$ there is $C>0$,
such that for every Fuchsian group $\Gamma$ with $\sup_z\sum_{\gamma\in\Gamma}\ue^{-(\beta+1/2)\,d(x_0,\gamma z)}\le P$, all $0<\hbar\le1$, all $\alpha\in(0,\alpha_0]$ with $\alpha^3\ge\kappa\hbar$, all $K,N\in\mathbb N_0$, all $b\in\partial\Hb$, all $a$
with $\norm a<\infty$, and all $0\le t\le\delta(\alpha/\hbar)^{1/2}$,
\[
\norm{\Delta^N\big[\cU(t)[a\ue^{i\varphi_b/\hbar}]_\Gamma-u^{(K)}(t)_\Gamma\big]}_{L^2(\Hb/\Gamma)}
\le C^{N+K+1}N^{2N}K!\,(\alpha\hbar)^{-2N}\Big(\frac{\hbar t}{\alpha^2}\Big)^{K+1}\norm a ,
\]
with $N^{2N}:=1$ at $N=0$.
(The supremum is finite for every Fuchsian group, by Lemma~\ref{lem:p1} at the exponent $\beta+\tfrac12>1$. So $\delta$ does not
depend on $\Gamma$, and $C$ depends on $\Gamma$ only through this Poincar\'e series bound; see the end of the proof of
Theorem~\ref{thm:dd-appx}.)
\end{theoremR}
On Schubert's window $t\le\delta\alpha^{3/2}\hbar^{-1/2}$ this implies the bound \cite[(1.30)]{S} for the norm $\norm a$,
with an extra factor $(\alpha\hbar)^4\hbar t/\alpha^2\le1$. The proof is in Section~\ref{subsec:dd-thm}.

\begin{corollaryR}\label{cor:dd15}
Under the hypotheses of Theorem~\ref{thm:dd14}, with $\delta\le\kappa^{1/2}/\max(1,C)$, put $\delta_K=C\hbar t/\alpha^2$ and let $\epsilon\in(0,1)$. For every integer $K$ with
$(1-\epsilon)/\delta_K\le K\le(1-\epsilon/2)/\delta_K$,
\[
\norm{\Delta^N\big[\cU(t)[a\ue^{i\varphi_b/\hbar}]_\Gamma-u^{(K)}(t)_\Gamma\big]}\le C'\epsilon^{-1/2}C^{N+1}N^{2N}(\alpha\hbar)^{-2N}\ue^{-(1-\epsilon)\alpha^2/(C\hbar t)}\norm a ,
\]
with $C'$ an absolute constant.
\end{corollaryR}
The window for $K$ is the one used in Schubert's proof of his Corollary~1.5; the window printed in its statement
does not match that proof.

\begin{corollaryR}\label{cor:dd16}
If moreover the injectivity radius of $\Hb/\Gamma$ is at least $r_0>0$, then pointwise
\[
\big|\Delta^N\big[\cU(t)[a\ue^{i\varphi_b/\hbar}]_\Gamma-u^{(K)}(t)_\Gamma\big]\big|\le C^{N+K+1}(2N)!\,K!\,(\alpha\hbar)^{-2N-4}\Big(\frac{\hbar t}{\alpha^2}\Big)^{K+1}\norm a ,
\]
with $C$ depending also on $r_0$, and correspondingly for the window of Corollary~\ref{cor:dd15}.
\end{corollaryR}

\medskip\noindent\textbf{Notation for the analytic estimates.}
Let $X=y\partial_x$, $Y=y\partial_y$, $\Delta=X^2+Y^2-Y$, $\Delta(s)=\ue^{-2s}X^2+Y^2-Y$,
$\Lambda=\sqrt{-\Delta}$, $X_0=\cosh d(x_0,\cdot)$, $\tilde\psi=\log X_0$, and $W_\beta=X_0^\beta$.
In the analytic estimates below, the unadorned amplitude norm uses the weight $X_0^\beta$:
\[
\norm{a}:=\norm{\ue^{\alpha\Lambda}X_0^\beta a}.
\]
By Lemma~\ref{lem:m2w} a finite norm with the weight $\ue^{\beta\langle d\rangle}$ at radius $\alpha$ gives a finite norm of this
kind at the radius $c_0\alpha$; Theorem~\ref{thm:dd14} is deduced from Theorem~\ref{thm:dd-appx} in this way at the end of
Section~\ref{subsec:dd-thm}. The two-time mollified propagator $V_\varepsilon(s_1,s_0)=V_\varepsilon(s_1)V_\varepsilon(s_0)^*$ is defined in
Appendix~\ref{app:kl}.

\subsection{Exact identities}\label{subsec:dd-id}

Let $V(t)$ be the unitary operator of \eqref{eq:factor}, let $V(t,s):=V(t)V(s)^{-1}$, the propagator of
$i\partial_tV=-\tfrac{\hbar}{2}\Delta(t)V$ (\cite[(1.22)]{S}), and let $c:=i\hbar/2$. By \eqref{eq:factor} (which is \cite[(1.23)]{S}),
\[
\cU(t)(a\,\ue^{i\varphi/\hbar})=\ue^{-it/(2\hbar)}\bigl(S(t)V(t)a\bigr)\ue^{i\varphi/\hbar},
\qquad (S(t)w)(x,y)=\ue^{-t/2}w(x,\ue^{t}y),
\]
with $S(t)$ unitary. Then
\[S(t)^{-1}XS(t)=\ue^{-t}X,\qquad S(t)^{-1}YS(t)=Y,\]
hence $S(t)^{-1}\Delta S(t)=\Delta(t)$.
Multiplication by $\ue^{i\varphi/\hbar}=y^{-i/\hbar}$ preserves every Sobolev space $H^k$, and so do $\cU(t)$ and $S(t)$.
By the factorization, so do $V(t)^{\pm1}$ and $V(t,s)$.

\begin{proposition}[Exact identities (D1), (D3), (D4)]\label{prop:dd-id}
\begin{enumerate}
\item[\textup{(D1)}] \emph{Dyson expansion.} Let $a$ be such that every word in $X,Y$ applied to $a$ lies in $L^2(\Hb)$, and put
\[
P_k(t)=\int_{t>s_1>\cdots>s_k>0}\Delta(s_1)\cdots\Delta(s_k)\,ds,\qquad a^{(K)}(t)=\sum_{k=0}^{K}c^kP_k(t)a .
\]
Then, with $\Sigma_{K+1}(t)=\{t>s_1>\cdots>s_{K+1}>0\}$,
\[
a-V(t)^*a^{(K)}(t)=c^{K+1}\int_{\Sigma_{K+1}(t)}V(s_1)^*\,\Delta(s_1)\cdots\Delta(s_{K+1})a\,ds .
\]
This $P_k$ is exactly the $P_k$ of \cite{S}
(the display before the semiclassical approximation of order $K$). Every amplitude $a$ of Theorem~\ref{thm:dd-appx} satisfies the
hypothesis, by Lemma~\ref{lem:sc}.
\item[\textup{(D3)}] \emph{Transport of $\Delta^N$} (the conjugated operator is that of \cite[Proposition~2.2 and the proof of Proposition~3.10, (3.31)--(3.34)]{S}). Put $\Delta_\varphi:=e^{-i\varphi/\hbar}\Delta e^{i\varphi/\hbar}$, so that
$\hbar^2\Delta_\varphi=\hbar^2\Delta-2i\hbar Y-1+i\hbar$. Put $\Delta_\varphi(s):=S(s)^{-1}\Delta_\varphi S(s)$, so that
$\hbar^2\Delta_\varphi(s)=\hbar^2\Delta(s)-2i\hbar Y-1+i\hbar$. Then for every $b$ such that every word in $X,Y$ applied to $b$ lies in $L^2(\Hb)$,
\[
\Delta_\varphi^N V(s)^*b=V(s)^*\Delta_\varphi(s)^Nb.
\]
\item[\textup{(D4)}] \emph{The error.} Let $E:=\Delta^N\bigl[\cU(t)(a\,\ue^{i\varphi/\hbar})_\Gamma-u^{(K)}(t)_\Gamma\bigr]$, with $u^{(K)}$ as in \cite{S}. If every word in $X,Y$ applied to $a$ lies in $L^2_\omega$, for a weight $\omega$ as in Section~\ref{subsec:dd-surf}, then
\[
\begin{gathered}
\norm{E}_{L^2(M)}\le\Bigl(\frac{\hbar}{2}\Bigr)^{K+1}\int_{\Sigma_{K+1}(t)}\bigl\|\mathcal V_{s_1}b'(s)\bigr\|_{L^2(M)}\,ds,\\
b'(s):=\Delta_\varphi(s_1)^N\Delta(s_1)\cdots\Delta(s_{K+1})a.
\end{gathered}
\]
Everything of high order is now inside $b'(s)$, which is a fixed word-polynomial in $X,Y$ applied to $a$.
\end{enumerate}
\end{proposition}

\begin{proof}
\medskip\noindent\textbf{(D1).}

$V(s)^*=V(0,s)$ and $\partial_sV(t,s)=-c\,V(t,s)\Delta(s)$ on $H^\infty$. So for $b$ with all words in $L^2$, $s\mapsto V(s)^*b$ is
differentiable in $L^2$ with derivative $-c\,V(s)^*\Delta(s)b$. (In integrated form this is Lemma~\ref{lem:surf-duh} for the trivial
group, the weight $\omega=1$ and the constant path $g(s)=b$: then $\mathcal V_sg=\ue^{i\varphi/\hbar}V(s)^*g$ by \eqref{eq:factor}; the proof
of that lemma does not use (D1).)

$P_k(s)a$ is a finite sum of continuous scalar functions of $s$ times fixed words applied to $a$; $\partial_sP_k(s)=\Delta(s)P_{k-1}(s)$
and $P_k(0)=0$ for $k\ge1$.

With $J_k(t):=\int_0^tV(s)^*\Delta(s)P_k(s)a\,ds$ this gives $V(t)^*P_k(t)a=J_{k-1}(t)-cJ_k(t)$ for $k\ge1$, and $V(t)^*a-a=-cJ_0(t)$.

Multiplying by $c^k$ and summing telescopes to $a-V(t)^*a^{(K)}(t)=c^{K+1}J_K(t)$. Expanding $P_K$ gives the integral over $\Sigma_{K+1}(t)$.

\medskip\noindent\textbf{(D3).}

For $c'\in H^2$, the commutation of the propagator with the Laplacian gives the following identity.

$\cU(s)$ commutes with $\Delta$, so $\cU(s)(e^{i\varphi/\hbar}\Delta_\varphi c')=\Delta\,\cU(s)(e^{i\varphi/\hbar}c')$.

By \cite[(1.23)]{S} this reads $S(s)V(s)\Delta_\varphi c'=\Delta_\varphi S(s)V(s)c'$, i.e. $V(s)\Delta_\varphi=\Delta_\varphi(s)V(s)$ on $H^2$.

Multiplying by $V(s)^*$ on both sides is legitimate because $V(s)^*$ preserves $H^2$. Iterating $N$ times gives the assertion, since $b$ lies in every $H^k$.

\medskip\noindent\textbf{(D4).}

By definition $u^{(K)}(t)=\ue^{-it/(2\hbar)}\ue^{i\varphi/\hbar}S(t)a^{(K)}(t)$, so $u^{(K)}(t)_\Gamma=\cU(t)\mathcal V_ta^{(K)}(t)$. Since $\cU(t)$ is unitary on $L^2(M)$ and commutes with $\Delta_M$, $\norm E=\norm{\Delta_M^N(\mathcal V_0a-\mathcal V_ta^{(K)}(t))}$. (No commutation of $\cU(t)$ with periodization is used.)

By Corollary~\ref{cor:surf-dyson}, $\mathcal V_0a-\mathcal V_ta^{(K)}(t)=c^{K+1}\int_{\Sigma_{K+1}(t)}\mathcal V_{s_1}(\Delta(s_1)\cdots\Delta(s_{K+1})a)\,ds$.

By Lemma~\ref{lem:surf-lap} $N$ times, $\Delta_M^N\mathcal V_{s_1}g=\mathcal V_{s_1}\Delta_\varphi(s_1)^Ng$. Closedness of $\Delta_M^N$ and Minkowski's inequality give the bound, with $|c|^{K+1}=(\hbar/2)^{K+1}$.
(D1) and (D3) are identities on $\Hb$; on $M$ the proof of (D4) uses their surface versions in Section~\ref{subsec:dd-surf}. \qedhere
\end{proof}

\subsection{The surface object}\label{subsec:dd-surf}

Periodization needs decay, and Proposition~\ref{prop:ws} gives weighted bounds only for $V_\varepsilon$, not for $V$. So $[\ue^{i\varphi/\hbar}V(s)^*g]_\Gamma$ cannot be read as a convergent sum, and the identity $\cU(t)u_\Gamma=\sum_\gamma(\cU(t)u)\circ\gamma^{-1}$ of \cite{S} is not available for $u=\ue^{i\varphi/\hbar}V(s)^*g$. We therefore define this object on $M$ through the surface propagator, and prove the Duhamel formula on $M$. All estimates of the proofs below are unchanged.

Let $\omega\ge1$ be a locally bounded measurable weight on $\Hb$, $L^2_\omega:=\{f:\omega f\in L^2(\Hb)\}$, and $\delta_sz:=x+i\ue^sy$. Assume:
\begin{itemize}
\item[(W1)] $f\mapsto f_\Gamma$ is bounded from $L^2_\omega$ to $L^2(M)$, the sum converging absolutely a.e.;
\item[(W2)] $\sup\{\omega(z)/\omega(\delta_sz):|s|\le T,\ z\in\Hb\}<\infty$ for every $T$.
\end{itemize}
Lemma~\ref{lem:t2} (its proof, $t=0$) gives (W1) for $\omega=W_\sigma$, $\sigma>1/2$. Proposition~\ref{os:p1} gives it for the periodization weight $\omega=X_0^{1/2}\ell^{\,c'}$, $\ell=\log(2X_0)/\log2$, of Section~\ref{sec:printed}. Since $d(z,\delta_sz)=|s|$ and $\tilde\psi$ is $1$-Lipschitz, (W2) holds with the explicit ratios $W_\sigma(z)/W_\sigma(\delta_sz)\le\ue^{\sigma|s|}$ and $\omega(z)/\omega(\delta_sz)\le\ue^{|s|/2}(1+|s|/\log 2)^{c'}$. By (W2), $S(s)$ is bounded on $L^2_\omega$, locally uniformly in $s$; it is strongly continuous there, by density of $C_c(\Hb)$.

Let $\mathcal D_\omega$ be the set of $g\in L^2_\omega$ with $Xg,Yg,X^2g,Y^2g\in L^2_\omega$ (weak derivatives). For $g\in L^2_\omega$ and $s\in\mathbb R$ put
\[
\mathcal V_sg:=\ue^{-is/(2\hbar)}\,\cU(-s)\,\bigl[\ue^{i\varphi/\hbar}S(s)g\bigr]_\Gamma\in L^2(M),
\]
with $\cU$ the propagator on $M$. Thus $\mathcal V_0g=[\ue^{i\varphi/\hbar}g]_\Gamma$ and $\norm{\mathcal V_sg}_{L^2(M)}=\norm{[\ue^{i\varphi/\hbar}S(s)g]_\Gamma}_{L^2(M)}$. This $\mathcal V_sg$ is what $[\ue^{i\varphi/\hbar}V(s)^*g]_\Gamma$ formally means by \cite[(1.23)]{S}.

\begin{lemma}[Laplacian on the weighted class]\label{lem:surf-lap}
For $g\in\mathcal D_\omega$ and $s\in\mathbb R$, $[\ue^{i\varphi/\hbar}S(s)g]_\Gamma\in\operatorname{Dom}(\Delta_M)$ and $\Delta_M[\ue^{i\varphi/\hbar}S(s)g]_\Gamma=[\ue^{i\varphi/\hbar}S(s)\Delta_\varphi(s)g]_\Gamma$. Hence $\mathcal V_sg\in\operatorname{Dom}(\Delta_M)$ and $\Delta_M\mathcal V_sg=\mathcal V_s\Delta_\varphi(s)g$.
\end{lemma}

\begin{proof}
\begin{itemize}
\item Put $f:=\ue^{i\varphi/\hbar}S(s)g$. Then $Xf=\ue^{-s}\ue^{i\varphi/\hbar}S(s)Xg$, $Yf=\ue^{i\varphi/\hbar}S(s)(Y-i/\hbar)g$ and $\Delta f=\ue^{i\varphi/\hbar}S(s)\Delta_\varphi(s)g$. So $f,Xf,Yf,\Delta f\in L^2_\omega$ by (W2), and $f\in\operatorname{Dom}(\Delta)$ because $\Delta$ is essentially self-adjoint on $C_c^\infty(\Hb)$.
\item Take radial $\chi_n=\rho(d(x_0,\cdot)/n)$, with $\rho\in C_c^\infty$, $\rho=1$ near $0$. Then $0\le\chi_n\le1$, $\chi_n\to1$, and $|X\chi_n|+|Y\chi_n|+|\Delta\chi_n|\le C/n$. Since $\Delta(\chi_nf)=\chi_n\Delta f+2(X\chi_n\,Xf+Y\chi_n\,Yf)+(\Delta\chi_n)f$, dominated convergence gives $\chi_nf\to f$ and $\Delta(\chi_nf)\to\Delta f$ in $L^2_\omega$.
\item Fix $n$. Essential self-adjointness gives test functions $F_m\to\chi_nf$ in graph norm. With a fixed $\theta\in C_c^\infty(\Hb)$, $\theta=1$ on $\operatorname{supp}\chi_n$, the test functions $\theta F_m$ converge to $\chi_nf$ in graph norm, since $\norm{\nabla F}^2=\langle F,-\Delta F\rangle$ controls the cross term. Their supports lie in the compact set $\operatorname{supp}\theta$, where $\omega$ is bounded.
\item For a test function $F$, $F_\Gamma\in\operatorname{Dom}(\Delta_M)$ and $\Delta_MF_\Gamma=(\Delta F)_\Gamma$. By (W1) and closedness of $\Delta_M$, first in $m$ and then in $n$, the claim for $f$ follows. The statement for $\mathcal V_s$ follows since $\cU(-s)$ commutes with $\Delta_M$.\qedhere
\end{itemize}
\end{proof}

\begin{lemma}[Duhamel on $M$]\label{lem:surf-duh}
Let $g\in C^1([0,T];L^2_\omega)$ with $g(s)\in\mathcal D_\omega$ and $Yg,X^2g,Y^2g\in C([0,T];L^2_\omega)$. Then, with $c=i\hbar/2$,
\[
\mathcal V_Tg(T)-\mathcal V_0g(0)=\int_0^T\mathcal V_s\bigl(g'(s)-c\,\Delta(s)g(s)\bigr)\,ds\qquad\text{in }L^2(M).
\]
\end{lemma}

\begin{proof}
\begin{itemize}
\item For $g_0$ with $g_0,Yg_0\in L^2_\omega$: $S(s)g_0-S(r)g_0=\int_r^sS(\tau)(Y-\tfrac12)g_0\,d\tau$. This holds in $L^2(\Hb)$ (generator of the dilation group), the right side converges in $L^2_\omega$ by strong continuity, and $L^2_\omega\subset L^2$ since $\omega\ge1$.
\item Hence $s\mapsto S(s)g(s)$ is $C^1$ into $L^2_\omega$, and by (W1) $h(s):=[\ue^{i\varphi/\hbar}S(s)g(s)]_\Gamma$ is $C^1$ into $L^2(M)$ with $h'=[\ue^{i\varphi/\hbar}S(s)((Y-\tfrac12)g+g')]_\Gamma$.
\item By Lemma~\ref{lem:surf-lap}, $h(s)\in\operatorname{Dom}(\Delta_M)$ and $\Delta_Mh=[\ue^{i\varphi/\hbar}S(s)\Delta_\varphi(s)g]_\Gamma$, continuous in $s$. So $\tfrac{d}{ds}\cU(-s)h(s)=\cU(-s)\bigl(h'(s)-\tfrac{i\hbar}{2}\Delta_Mh(s)\bigr)$.
\item Since $\Delta_\varphi(s)=\Delta(s)-\tfrac{2i}{\hbar}(Y-\tfrac12)-\hbar^{-2}$, we have $-\tfrac{i\hbar}{2}\Delta_\varphi(s)=-c\Delta(s)-(Y-\tfrac12)+\tfrac{i}{2\hbar}$. The $(Y-\tfrac12)$ terms cancel, and $\tfrac{i}{2\hbar}$ cancels the derivative of $\ue^{-is/(2\hbar)}$. Thus $\tfrac{d}{ds}\mathcal V_sg(s)=\mathcal V_s(g'(s)-c\Delta(s)g(s))$; integrate.\qedhere
\end{itemize}
\end{proof}

\begin{corollary}[Dyson expansion on $M$]\label{cor:surf-dyson}
If every word in $X,Y$ applied to $a$ lies in $L^2_\omega$, then with $a^{(K)}(t)=\sum_{k\le K}c^kP_k(t)a$,
\[
\mathcal V_0a-\mathcal V_ta^{(K)}(t)=c^{K+1}\int_{\Sigma_{K+1}(t)}\mathcal V_{s_1}\bigl(\Delta(s_1)\cdots\Delta(s_{K+1})a\bigr)\,ds.
\]
\end{corollary}

\begin{proof}
$P_k(s)a$ is a finite sum of continuous scalar functions times fixed words applied to $a$, so Lemma~\ref{lem:surf-duh} applies to it. Since $\partial_sP_{k}(s)=\Delta(s)P_{k-1}(s)$ for $k\ge1$ and $P_k(0)=0$, with $J_k(t):=\int_0^t\mathcal V_s\Delta(s)P_k(s)a\,ds$ we get $\mathcal V_tP_k(t)a=J_{k-1}(t)-cJ_k(t)$ for $k\ge1$, and $\mathcal V_ta-\mathcal V_0a=-cJ_0(t)$. Multiplying by $c^k$ and summing telescopes to $\mathcal V_0a-\mathcal V_ta^{(K)}(t)=c^{K+1}J_K(t)$. Expanding $P_K$ gives the integral over $\Sigma_{K+1}(t)$.
\end{proof}

\subsection{First-order Cauchy estimates (Lemma FC)}\label{subsec:dd-fc}

\begin{lemma}[FC]\label{lem:fc}
For $\tau_0<\pi/2$ there is $C_1=C_1(\tau_0)$ such that for $Z\in\{X,Y\}$ and $0\le\alpha''<\alpha'\le\tau_0$,
\[
\norm{\ue^{\alpha''\Lambda}Z\,\ue^{-\alpha'\Lambda}}_{L^2(\Hb)}\le\frac{C_1}{\alpha'-\alpha''}.
\]
\end{lemma}

\begin{proof}
\medskip\noindent\textbf{Reduction.}

We first Fourier-transform in $x$. On the mode $\ue^{ikx}$, $k\ne0$, put $z=|k|y$ and write the function of $y$
as $z^{1/2}f(z)$, with $f\in L^2(dz/z)$. Then $-\Delta-\tfrac14=-(z\partial_z)^2+z^2$,
$X=i\,\operatorname{sgn}(k)\,z$ is multiplication, and $Y=z\partial_z+\tfrac12$.
Thus $Y_0:=Y-\tfrac12=z\partial_z=:D_z$ is skew-adjoint on $L^2(dz/z)$.

In the unitary Kontorovich--Lebedev picture
\[G=w(\nu)^{1/2}F,\qquad F(\nu)=\int_0^\infty K_{i\nu}(z)f(z)\,dz/z,\qquad w(\nu)=\tfrac{2}{\pi^2}\nu\sinh\pi\nu\]
(\cite{Y}), $\Lambda$ acts as multiplication by $\Lambda(\nu)=(\nu^2+\tfrac14)^{1/2}$.

The operators are the same for every $k$, so bounds are uniform in $k$. On the $k=0$ mode, $X=0$, while $Y_0$ and $\Lambda$ commute: in $r=\log y$, $Y_0=\partial_r$ and $\Lambda=(-\partial_r^2+\tfrac14)^{1/2}$. So
\[\norm{\ue^{\alpha''\Lambda}Y_0\ue^{-\alpha'\Lambda}}\le\norm{\Lambda\ue^{-(\alpha'-\alpha'')\Lambda}}\]
there. The mode $k=0$ has measure zero in the direct integral.

Put $\gamma:=\alpha'-\alpha''$ and write
\[
\ue^{\alpha''\Lambda}Z\ue^{-\alpha'\Lambda}=\bigl(\ue^{\alpha''\Lambda}Z\ue^{-\alpha''\Lambda}\Lambda^{-1}\bigr)\bigl(\Lambda\ue^{-\gamma\Lambda}\bigr).
\]
The second factor has norm $\le\sup_{\lambda\ge1/2}\lambda\ue^{-\gamma\lambda}\le1/(e\gamma)$. It remains to bound the first factor uniformly for $\alpha''\le\tau_0$, and $Y=Y_0+\tfrac12$ adds at most $\tfrac12$.

\medskip\noindent\textbf{Spectral notation.} Put $p:=\pi(\nu+\nu')/2$, $q:=\pi(\nu-\nu')/2$, $x:=\nu-\nu'$. Since $\sinh\pi\nu\,\sinh\pi\nu'=\sinh(p+q)\sinh(p-q)$,
\[
(w(\nu)w(\nu'))^{1/2}=\tfrac{2}{\pi^2}(\nu\nu')^{1/2}(\sinh^2p-\sinh^2q)^{1/2}=\tfrac{2}{\pi^2}(\nu\nu')^{1/2}(\cosh^2p-\cosh^2q)^{1/2}.
\]

\medskip\noindent\textbf{Mellin product formula.} For $\operatorname{Re}s>|\operatorname{Re}a|+|\operatorname{Re}b|$,
\[
M_s(a,b):=\int_0^\infty z^{s-1}K_a(z)K_b(z)\,dz=2^{s-3}\,\frac{\Gamma\bigl(\tfrac{s+a+b}{2}\bigr)\Gamma\bigl(\tfrac{s+a-b}{2}\bigr)\Gamma\bigl(\tfrac{s-a+b}{2}\bigr)\Gamma\bigl(\tfrac{s-a-b}{2}\bigr)}{\Gamma(s)}.
\]
\emph{Proof of the formula.} Write $K_b=K_{-b}$ (\cite[10.27.3]{DLMF}) and use the integral representation
\cite[10.32.10]{DLMF} for both Bessel factors. The integration in $z$ is
\[
\int z^{m-1}\ue^{-z^2c/4}\,dz=\tfrac12(4/c)^{m/2}\Gamma(m/2),
\qquad m=s+a-b,\quad c=1/u+1/v.
\]
Substitute $u=rt$ and $v=r(1-t)$ in the remaining double integral. The $r$-integral gives
$\Gamma(\tfrac{s-a+b}{2})$, and the $t$-integral gives the beta function
$\Gamma(\tfrac{s-a-b}{2})\Gamma(\tfrac{s+a+b}{2})/\Gamma(s)$, yielding the formula above.
The absolute integrand is the same expression with $a,b$ replaced by $\operatorname{Re}a,\operatorname{Re}b$ (since $|(z/2)^a|=(z/2)^{\operatorname{Re}a}$). For $\operatorname{Re}s>|\operatorname{Re}a|+|\operatorname{Re}b|$ the same computation then gives a finite value with all Gamma arguments of positive real part, so Tonelli justifies every exchange. This covers $a=i\nu$, $b=i\nu'$ at $s=1$, and $b=i\nu'\pm1$ at $s=1+\eta$, $\eta>0$. \qed

\medskip\noindent\textbf{(X) Kernel of $z$.}

Since the measure is $dz/z$, the kernel of multiplication by $z$ is $(w w')^{1/2}M_1(i\nu,i\nu')$.

The identity $|\Gamma(\tfrac12+iy)|^2=\pi/\cosh\pi y$ gives
\[
Zt_1(\nu,\nu')=(ww')^{1/2}\frac{\pi^2}{4\cosh p\cosh q}=\frac{(\nu\nu')^{1/2}(\cosh^2p-\cosh^2q)^{1/2}}{2\cosh p\cosh q},
\]
so $|Zt_1|\le(\nu\nu')^{1/2}/(2\cosh q)$.

The kernel of $\ue^{\alpha''\Lambda}z\ue^{-\alpha''\Lambda}\Lambda^{-1}$ is bounded by $\ue^{\alpha''|x|}(\nu\nu')^{1/2}/(2\cosh(\pi x/2)\Lambda(\nu'))$.

Since $(\nu\nu')^{1/2}\le(\nu+\nu')/2\le\nu'+|x|/2$ and $\Lambda(\nu')\ge\max(\nu',1/2)$, we have
\[
(\nu\nu')^{1/2}/\Lambda(\nu')\le1+|x|.
\]

So the kernel is at most $g(x):=(1+|x|)\ue^{\tau_0|x|}/(2\cosh(\pi x/2))$, which is integrable for $\tau_0<\pi/2$. The Schur test gives the bound $\int g$.

\medskip\noindent\textbf{(Y) Kernel of $D_z$.} The integral $\int K_{i\nu}K'_{i\nu'}\,dz$ does not converge at $z=0$, so we regularize.

Let $F,G\in C_c^\infty((0,\infty))$ be the KL transforms of $f,g$ (the core $\mathcal C$).

On $\mathcal C$, $f(e^{-L})$ and $zg'(z)|_{z=e^{-L}}$ decay faster than any power of $L$ as $L\to\infty$: they are integrals of smooth compactly supported densities against $\ue^{\pm i\nu L}$-oscillations, plus $O(\ue^{-2L})$. This uses $K_{i\nu}=\tfrac12\pi(I_{-i\nu}-I_{i\nu})/\sin(i\nu\pi)$ (\cite[10.27.4]{DLMF}) with the series of $I_{\pm i\nu}$ (\cite[10.25.2]{DLMF}), and $1/\sin(i\nu\pi)$ is smooth on $\operatorname{supp}F\subset(0,\infty)$.

Hence $\langle f,D_zg\rangle=\int_0^\infty\bar f g'\,dz=\lim_{\eta\to0^+}\int z^\eta\bar fg'\,dz$, by dominated convergence.

For $\eta>0$ Fubini applies, and $K'_\nu=-(K_{\nu-1}+K_{\nu+1})/2$ (\cite[10.29.1]{DLMF}, with $\mathcal Z_\nu=\ue^{\nu\pi i}K_\nu$ in DLMF's convention, under which 10.29.1 reads $\mathcal Z_{\nu-1}+\mathcal Z_{\nu+1}=2\mathcal Z'_\nu$ and gives the minus sign for $K$) gives
\[
\begin{gathered}
\langle f,D_zg\rangle=\lim_{\eta\to0^+}\iint\bar F(\nu)G(\nu')\,w(\nu)w(\nu')\,I_\eta(\nu,\nu')\,d\nu\,d\nu',\\
I_\eta=-\tfrac12\bigl[M_{1+\eta}(i\nu,i\nu'+1)+M_{1+\eta}(i\nu,i\nu'-1)\bigr].
\end{gathered}
\]
Both terms are admissible, since $\operatorname{Re}s=1+\eta>0+1$.

In $M_{1+\eta}(i\nu,i\nu'+1)$ the factor $\Gamma((\eta+ix)/2)=(2/(\eta+ix))\Gamma(1+(\eta+ix)/2)$ is singular, and the other factors are smooth. In $M_{1+\eta}(i\nu,i\nu'-1)$ the singular factor is $\Gamma((\eta-ix)/2)$.

So $I_\eta=\frac{2}{\eta+ix}A_\eta+\frac{2}{\eta-ix}B_\eta$, where $A_\eta,B_\eta$ are smooth in $(\nu,\nu',\eta)$ on the support.

At $\eta=0$ and $x=0$, the reflection identity $\Gamma(1+w)\Gamma(-w)=-\pi/\sin\pi w$ gives
\[
\begin{aligned}
A_0&=-\tfrac12\cdot\tfrac14\,\Gamma(1+ip/\pi)\Gamma(-ip/\pi)=-\tfrac18\cdot\tfrac{i\pi}{\sinh p},\\
B_0&=-\tfrac12\cdot\tfrac14\,\Gamma(ip/\pi)\Gamma(1-ip/\pi)=-\tfrac18\cdot\tfrac{-i\pi}{\sinh p}.
\end{aligned}
\]
So $A_0+B_0=0$ on the diagonal.

The Sokhotski--Plemelj formula gives $\frac{2}{\eta\pm ix}\to\mp2i\,\mathrm{PV}\frac1x+2\pi\delta(x)$. The $\delta$-terms carry $A_0+B_0$ at $x=0$, so they cancel.

For $x\neq0$ the pointwise limit is $I_0=-\pi^2/(4\sinh p\,\sinh q)$ (same Gamma reflections). Hence, on $\mathcal C\times\mathcal C$,
\[
\langle f,D_zg\rangle=\mathrm{PV}\iint\bar F G\,w\,w'\,I_0,
\qquad\text{i.e.}\qquad
\begin{gathered}
\tilde D(\nu,\nu')=-\frac{(\nu\nu')^{1/2}\rho(\nu,\nu')}{2\sinh q},\\
\rho:=\Bigl(1-\frac{\sinh^2q}{\sinh^2p}\Bigr)^{1/2}\in[0,1].
\end{gathered}
\]

This kernel is real and antisymmetric, as $D_z$ must be.

As a consistency check, not used in the proof, $(\nu^2-\nu'^2)\tilde D=-2Zt_2$, where
$Zt_2=(w(\nu)w(\nu'))^{1/2}M_2(\nu,\nu')$ is the kernel of $z^2$, with $M_2$ as in
Remark~\ref{rem:fc-y-alt}. This agrees with $[\Lambda^2,D_z]=[z^2,D_z]=-2z^2$.

\medskip\noindent\textbf{(Y) bound.} The kernel of $Q:=\ue^{\alpha''\Lambda}D_z\ue^{-\alpha''\Lambda}\Lambda^{-1}$ is $\ue^{\alpha''(\Lambda(\nu)-\Lambda(\nu'))}\tilde D(\nu,\nu')/\Lambda(\nu')$. It is valid on $\mathcal C$, because $\ue^{\pm\alpha''\Lambda}$ and $\Lambda^{-1}$ preserve $\mathcal C$. Split $Q=Q_{\rm far}+Q_{\rm near}$, with $Q_{\rm near}:=Q\cdot\mathbf 1_{|x|\le1}$.

\textbf{$Q_{\rm far}$.} Here $|Q|\le(1+|x|)\ue^{\tau_0|x|}/(2\sinh(\pi|x|/2))$ on $|x|>1$. This is integrable, so the Schur test applies.

\textbf{$Q_{\rm near}$.} Write $1/(2\sinh(\pi x/2))=1/(\pi x)+r(x)$, with $r$ bounded on $|x|\le1$.

The $r$-part has a bounded kernel supported in $|x|\le1$, so the Schur test applies.

In the main part $-a(\nu,\nu')/(\pi(\nu-\nu'))$, put
\[a(\nu,\nu'):=\ue^{\alpha''(\Lambda(\nu)-\Lambda(\nu'))}(\nu\nu')^{1/2}\rho(\nu,\nu')/\Lambda(\nu')\]
Then $a(\nu,\nu')=a(\nu',\nu')+[a(\nu,\nu')-a(\nu',\nu')]$, with $a(\nu',\nu')=\nu'/\Lambda(\nu')\in[0,1]$.

The diagonal part is the truncated Hilbert transform $\mathrm{PV}\,\frac{1}{\pi x}\mathbf 1_{|x|\le1}$
composed with multiplication by $\nu'/\Lambda(\nu')$. This transform is bounded on $L^2(\mathbb R)$:
its Fourier multiplier $-\tfrac{2i}{\pi}\operatorname{Si}(|\xi|)\operatorname{sgn}\xi$ has modulus
at most $(2/\pi)\operatorname{Si}(\pi)<1.2$. Extending functions by zero from $(0,\infty)$ therefore gives the required bound.

For the off-diagonal remainder, write $s(\nu):=(\nu\nu')^{1/2}/\Lambda(\nu')$ and
$E:=\ue^{\alpha''(\Lambda(\nu)-\Lambda(\nu'))}$. Then
\[
a(\nu,\nu')-a(\nu',\nu')=(E-1)s\rho+(\rho-1)s+(s(\nu)-s(\nu')).
\]
For $|x|\le1$, the first term is controlled by $|E-1|\le\tau_0\ue^{\tau_0}|x|$ and
$s\rho\le s\le1+|x|\le2$. For the second, $\sinh|q|/\sinh p\le|q|/p$ when $|q|\le p$ gives
$1-\rho\le\sinh^2q/\sinh^2p\le|q|/p=|x|/(\nu+\nu')$, and hence
$(1-\rho)s\le|x|(\nu\nu')^{1/2}/((\nu+\nu')\Lambda(\nu'))\le|x|$.
Finally, $|s(\nu)-s(\nu')|=|\sqrt\nu-\sqrt{\nu'}|\sqrt{\nu'}/\Lambda(\nu')\le|x|^{1/2}$,
since $\sqrt{\nu'}/\Lambda(\nu')\le1$.
So the remainder kernel is $\le C|x|^{-1/2}$ on $|x|\le1$, which is integrable. The Schur test therefore bounds the remainder.
Combining the near and far bounds gives $\norm{Q}\le C(\tau_0)$ on $\mathcal C$, so $Q$ extends to a bounded operator. The bound for $\ue^{\alpha''\Lambda}D_z\ue^{-\alpha'\Lambda}$ follows on the dense set $\Lambda\ue^{-\gamma\Lambda}\mathcal C$, and then by closure. \qed
\end{proof}

\begin{remark}[Proof of (Y) without a principal value]\label{rem:fc-y-alt}
The case $Z=Y$ of Lemma~\ref{lem:fc} also follows without the kernel of $D_z$. For $A,B\in C_c^\infty((0,\infty))$ put $a=wA$, $b=wB$, $\hat a(t)=\int a(\nu)\ue^{i\nu t}\,d\nu$, $\psi_a(z)=\int a(\nu)K_{i\nu}(z)\,d\nu=\tfrac12\int\ue^{-z\cosh t}\hat a(t)\,dt$ and $S(A,B)=\langle\psi_a,D_z\psi_b\rangle_{L^2(dz/z)}$.
\begin{itemize}
\item By Fubini, $S(A,B)=-\tfrac14\iint\overline{\hat a(t)}\,\hat b(u)\,\frac{\cosh u}{\cosh t+\cosh u}\,dt\,du$, absolutely convergent.
\item Since $(\partial_u^2-\partial_t^2)\frac{\cosh u}{\cosh t+\cosh u}=2(\cosh t+\cosh u)^{-2}$, integrating by parts twice in each variable gives
\[S(\Lambda^2A,B)-S(A,\Lambda^2B)=-2\iint\bar A(\nu)B(\nu')\,w(\nu)w(\nu')\,M_2(\nu,\nu')\,d\nu\,d\nu',\]
where $M_2(\nu,\nu')=\int_0^\infty zK_{i\nu}K_{i\nu'}\,dz=\tfrac12\,\mathsf g(\tfrac{\nu+\nu'}2)\,\mathsf g(\tfrac{\nu-\nu'}2)$ and $\mathsf g(y)=\pi y/\sinh\pi y$ (the Mellin product formula at $s=2$).
\item The multipliers $\ue^{\pm s\Lambda}$ preserve $C_c^\infty((0,\infty))$. Integrating first along $\ue^{-t\Lambda}$ in both arguments over $t\in(0,\infty)$, then along $(\ue^{s\Lambda},\ue^{-s\Lambda})$ over $s\in(0,a)$, gives
\[S(\ue^{a\Lambda}A,\ue^{-a\Lambda}B)-S(A,B)=\iint\bar A(\nu)B(\nu')\,w(\nu)w(\nu')\,R_a(\nu,\nu')\,d\nu\,d\nu',\]
with $R_a=-2M_2\,\phi_a(\Lambda(\nu)-\Lambda(\nu'))/(\Lambda(\nu)+\Lambda(\nu'))$ and $\phi_a(\delta)=\int_0^a\ue^{s\delta}\,ds$.
\item Hence, on products in the KL picture and then by density, $\ue^{a\Lambda}Y\ue^{-(a+\gamma)\Lambda}=Y\ue^{-\gamma\Lambda}+T_{R_a}\ue^{-\gamma\Lambda}$, where $T_RG(\nu)=\int R(\nu,\nu')G(\nu')w(\nu')\,d\nu'$. The first term has norm at most $2/\gamma$, by $\norm{Yu}^2\le\norm{u}\,\norm{\Delta u}$. For the second, $\mathsf g(p')\cosh(\pi p')\le1+\pi p'$ with $p'=(\nu+\nu')/2$, together with the inequalities used in (X), gives
\[(w(\nu)w(\nu'))^{1/2}\,|R_a(\nu,\nu')|\,\ue^{-\gamma\Lambda(\nu')}\le C(\tau_0)\,\gamma^{-1}(1+|x|)\ue^{\tau_0|x|}\,\mathsf g(x/2),\qquad x=\nu-\nu',\]
which is integrable for $\tau_0<\pi/2$, and the Schur test concludes.
\end{itemize}
\end{remark}

\subsection{Unweighted Cauchy estimate for words (Lemma WC)}\label{subsec:dd-wc}

\begin{lemma}[WC]\label{lem:wc}
For a word $Z_1\cdots Z_m$ with $Z_i\in\{X,Y\}$, and $0\le\alpha''<\alpha\le\tau_0<\pi/2$:
\[
\norm{\ue^{\alpha''\Lambda}Z_1\cdots Z_m\ue^{-\alpha\Lambda}}\le\Bigl(\frac{C_1m}{\alpha-\alpha''}\Bigr)^m\le\Bigl(\frac{eC_1}{\alpha-\alpha''}\Bigr)^m m!.
\]
\end{lemma}

\begin{proof}
Take equally spaced radii $\alpha''=r_0<r_1<\cdots<r_m=\alpha$, and insert $\ue^{-r_i\Lambda}\ue^{r_i\Lambda}$.
Each factor $\ue^{r_{i-1}\Lambda}Z_i\ue^{-r_i\Lambda}$ is bounded by $C_1m/(\alpha-\alpha'')$ (Lemma~\ref{lem:fc}). The identity extends from $\mathcal C$ by closure, since each $\ue^{r_i\Lambda}$ is closed. Finally $m^m\le e^m m!$. \qed
\end{proof}

\subsection{Static weighted Cauchy estimates (Lemma SC)}\label{subsec:dd-sc}

Let $\mathcal A:=\operatorname{span}\{y,\;x,\;(x^2+1)/y\}$.
\begin{itemize}
\item $\mathcal A$ contains $X_0=\tfrac12\bigl((x^2+1)/y+y\bigr)$.
\item $\mathcal A$ is invariant under $X$ and $Y$: $Xy=0$, $Xx=y$, $X((x^2+1)/y)=2x$; $Yy=y$, $Yx=0$, $Y((x^2+1)/y)=-(x^2+1)/y$.
\item In this basis the $\ell^1$-operator norms are $\norm{X}\le2$ and $\norm{Y}\le1$.
\item Since $X_0\ge y/2$, $X_0\ge|x|$ (because $x^2+y^2+1-2y|x|=(|x|-y)^2+1$), and $X_0\ge(x^2+1)/(2y)$, every $v=c_1y+c_2x+c_3(x^2+1)/y$ satisfies $|v|\le2\norm{c}_1X_0$.
\item Hence $|Z_BX_0|\le2\cdot2^{|B|}X_0$ for every word $Z_B$.
\end{itemize}

\begin{lemma}[Faà di Bruno bound]\label{lem:fdb}
For $\beta\ge0$ and a word $Z_S$ of length $j$:
\[
|X_0^\beta Z_SX_0^{-\beta}|\le B_\beta\,j!\,\rho_\beta^{-j},\qquad \rho_\beta:=\tfrac14\log(3/2)
\]
(any fixed value below $\tfrac12\log(3/2)$ works), where $B_\beta$ depends only on $\beta$.
\end{lemma}

\begin{proof}
\begin{itemize}
\item Expand by induction on $j$: $Z_S\varphi(X_0)=\sum_{\pi}\varphi^{(|\pi|)}(X_0)\prod_{B\in\pi}Z_BX_0$, the sum over partitions $\pi$ of $S$, where each block $B$ inherits the order of $S$.
\item With $\varphi(u)=u^{-\beta}$, $|\varphi^{(k)}(u)|=(\beta)_ku^{-\beta-k}$. So the bound is $2^j\sum_kS(j,k)(\beta)_k2^k$.
\item The exponential generating function of $\sum_kS(j,k)(\beta)_k2^k$ is $(3-2\ue^z)^{-\beta}$, which is holomorphic on $|z|<\log(3/2)$.
\item Cauchy's estimate on $|z|=\rho<\log(3/2)$ gives the bound $C_{\beta,\rho}\,j!\,\rho^{-j}$. Combined with the factor $2^j$, this gives $\rho_\beta=\rho/2$. \qed
\end{itemize}
\end{proof}

\begin{lemma}[SC]\label{lem:sc}
For $\beta\ge0$, a word $Z_1\cdots Z_m$ and $0<\alpha\le\tau_0$, with $F:=X_0^\beta a$ and $\norm{\ue^{\alpha\Lambda}F}<\infty$:
\[
\norm{X_0^\beta Z_1\cdots Z_m a}\le B_\beta(m+1)\,m!\,(C_3/\alpha)^m\,\norm{\ue^{\alpha\Lambda}X_0^\beta a},
\qquad C_3:=\max(\ue C_1,\tau_0/\rho_\beta).
\]
\end{lemma}

\begin{proof}
\begin{itemize}
\item Leibniz for words of derivations: \[X_0^\beta Z_1\cdots Z_m(X_0^{-\beta}F)=\sum_{S\subseteq[m]}(X_0^\beta Z_SX_0^{-\beta})\,(Z_{S^c}F)\]
pointwise.
\item Bound each term by $B_\beta|S|!\rho_\beta^{-|S|}\norm{Z_{S^c}F}$, and the second factor by Lemma~\ref{lem:wc} with $\alpha''=0$: $\norm{Z_{S^c}F}\le(\ue C_1/\alpha)^{m-|S|}(m-|S|)!\,\norm{\ue^{\alpha\Lambda}F}$.
\item Since $\sum_{|S|=j}j!(m-j)!=m!$, the sum is $\le B_\beta\,m!\sum_j\rho_\beta^{-j}(\ue C_1/\alpha)^{m-j}\le B_\beta(m+1)m!(C_3/\alpha)^m$. \qed
\end{itemize}
\end{proof}

\subsection{\texorpdfstring{Bounds for $b'(s)$ (Lemma WB)}{Bounds for b'(s) (Lemma WB)}}\label{subsec:dd-wb}

\begin{lemma}[WB]\label{lem:wb}
Assume $0<\hbar\le1$ and $0<\alpha\le\alpha_0\le\min(1,\tau_0)$. Let $J$ be a word of length $\le n$, and let $\mathfrak D^J:=Z_1\cdots Z_m$ with $Z_i\in\{X,Y\}$ and $m=|J|$ be the ordered word it determines (the empty word is the identity). Uniformly in $s\in\Sigma_{K+1}(t)$:
\begin{itemize}
\item \textup{(WB-w)} $\norm{X_0^\beta\mathfrak D^Jb'(s)}\le B_\beta C_4^{N+K+n+1}(2N)!\,(n+2K+2)!\,\alpha^{-(n+2K+2)}(\alpha\hbar)^{-2N}\norm a$;
\item \textup{(WB-u)} for $0\le\alpha_1<\alpha_0'\le\alpha_0$ with $\alpha_0'-\alpha_1\ge\alpha/4$:
\[
\norm{\ue^{\alpha_1\Lambda}b'(s)}\le C_4^{N+K+1}(2N)!\,(2K+2)!\,\alpha^{-(2K+2)}(\alpha\hbar)^{-2N}\norm{\ue^{\alpha_0'\Lambda}a}.
\]
\end{itemize}
\end{lemma}

\begin{proof}
\begin{itemize}
\item Expand $b'$ as a sum over choices of one term from each factor:
  \begin{itemize}
  \item from each $\Delta(s_j)$: $\ue^{-2s_j}X^2$, $Y^2$ or $-Y$ (coefficients of modulus $\le1$, length $\le2$);
  \item from each $\Delta_\varphi(s_1)$: $\hbar^{-2}(-1+i\hbar)$ (length $0$, $|\text{coeff}|\le2\hbar^{-2}$), $-2i\hbar^{-1}Y$ (length $1$, $|\text{coeff}|=2\hbar^{-1}$), or one of the three length-$\le2$ terms (coefficient $\le1$).
  \end{itemize}
\item Suppose $n_0,n_1,n_2$ of the $N$ $\Delta_\varphi$-factors choose these three kinds. The coefficient is $\le2^{n_0+n_1}\hbar^{-2n_0-n_1}$, and the word length is $\le L:=n+\ell+2K+2$, with $\ell:=n_1+2n_2\le2N$.
\item Since $2n_0+n_1=2N-\ell$, the coefficient is $\le2^N\hbar^{-2N+\ell}$.
\item By Lemma~\ref{lem:sc}, and since $m\mapsto(m+1)m!(C_3/\alpha)^m$ is increasing, each term is $\le B_\beta(L+1)L!(C_3/\alpha)^L\norm a$.
\item Use $L!\le2^L\ell!(n+2K+2)!$, and then
\[
\begin{aligned}
\hbar^{-2N+\ell}\ell!(C_3/\alpha)^\ell&=\hbar^{-2N}\ell!(C_3\hbar/\alpha)^\ell\\
&\le\hbar^{-2N}(2N)!\max(1,C_3\hbar/\alpha)^{2N}\le(2N)!\,C_3^{2N}(\alpha\hbar)^{-2N},
\end{aligned}
\]
using $\alpha,\hbar\le1\le C_3$.
\item The per-factor weights for the $\Delta_\varphi$-factors are $2+2+1+1+1=7$ (weight $2$ for the length-$0$ and length-$1$ choices, weight $1$ for each of the three length-$\le2$ choices), so the weighted number of choices is $\le7^N3^{K+1}$, the factor $3^{K+1}$ coming from the $\Delta(s_j)$. Collecting the powers of $2$ and $7$ into $C_4$ gives (WB-w).
\item (WB-u) is the same count with Lemma~\ref{lem:wc} (gap $\ge\alpha/4$) in place of Lemma~\ref{lem:sc}, and $n=0$. \qed
\end{itemize}
\end{proof}

\subsection{The mollifier defect on analytic vectors (Lemma UD)}\label{subsec:dd-unw}

\begin{lemma}[UD]\label{lem:unw}
There is a constant $C_0$ such that for $0<\varepsilon\le1$, $0<a\le1$, $s\ge0$ and $z\in\operatorname{Dom}(\ue^{a\Lambda})$,
\[
\norm{(\Delta_\varepsilon(s)-\Delta(s))z}\le C_0\,a^{-4}\,\ue^{-a/(8\varepsilon)}\,\norm{\ue^{a\Lambda}z}.
\]
\end{lemma}

\begin{proof}
\begin{itemize}
\item \emph{The symbol.} For $\lambda\ge0$: $(1+\lambda^2)(1-j_\varepsilon(\lambda))\le5\,a^{-2}\ue^{-a/(4\varepsilon)}\ue^{a\lambda}$. Indeed
$0\le1-j_\varepsilon(\lambda)=\int g(\lambda-\mu)(1-\chi(\varepsilon\mu))\,d\mu\le\int_{|\mu|\ge1/\varepsilon}g(\lambda-\mu)\,d\mu$.
  \begin{itemize}
  \item If $\lambda\le1/(2\varepsilon)$, then $|\lambda-\mu|\ge1/(2\varepsilon)$ on the range of integration, so
  $1-j_\varepsilon(\lambda)\le2\ue^{-1/(8\varepsilon^2)}$. With $u=1/\varepsilon\ge1$ the left side is at most
  $2(1+u^2/4)\ue^{-u^2/8}\le5\ue^{-u/4}\le5\ue^{-a/(4\varepsilon)}$, because $2(1+u^2/4)\ue^{-(u^2-2u)/8}<5$ for $u\ge1$.
  \item If $\lambda>1/(2\varepsilon)$, then $1-j_\varepsilon\le1$ and $(1+\lambda^2)\ue^{-a\lambda/2}\le1+16/(\ue^2a^2)\le4a^{-2}$, so
  $1+\lambda^2\le4a^{-2}\ue^{a\lambda}\ue^{-a\lambda/2}\le4a^{-2}\ue^{a\lambda}\ue^{-a/(4\varepsilon)}$.
  \end{itemize}
\item \emph{The splitting.} $\Delta_\varepsilon(s)-\Delta(s)=-J_\varepsilon\Delta(s)(1-J_\varepsilon)-(1-J_\varepsilon)\Delta(s)$.
\item \emph{First term.} With $z_1:=(1-J_\varepsilon)z$, Lemma~\ref{lem:ell} and $\norm{J_\varepsilon}\le1$ give
$\norm{J_\varepsilon\Delta(s)z_1}\le2(\norm{\Delta z_1}+\norm{z_1})\le4\norm{(1+\Lambda^2)(1-j_\varepsilon(\Lambda))z}\le20\,a^{-2}\ue^{-a/(4\varepsilon)}\norm{\ue^{a\Lambda}z}$,
by the symbol bound.
\item \emph{Second term.} By Lemma~\ref{lem:wc} with $\tau_0=1$ (words of length $1$ and $2$, radii $a/2<a$),
$\norm{\ue^{(a/2)\Lambda}\Delta(s)z}\le3(1+4C_1/a)^2\norm{\ue^{a\Lambda}z}$. By the symbol bound with $a/2$ in place of $a$,
$\norm{(1-J_\varepsilon)w}\le20\,a^{-2}\ue^{-a/(8\varepsilon)}\norm{\ue^{(a/2)\Lambda}w}$. Take $w=\Delta(s)z$.
\end{itemize}
Adding the two terms gives the claim with $C_0=20+60(1+4C_1)^2$. \qed
\end{proof}

\subsection{Regularity and the base point}\label{subsec:dd-base}

The theorems are stated for amplitudes of finite norm and for every boundary point $b$. The next two lemmas reduce them to
$C^\infty$ amplitudes, the base point $i$ and $b=\infty$.

\begin{lemma}[regularity]\label{lem:lift}
Let $0<\alpha\le1$ and $F\in\operatorname{Dom}(\ue^{\alpha\Lambda})$. Then $F$ has a unique $C^\infty$ representative, and every word in
$X,Y$ applied to it lies in $L^2(\Hb)$. Consequently every $a$ with $\norm{\ue^{\alpha\Lambda}X_0^\beta a}<\infty$, and every
$a\in\operatorname{Dom}(\ue^{\alpha\Lambda})$, has a unique $C^\infty$ representative. The approximations $u^{(K)}$ are built from this
representative.
\end{lemma}

\begin{proof}
\begin{itemize}
\item By Lemma~\ref{lem:wc} at $\alpha''=0$ (with $\tau_0=1$), every word applied to $F$ in the sense of distributions lies in $L^2$. The same holds
for $F\circ p$ for every affine isometry $p(z)=p_1z+p_2$, $p_1>0$: these commute with $X$, $Y$ and $\Lambda$, so
$F\circ p\in\operatorname{Dom}(\ue^{\alpha\Lambda})$.
\item Let $w\in L^2$ be such that all words applied to $w$ lie in $L^2$, and let $c\in C_c^\infty(\Hb)$ be real. With the Euclidean
Laplacian $\Delta_{\mathrm E}=\partial_x^2+\partial_y^2=y^{-2}(X^2+Y^2-Y)$,
\[
\Delta_{\mathrm E}(cw)=(\Delta_{\mathrm E}c)\,w+2y^{-2}\bigl((Xc)(Xw)+(Yc)(Yw)\bigr)+c\,y^{-2}(X^2+Y^2-Y)w,
\]
a finite sum of terms $c'w'$ of the same kind. So $\Delta_{\mathrm E}^n(cw)\in L^2(\mathbb R^2)$ for every $n$; on $\operatorname{supp}c$
the measures $dx\,dy$ and $d\nu$ are comparable.
\item Hence $|\xi|^{2n}\widehat{cw}\in L^2(\mathbb R^2)$ for every $n$, all moments of $\widehat{cw}$ are integrable, and $cw$ has a
$C^\infty$ representative.
\item Take $c=1$ near $i$. This gives a $C^\infty$ representative of $F\circ p$ near $i$ for every $p$, that is, of $F$ near
every point. Two continuous representatives on an open set coincide, so they glue to a $C^\infty$ representative of $F$, which is unique.
\item For the consequences take $F=X_0^\beta a$, respectively $F=a$; $X_0^{-\beta}$ is smooth and positive. \qed
\end{itemize}
\end{proof}

\begin{lemma}[base point]\label{lem:base}
Let $x_0\in\Hb$ and $b\in\partial\Hb$.
\begin{itemize}
\item[\textup{(i)}] There is $g\in SL(2,\mathbb R)$, unique up to sign, with $g\,x_0=x_0$ and $g\,\infty=b$. We normalize
$\varphi_b:=\varphi_\infty\circ g^{-1}$. The construction of \cite[(1.20)--(1.29)]{S} is covariant under isometries: the initial state and
the approximations for $b$ are $v\mapsto v\circ g^{-1}$ applied to those for $\infty$ with the amplitude $a\circ g$.
\item[\textup{(ii)}] Let $p(z)=(\operatorname{Im}x_0)\,z+\operatorname{Re}x_0$ and $q:=g\,p$, so $q\,i=x_0$ and $q\,\infty=b$. For an amplitude $a$ put
$a_q:=a\circ q$ and $\Gamma_q:=q^{-1}\Gamma q$. Then, with $c:=\ue^{-i\log(\operatorname{Im}x_0)/\hbar}$,
\[
(a\,\ue^{i\varphi_b/\hbar})\circ q=c\,a_q\,\ue^{i\varphi_\infty/\hbar},\qquad u^{(K)}_b(t)[a]\circ q=c\,u^{(K)}_\infty(t)[a_q],
\]
where the subscript names the boundary point of the construction.
\item[\textup{(iii)}] The map $f\mapsto f\circ q$ is unitary from $L^2(\Hb/\Gamma)$ to $L^2(\Hb/\Gamma_q)$, intertwines the periodizations over $\Gamma$ and
$\Gamma_q$, and commutes with $\Delta$, $\cU(t)$ and $\ue^{\alpha\Lambda}$. The swapped and printed norms of $a$ with base point $x_0$ equal those of $a_q$
with base point $i$. For every $s$, $\sup_z\sum_{\gamma\in\Gamma_q}\ue^{-s\,d(i,\gamma z)}=\sup_z\sum_{\gamma\in\Gamma}\ue^{-s\,d(x_0,\gamma z)}$.
\end{itemize}
Consequently the error norm of $(\Gamma,x_0,b,a)$ equals that of $(\Gamma_q,i,\infty,a_q)$, and an estimate proved for base point $i$ and
$b=\infty$, with constants depending on the group only through quantities invariant under $\Gamma\mapsto\Gamma_q$, $x_0\mapsto i$, holds
for every $x_0$ and $b$.
\end{lemma}

\begin{proof}
\begin{itemize}
\item (i) The stabilizer of $x_0$ in $PSL(2,\mathbb R)$ acts simply transitively on $\partial\Hb$. The phase $\varphi_b$, the transport and
the operators $\Delta(s)$ for $b$ are defined from the geodesics ending at $b$, so they are the transports by $g$ of those for $\infty$.
\item (ii) $\varphi_b\circ q=\varphi_\infty\circ p=\varphi_\infty-\log\operatorname{Im}x_0$, since $\operatorname{Im}p(z)=(\operatorname{Im}x_0)\operatorname{Im}z$.
By (i), composing with $g$ turns the objects for $b$ and $a$ into those for $\infty$ and $a\circ g$. The affine map $p$ commutes with
$X$, $Y$ and $S(t)$, hence with $\Delta(s)$ and the $P_k(t)$, and it changes $\ue^{i\varphi_\infty/\hbar}$ by the constant factor $c$. This
gives both identities.
\item (iii) $q$ is an isometry of $\Hb$ with $q\,i=x_0$, so $d(x_0,qz)=d(i,z)$ and the weights $X_0$, $\langle d\rangle$ based at $x_0$ pull
back to those based at $i$. The spectral functions of $-\Delta$ commute with composition by isometries. For the last identity,
$d(i,q^{-1}\gamma qz)=d(x_0,\gamma(qz))$. \qed
\end{itemize}
\end{proof}

\begin{remark}
The hypothesis $g\,x_0=x_0$ is used in (iii): the Poincar\'e sum of $\Gamma$ is not bounded uniformly in the base point when $\Gamma$
has parabolic elements, so the base point of the weight must be the point at which the sum is taken.
\end{remark}

\subsection{Approximation with the power weight \texorpdfstring{$X_0^\beta$}{X0}}\label{subsec:dd-thm}

\begin{theorem}[approximation with the power weight $X_0^\beta$]\label{thm:dd-appx}
Let $\Gamma$ be any Fuchsian group, $\beta>1/2$ and $\kappa>0$, and $\alpha_0\in(0,\min(1,\pi/4))$. Let $\chi$ be the cutoff of \cite[(4.20)]{S} with the standing hypotheses of Appendix~\ref{app:kl} (real, even, $0\le\chi\le1$).
There are $C,\delta>0$, depending only on $\beta,\kappa,\alpha_0,\chi,\Gamma$, such that the following holds for all $0<\hbar\le1$, all $\alpha\in(0,\alpha_0]$ with $\alpha^3\ge\kappa\hbar$, all $K,N\in\mathbb N_0$, every $b\in\partial\Hb$, every $a$ with $\norm a=\norm{\ue^{\alpha\Lambda}X_0^\beta a}<\infty$, and all $0\le t\le\delta(\alpha/\hbar)^{1/2}$:
\[
\norm{\Delta^N\bigl[\cU(t)(a\,\ue^{i\varphi_b/\hbar})_\Gamma-u^{(K)}(t)_\Gamma\bigr]}_{L^2(M)}\le C^{N+K+1}\,N^{2N}\,K!\,(\alpha\hbar)^{-2N}\Bigl(\frac{\hbar t}{\alpha^2}\Bigr)^{K+1}\norm a,
\]
with $N^{2N}:=1$ at $N=0$.
\end{theorem}

\begin{remark}
This is the power-weight form of the swapped-class estimate. Lemma~\ref{lem:m2w} transfers it to the norm of
Theorem~\ref{thm:dd14}; the deduction follows the proof below. Section~\ref{sec:changes} compares the resulting bounds with \cite{S}.
\end{remark}

\begin{proof}
\medskip\noindent\textbf{Reductions.} By Lemma~\ref{lem:base} we may take $x_0=i$ and $b=\infty$: the group enters the estimates below only
through the Poincar\'e series bound at the base point. By Lemma~\ref{lem:lift}, $a$ is $C^\infty$. We write $\varphi=\varphi_\infty$.

\medskip\noindent\textbf{Parameters.} All lemmas are used with $\tau_0=1$. Let $c_M$ be the constant of Lemma~\ref{lem:mw} and $c_A$ that of
Proposition~\ref{prop:a0}. Put
\[
a_0:=c_M\alpha,\qquad a_1:=\tfrac34a_0,\qquad a_2:=\tfrac12a_0,\qquad\varepsilon:=(\alpha\hbar)^{1/2}/2,\qquad x:=(\alpha/\hbar)^{1/2},
\]
so that $\varepsilon\le\tfrac12$, $\varepsilon/\hbar=x/2$ and $\alpha/\varepsilon=2x$. Fix $\gamma:=(2\beta+1)/4\in(1/2,\beta)$ and put $\theta:=\gamma/\beta\in(0,1)$.

Use (D4) with $\omega=W_\gamma$; every word applied to $a$ lies in $L^2_\omega$ by Lemma~\ref{lem:sc}, since $W_\gamma\le X_0^\beta$. For each $\mathbf s=(s_1,\dots,s_{K+1})\in\Sigma_{K+1}(t)$ let
\[
g(t'):=V_\varepsilon(t',s_1)\,b',\quad0\le t'\le s_1,\qquad w:=g(0)=V_\varepsilon(s_1)^*b',
\]
and write
\[
\mathcal V_{s_1}b'=\mathcal V_0w+\bigl(\mathcal V_{s_1}g(s_1)-\mathcal V_0g(0)\bigr).
\]

\medskip\noindent\textbf{(I) Mollified part.}
\begin{itemize}
\item Lemma~\ref{lem:t2} at $t=0$, and Proposition~\ref{prop:ws} at the weight $X_0^\gamma$ (with $\hbar s_1\le\varepsilon$), give
\[
\norm{\mathcal V_0w}=\norm{[e^{i\varphi/\hbar}V_\varepsilon(s_1)^*b']_\Gamma}\le C\norm{W_\gamma V_\varepsilon(0,s_1)b'}\le C\ue^{G_\gamma/2}\norm{W_\gamma b'}\le C\norm{X_0^\beta b'},
\]
using $W_\gamma=X_0^\gamma\le X_0^\beta$ (since $X_0\ge1$ and $\gamma\le\beta$).
\item By (WB-w) with $n=0$:
\[
\text{(I)}\le C\,C_4^{N+K+1}(2N)!\,(2K+2)!\,\alpha^{-(2K+2)}(\alpha\hbar)^{-2N}\norm a.
\]
\end{itemize}

\medskip\noindent\textbf{(II) Mollifier part.}
\begin{itemize}
\item By (WB-w) with $n\le2$, $X_0^\beta Zb'\in L^2$ for $Z\in\mathcal S=\{1,X,Y,X^2,Y^2\}$. So Lemma~\ref{lem:path}(c), with the exponent
$\gamma$, shows that $g$ satisfies the hypotheses of Lemma~\ref{lem:surf-duh} for $\omega=W_\gamma$. Since $g'=c\,\Delta_\varepsilon(t')g$,
Duhamel on $M$ gives $\mathcal V_{s_1}g(s_1)-\mathcal V_0g(0)=-c\int_0^{s_1}\mathcal V_{t'}d(t')\,dt'$, with
\[
d(t'):=(\Delta(t')-\Delta_\varepsilon(t'))\,g(t').
\]
\item By the unitarity of $\cU(t')$ on $M$ and Lemma~\ref{lem:t2} (exponent $\gamma$),
\[\norm{\mathcal V_{t'}d(t')}=\norm{[e^{i\varphi/\hbar}S(t')d(t')]_\Gamma}\le C\ue^{\gamma t'}\norm{W_\gamma d(t')}.\]
\item H\"older on $L^2$: $\norm{W_\gamma d}\le\norm{W_\beta d}^\theta\norm{d}^{1-\theta}$, since $W_\gamma=W_\beta^\theta$.
\item Weighted factor: by Lemma~\ref{lem:path}(b) with the exponent $\beta$, and (WB-w) with $n\le2$ (five words; the bound increases with $n$),
\begin{align*}
\norm{W_\beta d(t')}&\le C\varepsilon^{-3}\sum_{Z\in\mathcal S}\norm{X_0^\beta Zb'}\\
&\le C\varepsilon^{-3}\,C_4^{N+K+3}(2N)!\,(2K+4)!\,\alpha^{-(2K+4)}(\alpha\hbar)^{-2N}\norm a.
\end{align*}
\item Unweighted factor.
\begin{itemize}
\item Lemma~\ref{lem:mw} gives $\norm{\ue^{a_0\Lambda}a}\le C_*B_\beta\norm a$.
\item (WB-u), applied with $a_0$ in the role of $\alpha$ and the radii $a_1<a_0$ (gap $a_0/4$), gives
$\norm{\ue^{a_1\Lambda}b'}\le C_4^{N+K+1}(2N)!(2K+2)!\,a_0^{-(2K+2)}(a_0\hbar)^{-2N}\norm{\ue^{a_0\Lambda}a}$.
\item Proposition~\ref{prop:a0}, from the radius $a_1$ to $a_2$ (so $(a_1-a_2)/a_1=\tfrac13$) and for $|t'-s_1|\le t\le\tfrac13c_A\varepsilon/\hbar$, gives
$\norm{\ue^{a_2\Lambda}g(t')}\le\norm{\ue^{a_1\Lambda}b'}$.
\item Lemma~\ref{lem:unw} with $a=a_2$ gives $\norm{d(t')}\le C_0a_2^{-4}\ue^{-a_2/(8\varepsilon)}\norm{\ue^{a_2\Lambda}g(t')}$, and $a_2/(8\varepsilon)=c_Mx/8$.
\end{itemize}
\item Combining: put $F:=(C_4/c_M^2)^{N+K+3}(2N)!\,(2K+4)!\,\alpha^{-(2K+4)}(\alpha\hbar)^{-2N}\norm a$. Since $c_M\le1\le C_4$ and $\alpha\le1$, the weighted
factor is at most $C\varepsilon^{-3}F$ and the unweighted one at most $C\alpha^{-4}\ue^{-c_Mx/8}F$. After the H\"older split, using $\varepsilon,\alpha\le1$,
\[
\text{(II)}\le C\,F\,R,\qquad R:=\hbar t\,\ue^{\gamma t}\,\varepsilon^{-3}\,\alpha^{-4}\,\ue^{-(1-\theta)c_Mx/8}.
\]
\item For $t\le\delta x$ with $\delta\le(1-\theta)c_M/(16\gamma)$: $\ue^{\gamma t}\le\ue^{(1-\theta)c_Mx/16}$. Since $\alpha^3\ge\kappa\hbar$ and $\alpha\le1$,
\[
\alpha^{-1}\le x/\sqrt\kappa,\qquad\varepsilon^{-1}=2x/\alpha\le2x^2/\sqrt\kappa,\qquad\hbar t\le\delta\hbar x=\delta\alpha/x\le\delta/x .
\]
Hence $R\le C(\kappa)\,x^{9}\,\ue^{-(1-\theta)c_Mx/16}\le C(\kappa)\,x^{-2}$, and $x^{-2}=\hbar/\alpha\le\alpha^2/\kappa$. So $R\le C(\kappa)\alpha^2$, and
(II) is at most $C\alpha^2F$.
\end{itemize}

\medskip\noindent\textbf{Collecting.}
\begin{itemize}
\item $\int_{\Sigma_{K+1}(t)}ds=t^{K+1}/(K+1)!$.
\item $(2K+4)!/(K+1)!\le C^{K+1}K!$ and $(2K+2)!/(K+1)!\le C^{K+1}K!$, since $(2K+4)!\le4^{K+2}((K+2)!)^2$.
\item $(2N)!\le4^NN^{2N}$.
\item $\hbar^{K+1}t^{K+1}\alpha^{-(2K+2)}=(\hbar t/\alpha^2)^{K+1}$, and $\alpha^2F$ carries $\alpha^{-(2K+2)}$.
\end{itemize}
Hence
\[
\norm E\le C^{N+K+1}N^{2N}K!\,(\alpha\hbar)^{-2N}\Bigl(\frac{\hbar t}{\alpha^2}\Bigr)^{K+1}\norm a.
\]
\qed

\medskip\noindent\textbf{Time windows used.}
\begin{itemize}
\item Weighted propagation and the path bound: $\hbar s_1\le\varepsilon$, that is $t\le x/2$.
\item $A_0$: $t\le\tfrac13c_A\varepsilon/\hbar=c_Ax/6$.
\item Transport: $\ue^{\gamma t}$ is absorbed for $\delta\le(1-\theta)c_M/(16\gamma)$.
\end{itemize}
So $\delta:=\min\bigl(\tfrac12,\;c_A/6,\;(1-\theta)c_M/(16\gamma)\bigr)$ works. The factor $\kappa^{1/2}/\max(1,C)$ is used only in Corollary~\ref{cor:dd15-appx}, to ensure $\delta_K\le1$.

\medskip\noindent\textbf{Dependence on $\Gamma$.} The group enters only through Lemma~\ref{lem:t2}, used in (I) and (II) with exponent $\gamma$, whose
constant $C(\Gamma,\gamma)=2^\gamma P_\Gamma(2\gamma)^{1/2}$ depends on $\Gamma$ only through the Poincar\'e series bound $P_\Gamma(2\gamma)$, $2\gamma=\beta+\tfrac12$.
All other inputs are estimates on $\Hb$. So $\delta$ does not depend on $\Gamma$, and $C$ depends on $\Gamma$ only through $P_\Gamma(\beta+\tfrac12)$;
this is the uniformity claimed in Theorem~\ref{thm:dd14}. (It concerns this theorem only: the extra cap $\delta\le\kappa^{1/2}/\max(1,C)$ of
Corollary~\ref{cor:dd15-appx} depends on $\Gamma$ through $C$.)
\end{proof}

\begin{proof}[Proof of Theorem~\ref{thm:dd14}]
Let $c_0$ and $C'$ be the constants of Lemma~\ref{lem:m2w}, and let $\norm a$ denote the norm of Section~\ref{sec:swapped}, with the weight
$\ue^{\beta\langle d\rangle}$. If $\alpha\le\alpha_0$ and $\alpha^3\ge\kappa\hbar$, then $\bar\alpha:=c_0\alpha\le c_0\alpha_0$,
$\bar\alpha^3\ge c_0^3\kappa\hbar$ and $\norm{\ue^{\bar\alpha\Lambda}X_0^\beta a}\le C'\norm a$. Apply Theorem~\ref{thm:dd-appx} with
$c_0\alpha_0$, $c_0^3\kappa$ and $\bar\alpha$ in place of $\alpha_0$, $\kappa$ and $\alpha$, and let $C_{\rm int},\delta_{\rm int}$ be its constants. Its
window $t\le\delta_{\rm int}(\bar\alpha/\hbar)^{1/2}$ is the window $t\le\delta(\alpha/\hbar)^{1/2}$ with $\delta:=c_0^{1/2}\delta_{\rm int}$, and
$(\bar\alpha\hbar)^{-2N}(\hbar t/\bar\alpha^2)^{K+1}=c_0^{-2N-2K-2}(\alpha\hbar)^{-2N}(\hbar t/\alpha^2)^{K+1}$. So the bound holds with
$C:=C_{\rm int}\,c_0^{-2}\max(1,C')$. The dependence on $\Gamma$ is that of Theorem~\ref{thm:dd-appx}.
\end{proof}

\subsection{Optimal truncation and pointwise estimates}\label{subsec:dd-cor}

\begin{corollary}[optimal truncation with the power weight $X_0^\beta$]\label{cor:dd15-appx}
Under the hypotheses of Theorem~\ref{thm:dd-appx}, with $\delta\le\kappa^{1/2}/\max(1,C)$, let $C$ be its constant, put $\delta_K:=C\hbar t/\alpha^2$, and let $\varepsilon\in(0,1)$. Then for every integer $K$ with $(1-\varepsilon)/\delta_K\le K\le(1-\varepsilon/2)/\delta_K$,
\[
\norm{\Delta^N\bigl[\cU(t)(a\,\ue^{i\varphi_b/\hbar})_\Gamma-u^{(K)}(t)_\Gamma\bigr]}_{L^2(M)}\le C'\varepsilon^{-1/2}C^{N+1}N^{2N}(\alpha\hbar)^{-2N}\ue^{-(1-\varepsilon)\alpha^2/(C\hbar t)}\norm a,
\]
with $C'$ an absolute constant.
\end{corollary}

\begin{proof}[Proof of Corollary~\ref{cor:dd15-appx} (Schubert's computation for his Corollary 1.5 \cite{S})]
\begin{itemize}
\item The bound of Theorem~\ref{thm:dd-appx} is $C^NN^{2N}(\alpha\hbar)^{-2N}\,\delta_K^{K+1}K!\,\norm a$, because $C^{K+1}(\hbar t/\alpha^2)^{K+1}=\delta_K^{K+1}$.
On the window, $K\ge(1-\varepsilon)/\delta_K>0$, so $K\ge1$.
\item \[\delta_K^{K+1}K!\le\delta_K\cdot e\,K^{1/2}\ue^{-K}(K\delta_K)^K\] (Stirling).
\item $K\delta_K\le1-\varepsilon/2$ gives $(K\delta_K)^K\le\ue^{-K\varepsilon/2}$, and $\sup_KK^{1/2}\ue^{-K\varepsilon/2}\le(e\varepsilon)^{-1/2}$.
\item $K\ge(1-\varepsilon)/\delta_K$ gives $\ue^{-K}\le\ue^{-(1-\varepsilon)\alpha^2/(C\hbar t)}$.
\item $\hbar t/\alpha^2\le\delta\,\hbar^{1/2}\alpha^{-3/2}\le\delta\kappa^{-1/2}\le1/\max(1,C)$, by $t\le\delta(\alpha/\hbar)^{1/2}$, $\alpha^3\ge\kappa\hbar$ and
$\delta\le\kappa^{1/2}/\max(1,C)$. Hence $\delta_K\le1$ and $C^N\delta_K=C^{N+1}\hbar t/\alpha^2\le C^{N+1}$, and the bound holds with $C'=\ue^{1/2}$. \qed
\end{itemize}
\end{proof}

\begin{remark}
In \cite[Corollary~1.5]{S} the window is printed as $K\in[C\alpha^2/\hbar t-\varepsilon,\;C\alpha^2/\hbar t-\varepsilon/2]$, with the constant $C$ of
\cite[(1.30)]{S}. The proof there uses $K\delta_K\in[1-\varepsilon,1-\varepsilon/2]$, which is the window used here. If the window contains no integer
the statement is empty; at $t=0$ the error is $0$.
\end{remark}

\begin{corollary}[pointwise approximation with the power weight $X_0^\beta$]\label{cor:dd16-appx}
Assume, in addition, that $M$ has injectivity radius $\ge r_0>0$. Then, under the hypotheses of Theorem~\ref{thm:dd-appx}, the error
$\cU(t)(a\,\ue^{i\varphi_b/\hbar})_\Gamma-u^{(K)}(t)_\Gamma$ lies in $\operatorname{Dom}(\Delta_M^N)$, its image under $\Delta^N$ has a continuous
representative, and at every point of $M$
\[
\Bigl|\Delta^N\bigl[\cU(t)(a\,\ue^{i\varphi_b/\hbar})_\Gamma-u^{(K)}(t)_\Gamma\bigr]\Bigr|\le C^{N+K+1}(2N)!\,K!\,(\alpha\hbar)^{-2N-4}\Bigl(\frac{\hbar t}{\alpha^2}\Bigr)^{K+1}\norm a,
\]
with $C$ depending also on $r_0$. In the same way Corollary~\ref{cor:dd15-appx} gives the pointwise exponential bound with $(\alpha\hbar)^{-2N-4}$.
\end{corollary}

\begin{proof}
\begin{itemize}
\item If $d(\tilde z,\gamma\tilde z)\ge2r_0$ for all $\gamma\ne1$, the balls $B(\gamma\tilde z,r_0)$ are disjoint, and those with
$d(\tilde z,\gamma\tilde z)<2$ lie in $B(\tilde z,2+r_0)$. So the multiplicity of Lemma~\ref{pc:g43} satisfies $m(z)\le C(r_0)$, and that lemma gives, for
$u\in\operatorname{Dom}(\Delta_M)$, a continuous representative with
$\sup_M|u|\le C(r_0)\bigl(\norm{\Delta u}_{L^2(M)}+\norm u_{L^2(M)}\bigr)$.
\item Apply this to $u=\Delta^NE$, with Theorem~\ref{thm:dd-appx} at $N$ and at $N+1$. Use $N^{2N}\le(\ue/2)^{2N}(2N)!$,
$(N+1)^{2N+2}\le54^{N+1}(2N)!$ and $\alpha\hbar\le1$. \qed
\end{itemize}
\end{proof}

\begin{remark}
The proof uses one Laplacian and gives the exponent $-2N-2$ of $\alpha\hbar$. We state the corollary with $-2N-4$, the exponent that
enters the comparison with \cite[(1.34)]{S} in Section~\ref{sec:changes}.
\end{remark}

\begin{remark}[Cusped $M$]
The Sobolev inequality \cite[(5.91)]{S} fails on $M$ with cusps (Remark~\ref{rem:g4}). On such $M$ the same proof, with Lemma~\ref{pc:g43} as it stands, gives the pointwise bound with the additional factor $C\,m(z)^{1/2}$, as in Corollary~\ref{pc:cor16}.
\end{remark}

% The corollary body ends before the remarks.

\subsection{Lemma EX: the example amplitudes of \texorpdfstring{\cite{S}}{[S]} are covered}\label{subsec:dd-ex}

\begin{lemma}[EX]\label{lem:ex}
Fix $r_1<\tfrac12$. There are constants $C_H$ and $C_{\mathrm E}$, depending only on $r_1$, such that for $\mu\ge1$ and
$0<\alpha\le\min(r_1,\mu^{-1/2})$,
\[
\norm{\ue^{\alpha\Lambda}\ue^{-\mu\langle d\rangle}}_{L^2(\Hb)}\le(\pi\ue)^{1/2}\,\ue^{C_H\mu\alpha^2}\norm{\ue^{-\mu\langle d\rangle}}_{L^2(\Hb)}\le C_{\mathrm E}\,\mu^{-1/2}.
\]
\end{lemma}

\begin{proof}
We use the notation of Section~\ref{subsec:gutzmer}, with $x_0=i$.

\emph{The function $H$.} Let $G(z):=-\arccos(z)^2$ and $H(z):=(1+G(z))^{1/2}-1$. The function $G$ is holomorphic on
$\mathbb C\setminus(-\infty,-1]$ and equals $d^2$ at $z=\cosh d$, so $\langle d\rangle=H(X_0)$ on $\Hb$. The number $1+G(z)$ is real and at most $0$ only for
real $z\le\cos1$, so $H$ is holomorphic on the complement of $(-\infty,\cos1]$, real on $(\cos1,\infty)$ and nondecreasing there.
For $0<\rho<\tfrac12$ let $S_\rho:=\{\operatorname{Re}z\ge\cos2\rho,\ |\arg z|\le2\rho\}$, and fix $r_1<r_1'<r_2<\tfrac12$ depending only on $r_1$.
Since $\cos2r_2>\cos1$, $H$ is holomorphic near $S_{r_2}$. As $|z|\to\infty$ in $S_{r_2}$ we have $G=L^2$ with
$L=\log(z+(z^2-1)^{1/2})=\log(2z)+O(|z|^{-2})$, hence $H=L-1+O(1/|L|)$; so $M_H:=\sup_{S_{r_2}}|H(z)-\log z|<\infty$. There is $c>0$,
depending only on $r_1',r_2$, such that the disc of radius $c|z|$ about any $z\in S_{r_1'}$ lies in $S_{r_2}$, and Cauchy's estimates for
$H-\log$ on these discs give
\[
|H'(z)|\le C_1/|z|,\qquad|H''(z)|\le C_1/|z|^2\qquad(z\in S_{r_1'}).
\]

\emph{The orbit norms.} By Lemma~\ref{lem:gutzmer}(b), on the orbit $G\exp(i\rho h_0)x_0$, $\rho\le r_1'$, the extension of $X_0$ takes the
values $\tilde X_0=cA+isB$ with $c=\cos2\rho$, $s=\sin2\rho$, $A=X_0(gx_0)\ge1$ and $|B|\le A$; these lie in $S_\rho\subset S_{r_1'}$, and so does
the vertical segment from $cA$ to $\tilde X_0$. So $f:=\ue^{-\mu\langle d\rangle}=\ue^{-\mu H(X_0)}$ extends holomorphically to $\Xi_{r_1'}$. Since $H'(cA)$ is
real, Taylor's formula on that segment gives $\operatorname{Re}H(\tilde X_0)\ge H(cA)-\tfrac12\sup|H''|\,s^2B^2\ge H(cA)-C_1s^2/(2c^2)$, and
$H(A)-H(cA)\le(1-c)A\sup_{[cA,A]}|H'|\le C_1(1-c)/c$. With $1-c\le2\rho^2$, $s^2\le4\rho^2$ and $c\ge\cos1>\tfrac12$,
\[
\operatorname{Re}H(\tilde X_0)\ge\langle d\rangle(gx_0)-C_H\rho^2,\qquad C_H:=12\,C_1 .
\]
Hence $|f(g\exp(i\rho h_0)x_0)|\le\ue^{C_H\mu\rho^2}f(gx_0)$ and $N_\rho(f)\le\ue^{C_H\mu\rho^2}\norm f$ for $\rho\le r_1'$. In particular
$\sup_{\rho<r_1'}N_\rho(f)<\infty$, and Lemma~\ref{lem:gutzmer}(a), with $R=r_1'$ and $r=\alpha\le r_1$, gives the first inequality.

\emph{The $L^2$ norm.} In geodesic polar coordinates $\norm f^2=2\pi\int_0^\infty\ue^{-2\mu\langle\rho\rangle}\sinh\rho\,d\rho$, with
$\langle\rho\rangle=(1+\rho^2)^{1/2}-1$. For $\rho\le1$, $\langle\rho\rangle\ge\rho^2/(1+\sqrt2)$ and $\sinh\rho\le2\rho$; for $\rho\ge1$,
$\langle\rho\rangle\ge\rho-1$ and $\sinh\rho\le\ue^\rho$. So for $\mu\ge1$
\[
\norm f^2\le2\pi\Bigl(\int_0^\infty2\rho\,\ue^{-2\mu\rho^2/(1+\sqrt2)}d\rho+\ue\int_0^\infty\ue^{-(2\mu-1)u}du\Bigr)\le\frac{2\pi(2+\ue)}{\mu},
\]
and the second inequality holds with $C_{\mathrm E}:=(\pi\ue)^{1/2}\ue^{C_H}(2\pi(2+\ue))^{1/2}$, because $\mu\alpha^2\le1$.
\end{proof}

\begin{remark}[Consequence for the example of \texorpdfstring{\cite{S}}{[S]}]
\cite[Remark~1.7, (1.38)]{S} gives the example $a_\hbar:=\hbar^{-\delta/2}\ue^{-\hbar^{-\delta}\langle d\rangle}$, said there to lie in $H_{\alpha,\beta}$ with $\beta=1$ and $\alpha=\hbar^\delta$, and to give long times for $\delta<1/3$. In the swapped $\langle d\rangle$-class, $\norm{a_\hbar}^{\rm sw}_{\alpha,\beta}=\hbar^{-\delta/2}\norm{\ue^{\alpha\Lambda}\ue^{-(\hbar^{-\delta}-\beta)\langle d\rangle}}$. Apply Lemma~\ref{lem:ex} with $\mu=\hbar^{-\delta}-\beta$ and $\alpha_\hbar:=(\hbar^{-\delta}-\beta)^{-1/2}\asymp\hbar^{\delta/2}$:
\begin{itemize}
\item $\norm{a_\hbar}^{\rm sw}_{\alpha_\hbar,\beta}\le C$ uniformly in $\hbar$, for every fixed $\beta$ (not only $\beta=1$);
\item the swapped $\langle d\rangle$-class transfers to the $\tilde\psi$-class by Lemma~\ref{lem:m2w}, at the radius $c_0\alpha_\hbar$ with the fixed ratio $c_0$;
\item so Theorem~\ref{thm:dd-appx} (any $\beta>1/2$, every Fuchsian $\Gamma$) applies when $(c_0\alpha_\hbar)^3\ge\kappa\hbar$, i.e. for $0<\delta<2/3$, and for $\hbar$ small enough that $\mu=\hbar^{-\delta}-\beta\ge1$ and $\alpha_\hbar\le r_1$. The case $\delta=0$ is not covered: $\ue^{-\langle d\rangle}$ has infinite $\ue^{\beta\langle d\rangle}$-weighted $L^2$ norm for every $\beta>1/2$. The time window is $t\le\delta'(\alpha_\hbar/\hbar)^{1/2}\asymp\hbar^{\delta/4-1/2}$.
\end{itemize}
The radius $\alpha=\hbar^\delta$ of \cite{S} gave $\delta<1/3$ and $t\ll\hbar^{-(1-3\delta)/2}$; the holomorphic tube bound gives the larger radius $\hbar^{\delta/2}$. The example is therefore covered. The printed (original-order) class is treated in Section~\ref{sec:printed}; the amplitudes named in \cite{S} are covered already by the swapped class.
\end{remark}

\section{Periodization by Cauchy--Schwarz}\label{sec:t2}

\begin{lemma}[P1]\label{lem:p1}
For every Fuchsian $\Gamma$ and $s>1$: $\sup_{z}\sum_{\gamma}\ue^{-s\,d(x_0,\gamma z)}<\infty$.
\end{lemma}
\begin{proof}
The orbit $\Gamma x_0$ is uniformly separated, and its stabilizer is finite. Packing balls of a fixed radius
around orbit points gives $\#\{\gamma:d(z,\gamma x_0)<n+1\}\le C\ue^{n}$, uniformly in $z$.
\end{proof}

\begin{lemma}[T2]\label{lem:t2}
For every Fuchsian $\Gamma$, every $s>1/2$, $t\ge0$, every real phase $\varphi$ and every measurable $w$ with $\ue^{s\tilde\psi}w\in L^2(\Hb)$:
\[
\norm{\bigl[(S(t)w)\,\ue^{i\varphi/\hbar}\bigr]_\Gamma}_{L^2(M)}\le C(\Gamma,s)\,\ue^{st}\,\norm{\ue^{s\tilde\psi}w}_{L^2(\Hb)},
\]
with $C(\Gamma,s)=2^sP_\Gamma(2s)^{1/2}$, where $P_\Gamma(\sigma):=\sup_z\sum_{\gamma\in\Gamma}\ue^{-\sigma d(x_0,\gamma z)}$ is the Poincar\'e series bound of Lemma~\ref{lem:p1}.
\end{lemma}

\begin{proof}
\begin{itemize}
\item \emph{Pointwise Cauchy--Schwarz.} For $f$ on $\Hb$ and $z\in\Hb$:
\[
\Bigl|\sum_\gamma f(\gamma z)\Bigr|^2\le\Bigl(\sum_\gamma\ue^{-2sd(\gamma z,x_0)}\Bigr)\Bigl(\sum_\gamma\ue^{2sd(\gamma z,x_0)}|f(\gamma z)|^2\Bigr).
\]
The first factor is $\le P_\Gamma(2s)$ for every $z$, by Lemma~\ref{lem:p1} with exponent $2s>1$ (since $d(\gamma z,x_0)=d(z,\gamma^{-1}x_0)$).
\item \emph{Unfolding.} Integrate over a fundamental domain (Lemma 3.1 of \cite{S}, Tonelli). This gives
\[
\norm{[f]_\Gamma}^2_{L^2(M)}\le P_\Gamma(2s)\int_\Hb\ue^{2sd(z,x_0)}|f(z)|^2\,d\nu(z).
\]
Prove this first for finite partial sums, then pass to the limit by monotone convergence. In particular $\sum_\gamma|f\circ\gamma|$ converges a.e. and in $L^2(M)$. Torsion-freeness is not needed: the fixed sets of elliptic elements are null, so unfolding over a fundamental domain carries no stabilizer factor, and any multiplicity is already counted in the orbit series $P_\Gamma(2s)$.
\item \emph{Transport.} Take $f=(S(t)w)\ue^{i\varphi/\hbar}$, so $|f|=|S(t)w|$.
  \begin{itemize}
  \item $(S(t)w)(z)=\ue^{-t/2}w(z_t)$, where $z_t$ is at distance $t$ from $z$, so $d(z,x_0)\le d(z_t,x_0)+t$.
  \item $S(t)$ is unitary, so $\int\ue^{2sd(z)}|S(t)w(z)|^2\,d\nu(z)\le\ue^{2st}\int\ue^{2sd(z')}|w(z')|^2\,d\nu(z')$.
  \end{itemize}
\item Finally $\ue^{sd}\le2^s\ue^{s\tilde\psi}$, since $\tilde\psi=\log\cosh d\ge d-\log2$. \qed
\end{itemize}
\end{proof}

% End of dd_appendix.tex

% app_o1_static.tex : the printed class at the critical weight (wave-group route).
% Input from paper.tex; no preamble. Macros below are prefixed \os.
\newcommand{\osomc}{\omega_c}        % critical weight X_0^{1/2} l^c
\newcommand{\osomp}{\omega}          % periodization weight X_0^{1/2} l^{c'}
\newcommand{\osX}{X_0}               % X_0 = cosh d(x_0,\cdot)
\newcommand{\osL}{\ell}              % l = log(2X_0)/log 2
\newcommand{\osbr}[1]{\langle #1\rangle}  % <r> = (1+r^2)^{1/2}
\newcommand{\osHK}{\widetilde p}     % continuation of the heat kernel to complex distance
\newcommand{\osLam}{\Lambda}         % Lambda = (-Delta)^{1/2}

\section{The printed class at the critical weight}\label{app:o1st}

Throughout this appendix $x_0=i$, $d=d(i,\cdot)$, $\osX=\cosh d$, $\psi:=\log\osX\ge0$, $\osL:=\log(2\osX)/\log2=1+\psi/\log2\ge1$,
$\osLam=(-\Delta)^{1/2}$ and $D=(-\Delta-\tfrac14)^{1/2}$. For $0\le c\le1$ the critical weight is
\[
\osomc:=\osX^{1/2}\osL^{\,c};
\]
the periodization weight of Section~\ref{sec:printed} is $\osomp=\omega_{c'}$. The shells are $A_n:=\{n\le\psi<n+1\}$, $n=0,1,2,\dots$, and
$1_n$ denotes multiplication by the indicator of $A_n$. Since $1+n\le\osL\le3(1+n)$ on $A_n$,
\begin{equation}\label{os:shell}
\ue^{n/2}(1+n)^c\le\osomc\le3\,\ue^{1/2}\,\ue^{n/2}(1+n)^c\qquad\text{on }A_n .
\end{equation}
For $m\ge\tfrac12$ put $\Lambda_m:=(m^2-\tfrac14-\Delta)^{1/2}=(D^2+m^2)^{1/2}$; thus $\Lambda_{1/2}=\osLam$ and $\Lambda_m\ge\osLam$.

\subsection{Profiles}

\begin{lemma}[two integrals]\label{os:int}
For $a,b>0$,
\[
\int_0^\infty\ue^{-a^2u^2-b^2/u^2}\,du=\frac{\sqrt\pi}{2a}\,\ue^{-2ab},\qquad\int_0^\infty u^{-2}\,\ue^{-a^2u^2-b^2/u^2}\,du=\frac{\sqrt\pi}{2b}\,\ue^{-2ab}.
\]
\end{lemma}
\begin{proof}
Let $F(u)$ be the integrand of the first integral. The map $u\mapsto v=au-b/u$ is a bijection of $(0,\infty)$ onto $\mathbb R$ with
$dv=(a+b/u^2)\,du$ and $a^2u^2+b^2/u^2=v^2+2ab$, so $\int_0^\infty(a+b/u^2)F\,du=\sqrt\pi\,\ue^{-2ab}$. The substitution $u\mapsto b/(au)$
preserves $F$ and exchanges $a\int F$ with $b\int u^{-2}F$, so the two are equal.
\end{proof}

For $m\ge\tfrac12$ and $\tau>0$ define the even positive function
\[
g^m_\tau(r):=\frac{\tau}{2\pi}\int_0^\infty u^{-3}\exp\Bigl(-m^2u^2-\frac{r^2+\tau^2}{4u^2}\Bigr)du .
\]

\begin{lemma}[profiles]\label{os:prof}
Let $m\ge\tfrac12$ and $\tau,s>0$.
\begin{itemize}
\item[(i)] $\int_{\mathbb R}g^m_\tau(r)\cos(r\nu)\,dr=\ue^{-\tau(\nu^2+m^2)^{1/2}}$ for real $\nu$. Hence $\ue^{-\tau\Lambda_m}=\int_{\mathbb R}g^m_\tau(r)\cos(rD)\,dr$.
\item[(ii)] For $\kappa\ge0$ with $2\kappa^2<m^2$, $\int_{\mathbb R}g^m_\tau(r)\,\ue^{\kappa|r|}\,dr\le\sqrt2$.
\item[(iii)] The function $h^m_s:=-\partial_sg^m_s$ satisfies $\Lambda_m\ue^{-s\Lambda_m}=\int_{\mathbb R}h^m_s(r)\cos(rD)\,dr$ and
$|h^m_s|\le(\sqrt2\,c_H/s)\,g^m_{s/\sqrt2}$, with $c_H:=1+4/\ue$. Hence $\int_{\mathbb R}|h^m_s(r)|\,\ue^{\kappa|r|}\,dr\le2c_H/s$ for $\kappa$ as
in (ii).
\item[(iv)] For $m=\tfrac12$ and $0<\tau\le1$: $\int_{|r|\ge T}g^{1/2}_\tau(r)\,dr\le B_0\,\ue^{-T/2}(1+T)^{-3/2}$ for all $T\ge0$, with
an absolute constant $B_0$.
\item[(v)] The function $G_m(r):=\frac1{2\pi}\int_0^\infty\bigl(\ue^{-u^2/4}-\ue^{-m^2u^2}\bigr)u^{-3}\ue^{-r^2/(4u^2)}\,du\ge0$ is integrable,
$\int_{\mathbb R}G_m(r)\cos(r\nu)\,dr=(\nu^2+m^2)^{1/2}-(\nu^2+\tfrac14)^{1/2}$, and
$\int_{|r|\ge T}G_m\le B_m\,\ue^{-T/2}(1+T)^{-3/2}$ for all $T\ge0$, with a constant $B_m$ depending only on $m$.
Hence $\Lambda_m-\osLam=\int_{\mathbb R}G_m(r)\cos(rD)\,dr$.
\end{itemize}
\end{lemma}
\begin{proof}
We use, for $u>0$ and $T>0$,
\begin{gather*}
\int_{\mathbb R}\ue^{-r^2/(4u^2)}\cos(r\nu)\,dr=2\sqrt\pi\,u\,\ue^{-u^2\nu^2},\qquad
\int_{\mathbb R}\ue^{-r^2/(4u^2)}\,\ue^{\kappa|r|}\,dr\le\sqrt2\cdot2\sqrt\pi\,u\,\ue^{2\kappa^2u^2},\\
\int_{|r|\ge T}\ue^{-r^2/(4u^2)}\,dr\le\frac{4u^2}{T}\,\ue^{-T^2/(4u^2)} .
\end{gather*}
The second holds because $\kappa|r|\le r^2/(8u^2)+2\kappa^2u^2$ and $\int_{\mathbb R}\ue^{-r^2/(8u^2)}dr=\sqrt2\cdot2\sqrt\pi\,u$.
All integrands are positive or dominated, so Tonelli and Fubini apply.

(i) The left side is $\frac{\tau}{\sqrt\pi}\int_0^\infty u^{-2}\exp(-(\nu^2+m^2)u^2-\tau^2/(4u^2))\,du$, and Lemma~\ref{os:int} with
$a=(\nu^2+m^2)^{1/2}$, $b=\tau/2$ gives the value. The operator identity is the spectral calculus of $D$.

(ii) In the same way, with $a=(m^2-2\kappa^2)^{1/2}>0$, $\int g^m_\tau(r)\,\ue^{\kappa|r|}\,dr\le\sqrt2\,\ue^{-\tau a}\le\sqrt2$.

(iii) Differentiating under the integral,
$h^m_s(r)=\frac1{2\pi}\int_0^\infty\bigl(\frac{s^2}{2u^5}-\frac1{u^3}\bigr)\exp(-m^2u^2-\frac{r^2+s^2}{4u^2})\,du$. With $x=s^2/u^2$,
$|x/2-1|\,\ue^{-x/4}\le(1+x/2)\ue^{-x/4}\le c_H\,\ue^{-x/8}$, which gives the domination. The cosine transform of $h^m_s$ is $-\partial_s$ of
(i), and the last bound follows from (ii) for $g^m_{s/\sqrt2}$.

(iv) The mass is $\ue^{-\tau/2}\le1$ by (i) at $\nu=0$. For $T>0$, with $\varrho:=(T^2+\tau^2)^{1/2}$ and
$\Phi(u):=\exp(-u^2/4-\varrho^2/(4u^2))$,
\begin{align*}
\int_{|r|\ge T}g^{1/2}_\tau&\le\frac{2\tau}{\pi T}\int_0^\infty u^{-1}\Phi\,du\le\frac{2\tau}{\pi T}\Bigl(\int_0^\infty u^{-2}\Phi\Bigr)^{1/2}\Bigl(\int_0^\infty\Phi\Bigr)^{1/2}\\
&=\frac{2\tau}{\sqrt\pi\,T}\,\varrho^{-1/2}\ue^{-\varrho/2}\le\frac{2\tau}{\sqrt\pi}\,T^{-3/2}\ue^{-T/2},
\end{align*}
by Lemma~\ref{os:int} with $a=\tfrac12$, $b=\varrho/2$. Use this for $T\ge1$ and the mass bound for $T<1$; $B_0=2^{3/2}\ue^{1/2}$ works.

(v) From $\int_0^\infty u^{-2}(1-\ue^{-Au^2})\,du=\sqrt{\pi A}$ (integrate by parts),
\[
(\nu^2+m^2)^{1/2}-(\nu^2+\tfrac14)^{1/2}=\pi^{-1/2}\int_0^\infty u^{-2}\bigl(\ue^{-u^2/4}-\ue^{-m^2u^2}\bigr)\ue^{-u^2\nu^2}\,du,
\]
which is the cosine transform of $G_m$. At $\nu=0$ this gives $\int G_m=m-\tfrac12$. Since $\ue^{-u^2/4}-\ue^{-m^2u^2}\le\ue^{-u^2/4}$, the tail of
$G_m$ is at most $\frac{2}{\pi T}\int_0^\infty u^{-1}\exp(-u^2/4-T^2/(4u^2))\,du\le\frac2{\sqrt\pi}T^{-3/2}\ue^{-T/2}$, as in (iv). With the mass
bound for $T<1$ this gives the tail bound.
\end{proof}

\subsection{Finite propagation speed and the multiplier theorem}

\begin{lemma}[finite propagation speed]\label{os:speed}
Let $f,h\in L^2(\Hb)$ with $\operatorname{supp}f\subset\{\psi\le a\}$ and $\operatorname{supp}h\subset\{\psi\ge b\}$. Then
$\langle h,\cos(rD)f\rangle=0$ and $\langle f,\cos(rD)h\rangle=0$ for $|r|<b-a$.
\end{lemma}
\begin{proof}
For $\lambda>0$ apply Lemma~\ref{lem:wave}(b) to $W_{\lambda,n}$, with $l=\lambda$, and let $n\to\infty$:
$\norm{\osX^\lambda\cos(rD)f}\le\ue^{(\lambda+1/2)|r|}\norm{\osX^\lambda f}\le\ue^{(\lambda+1/2)|r|+\lambda a}\norm f$. Hence
$\norm{1_{\{\psi\ge b\}}\cos(rD)f}\le\ue^{|r|/2}\ue^{-\lambda(b-a-|r|)}\norm f$, which tends to $0$ as $\lambda\to\infty$. The second identity follows
because $\cos(rD)$ is self-adjoint.
\end{proof}

\begin{theorem}[multiplier theorem at the critical weight]\label{os:cw}
Let $\tfrac12<c<1$. Let $g\in L^1(\mathbb R)$ be real with
\[
\int_{|r|\ge T}|g(r)|\,dr\le B\,\ue^{-T/2}(1+T)^{-3/2}\qquad\text{for all }T\ge0,
\]
and let $\mathsf G:=\int_{\mathbb R}g(r)\cos(rD)\,dr$, a bounded operator on $L^2(\Hb)$. Then for all $f\in L^2(\Hb)$,
\[
\norm{\osomc\,\mathsf Gf}\le C_c\,B\,\norm{\osomc f},
\]
with $C_c$ depending only on $c$.
\end{theorem}
\begin{proof}
Put $f_n=1_nf$ and $a_n=\norm{\osomc f_n}$, and let $T(l,n):=\max(|l-n|-1,0)$. By Lemma~\ref{os:speed}, for $h$ supported in $A_l$ the function
$r\mapsto\langle h,\cos(rD)f_n\rangle$ vanishes on $|r|<T(l,n)$. Hence
$\norm{1_l\mathsf Gf_n}\le\int_{|r|\ge T(l,n)}|g|\,\norm{f_n}\le B\,\ue^{-T/2}(1+T)^{-3/2}\norm{f_n}$, and by \eqref{os:shell}
\[
\norm{\osomc1_l\mathsf Gf_n}\le\mathsf B(l,n)\,a_n,\qquad\mathsf B(l,n):=C\,B\,\Bigl(\frac{1+l}{1+n}\Bigr)^c(1+|l-n|)^{-3/2}\times
\begin{cases}1,&l\ge n,\\ \ue^{-(n-l)},&l<n.\end{cases}
\]
Since $\mathsf Gf=\sum_n\mathsf Gf_n$ in $L^2$, $\norm{\osomc\mathsf Gf}^2\le\sum_l\bigl(\sum_n\mathsf B(l,n)a_n\bigr)^2$. We apply the Schur test on
$\mathbb N_0$ with the weights $p_n=(1+n)^{-1/2}$.
\begin{itemize}
\item $\sum_n\mathsf B(l,n)p_n\le C'B\,p_l$. For $n\ge l/2$: $((1+l)/(1+n))^c\le2^c$, $p_n\le\sqrt2\,p_l$, and $\sum_k(1+k)^{-3/2}<\infty$. For $n<l/2$:
$\mathsf B(l,n)p_n\le CB(1+l)^{c-3/2}(1+n)^{-c-1/2}$; the sum over $n$ converges because $c>\tfrac12$, and $(1+l)^{c-3/2}\le p_l$ because $c<1$.
\item $\sum_l\mathsf B(l,n)p_l\le C'B\,p_n$. For $l<n$: $((1+l)/(1+n))^cp_l\le p_n$ because $c>\tfrac12$. For $n\le l\le2n+1$: $((1+l)/(1+n))^c\le2^c$ and
$p_l\le p_n$. For $l>2n+1$: $\mathsf B(l,n)p_l\le CB(1+n)^{-c}(1+l)^{c-2}$, and $\sum_{l>2n+1}(1+l)^{c-2}\le C(1+n)^{c-1}$ because $c<1$.
\end{itemize}
So $\sum_l(\sum_n\mathsf B(l,n)a_n)^2\le(C'B)^2\sum_na_n^2=(C'B)^2\norm{\osomc f}^2$.
\end{proof}

\begin{corollary}\label{os:cwp}
Let $\tfrac12<c<1$. For $0<\tau\le1$ and $f\in L^2$, $\norm{\osomc\,\ue^{-\tau\osLam}f}\le K_c\norm{\osomc f}$. For $m\ge\tfrac12$,
$\norm{\osomc(\Lambda_m-\osLam)f}\le K_{m,c}\norm{\osomc f}$.
\end{corollary}
\begin{proof}
Theorem~\ref{os:cw} with $g=g^{1/2}_\tau$ and Lemma~\ref{os:prof}(i),(iv), respectively with $g=G_m$ and Lemma~\ref{os:prof}(v).
\end{proof}

\subsection{The data of the printed class}

\begin{proposition}[factorization of the data]\label{os:data}
Let $\tfrac12<c<1$, $m\ge\tfrac12$ and $0<\sigma<\alpha\le1$ with $\sigma m\le\tfrac12$, and put
\[
\mathsf R:=\ue^{-(\alpha-\sigma)\osLam}\,\ue^{\sigma(\Lambda_m-\osLam)},
\]
a bounded operator. Then $\ue^{-\sigma\Lambda_m}\mathsf R=\ue^{-\alpha\osLam}$, and
$\norm{\osomc\mathsf Rf}\le C_{\mathrm d}\norm{\osomc f}$ for all $f\in L^2$, with a constant $C_{\mathrm d}$ depending only on $m$ and $c$.
Consequently, if $\beta>\tfrac12$ and $\norm{\ue^{\beta\langle d\rangle}\ue^{\alpha\osLam}a}<\infty$, then $\tilde f:=\mathsf R\,\ue^{\alpha\osLam}a$ satisfies
\[
a=\ue^{-\sigma\Lambda_m}\tilde f,\qquad\norm{\osomc\tilde f}\le C\,\norm{\ue^{\beta\langle d\rangle}\ue^{\alpha\osLam}a},
\]
with $C=C(m,c,\beta)$.
\end{proposition}
\begin{proof}
The identity is the bounded functional calculus: $-\sigma\Lambda_m+\sigma(\Lambda_m-\osLam)-(\alpha-\sigma)\osLam=-\alpha\osLam$. The operator
$\mathsf B:=\Lambda_m-\osLam$ is bounded, and $\ue^{\sigma\mathsf B}=\sum_k\sigma^k\mathsf B^k/k!$ in operator norm. By Corollary~\ref{os:cwp},
$\norm{\osomc\mathsf B^kf}\le K_{m,c}^k\norm{\osomc f}$, so the partial sums are bounded in the weighted norm by $\ue^{\sigma K_{m,c}}\norm{\osomc f}$, and
Fatou's lemma gives $\norm{\osomc\ue^{\sigma\mathsf B}f}\le\ue^{K_{m,c}}\norm{\osomc f}$. Then apply Corollary~\ref{os:cwp} to
$\ue^{-(\alpha-\sigma)\osLam}$; so $C_{\mathrm d}=K_c\,\ue^{K_{m,c}}$ works. (The condition $\sigma m\le\tfrac12$ is not used in this argument; it
is the range in which the statement is formalized and in which it is applied.) For the consequence put $f=\ue^{\alpha\osLam}a$, so $a=\ue^{-\alpha\osLam}f$. Since $\psi\le d\le\langle d\rangle+1$ and
$\beta>\tfrac12$, $\osomc\le\ue^{d/2}(1+d/\log2)\le C_\beta\,\ue^{\beta\langle d\rangle}$.
\end{proof}

\subsection{Words at a power weight}

\begin{lemma}[the even semigroup at a bounded weight]\label{os:words}
Let $m\ge\tfrac12$ and $\rho\ge0$ with $2(\rho+\tfrac12)^2<m^2$, and let $W$ be a bounded weight (Appendix~\ref{app:ws}) with
$(X\log W)^2+(Y\log W)^2\le\rho^2$. Put $\mu:=m^2-\tfrac14$, $C_0:=2\ue\,c_H$ and $Q:=4+2\rho$.
\begin{itemize}
\item[(a)] For $\tau,s>0$ and $u\in L^2$: $\norm{W\ue^{-\tau\Lambda_m}u}\le\sqrt2\norm{Wu}$ and
$\norm{W\Lambda_m\ue^{-s\Lambda_m}u}\le(2c_H/s)\norm{Wu}$.
Hence for $j\ge1$ and $\sigma>0$, $\norm{W(\mu-\Delta)^j\ue^{-\sigma\Lambda_m}u}\le(4jc_H/\sigma)^{2j}\norm{Wu}$.
\item[(b)] For $0<\sigma\le1$, every word $Z_S$ of length $k$ in $X,Y$, and $h\in L^2$,
\[
\norm{WZ_S\,\ue^{-\sigma\Lambda_m}h}\le\sqrt2\,Q^k\,k!\,(C_0/\sigma)^k\,\norm{Wh}.
\]
\end{itemize}
In particular, let $\gamma\ge0$ with $2(\gamma+\tfrac12)^2<m^2$ and $\osX^\gamma h\in L^2$. Then, with $\rho=\gamma$ in the definition of $Q$,
\[
\norm{\osX^\gamma Z_S\,\ue^{-\sigma\Lambda_m}h}\le\sqrt2\,Q^k\,k!\,(C_0/\sigma)^k\,\norm{\osX^\gamma h}.
\]
Moreover $\norm{\ue^{\sigma\osLam}\ue^{-\sigma\Lambda_m}}\le1$.
\end{lemma}
\begin{proof}
(a) By Lemma~\ref{os:prof}(i), Lemma~\ref{lem:wave}(b) and Lemma~\ref{os:prof}(ii) with $\kappa=\rho+\tfrac12$,
$\norm{W\ue^{-\tau\Lambda_m}u}\le\int g^m_\tau(r)\ue^{\kappa|r|}dr\,\norm{Wu}\le\sqrt2\norm{Wu}$. The same with $h^m_s$ and Lemma~\ref{os:prof}(iii) gives the second
bound. For the third write $(\mu-\Delta)^j\ue^{-\sigma\Lambda_m}=(\Lambda_m\ue^{-(\sigma/2j)\Lambda_m})^{2j}$.

(b) Let $v_j:=(\mu-\Delta)^j\ue^{-\sigma\Lambda_m}h$. Since $\Lambda_m\ge\osLam$, $v_j\in\operatorname{Dom}(\ue^{\sigma'\osLam})$ for $\sigma'<\sigma$, so $v_j$ is smooth
and all words applied to $v_j$ lie in $L^2$ (Lemmas~\ref{lem:wc} and~\ref{lem:lift}). Put
\[
A_{k,j}:=\max_{|S|=k}\norm{WZ_Sv_j},\qquad\mathsf B_{k+1,j}:=\max_{|S|=k}\bigl(\norm{WXZ_Sv_j}^2+\norm{WYZ_Sv_j}^2\bigr)^{1/2},
\]
so $A_{k+1,j}\le\mathsf B_{k+1,j}$. Fix $S$ with $|S|=k$ and $w=Z_Sv_j$. Integration by parts gives
\begin{multline*}
\norm{WXw}^2+\norm{WYw}^2=\mathrm{Re}\langle W^2w,(\mu-\Delta)w\rangle-\mu\norm{Ww}^2\\
-\mathrm{Re}\langle X(W^2)w,Xw\rangle-\mathrm{Re}\langle Y(W^2)w,Yw\rangle .
\end{multline*}
Since $[\Delta,X]=2XY$ and $[\Delta,Y]=-2X^2$, $[\Delta,Z_S]$ is a sum of $k$ words of length $k+1$ with coefficients $\pm2$, and
$(\mu-\Delta)w=Z_Sv_{j+1}-[\Delta,Z_S]v_j$. The last two terms are at most $2\rho\norm{Ww}(\norm{WXw}^2+\norm{WYw}^2)^{1/2}$ together. Hence
\[
\mathsf B_{k+1,j}^2\le A_{k,j}A_{k,j+1}+(2k+2\rho)A_{k,j}\,\mathsf B_{k+1,j},
\]
and therefore
\[
A_{k+1,j}\le\mathsf B_{k+1,j}\le(2k+2\rho)A_{k,j}+(A_{k,j}A_{k,j+1})^{1/2}.
\]
We claim $A_{k,j}\le\sqrt2\,Q^k(C_0/\sigma)^{k+2j}(k+2j)!\,\norm{Wh}$ for all $k,j\ge0$. For $k=0$ this is (a), because
$(4jc_H/\sigma)^{2j}\le(2\ue c_H/\sigma)^{2j}(2j)!$. If it holds for $k$ and all $j$, then with $n=k+2j$, using $(n!\,(n+2)!)^{1/2}\le\sqrt2\,(n+1)!$,
$(2k+2\rho)\,n!\le(2+2\rho)(n+1)!$ and $(C_0/\sigma)^n\le(C_0/\sigma)^{n+1}/C_0$,
\[
A_{k+1,j}\le\sqrt2\,Q^k(C_0/\sigma)^{n+1}(n+1)!\,\bigl((2+2\rho)/C_0+\sqrt2\bigr)\norm{Wh},
\]
which is the claim for $k+1$, because $C_0\ge1$ and so $(2+2\rho)/C_0+\sqrt2\le Q$. The case $j=0$ is (b). The statement for $\osX^\gamma$ follows with $W=W_{\gamma,n}$ and $n\to\infty$, and the last
statement from $\Lambda_m\ge\osLam$.
\end{proof}

\subsection{Transfer from a power weight to the critical weight}

\begin{lemma}[transfer over shells]\label{os:transfer}
Let $0\le c\le1$. Let $T:L^2(\Hb)\to L^2(\Hb)$ be additive, and let $0<M_0\le M_1$ be such that
\begin{itemize}
\item[(a)] $\norm{Tu}\le M_0\norm u$ for all $u\in L^2$, and
\item[(b)] $\norm{\osX\,Tu}\le M_1\norm{\osX u}$ for every $u$ supported in one shell $A_n$.
\end{itemize}
Then for every $h\in L^2$,
\[
\norm{\osomc\,Th}\le K\,M_0^{1/4}M_1^{3/4}\,\norm{\osomc h},
\]
with a constant $K$ depending only on $c$.
\end{lemma}
\begin{proof}
Let $u=1_nh$. By (a), $\norm{1_lTu}\le M_0\norm u$. By (b), since $\osX\ge\ue^l$ on $A_l$ and $\osX\le\ue^{n+1}$ on $A_n$,
$\norm{1_lTu}\le\ue\,M_1\ue^{-(l-n)}\norm u$. For $l\ge n$, $\min(x,y)\le x^{1/4}y^{3/4}$ gives
$\norm{1_lT1_n}\le\ue^{3/4}M_0^{1/4}M_1^{3/4}\ue^{-3(l-n)/4}$; for $l<n$ we use $M_0\le M_0^{1/4}M_1^{3/4}$. By \eqref{os:shell} and $c\le1$,
$\sup_{A_l}\osomc/\inf_{A_n}\osomc\le3\ue^{1/2}\ue^{(l-n)/2}(1+(l-n)_+)$. So with $a_n:=\norm{\osomc1_nh}$,
\[
\norm{\osomc1_lT1_nh}\le M_0^{1/4}M_1^{3/4}\,k(l-n)\,a_n,\qquad k(j):=\begin{cases}3\ue^{5/4}(1+j)\ue^{-j/4},&j\ge0,\\3\ue^{1/2}\ue^{-|j|/2},&j<0.\end{cases}
\]
Since $T$ is additive and bounded, $Th=\sum_nT1_nh$ in $L^2$, and therefore
$\norm{\osomc Th}^2\le\sum_l\bigl(\sum_n\norm{\osomc1_lT1_nh}\bigr)^2\le M_0^{1/2}M_1^{3/2}\bigl(\sum_jk(j)\bigr)^2\sum_na_n^2$ by Young's inequality for
sequences.
\end{proof}

\subsection{Periodization at the critical weight}

\begin{proposition}[periodization and transport at the critical weight]\label{os:p1}
Let $\Gamma$ be any Fuchsian group (cusps allowed), $\tfrac12<c'\le1$, $\osomp=\omega_{c'}$, $t\ge0$, $\hbar>0$, $\varphi=\varphi_\infty=-\log y$, and
$w\in L^2(\Hb)$ with $\osomp w\in L^2(\Hb)$. Then the sum defining $[(S(t)w)\ue^{i\varphi/\hbar}]_\Gamma$ converges absolutely almost everywhere, and
\[
\bigl\|[(S(t)w)\ue^{i\varphi/\hbar}]_\Gamma\bigr\|_{L^2(M)}\le C(\Gamma,c')\,\ue^{t/2}(1+t/\log2)^{c'}\,\|\osomp w\|_{L^2(\Hb)}.
\]
In particular $\osomp$ satisfies \textup{(W1)} and \textup{(W2)} of Section~\ref{subsec:dd-surf}, with
$\osomp(z)/\osomp(\delta_sz)\le\ue^{|s|/2}(1+|s|/\log2)^{c'}\le(2/\log2)\,\ue^{|s|}$.
\end{proposition}
\begin{proof}
\emph{Packing.} Since $\Gamma$ is discrete and $x_0$ has finite stabilizer, the orbit $\Gamma x_0$ has a fixed positive separation $2r_0$ after
removing stabilizer multiplicity (proper discontinuity). Balls $B(\gamma x_0,r_0)$ are disjoint up to bounded multiplicity, and those with
$d(z,\gamma x_0)<n+1$ lie in $B(z,n+1+r_0)$. Comparing hyperbolic areas gives $\#\{\gamma:d(z,\gamma x_0)<n+1\}\le C_\Gamma\ue^n$ uniformly in $z\in\Hb$, cusps
included.

\emph{Pointwise Cauchy--Schwarz.} Since $\osomp(w)^2\asymp\ue^{d(w,x_0)}\osbr{d(w,x_0)}^{2c'}$ and $d(\gamma z,x_0)=d(z,\gamma^{-1}x_0)$,
reindexing $\gamma\mapsto\gamma^{-1}$ gives
\[
\sum_\gamma\osomp(\gamma z)^{-2}\le C\sum_\gamma\ue^{-d(z,\gamma x_0)}\osbr{d(z,\gamma x_0)}^{-2c'}\le C\sum_{n\ge0}\ue^{n}\ue^{-n}(1+n)^{-2c'}<\infty,
\]
since $2c'>1$. By Cauchy--Schwarz, $(\sum_\gamma|f(\gamma z)|)^2\le C\sum_\gamma\osomp(\gamma z)^2|f(\gamma z)|^2$ with $f=(S(t)w)\ue^{i\varphi/\hbar}$.

\emph{Unfolding.} Integrating over a fundamental domain (Tonelli; a finite stabilizer changes only the constant) gives
$\|[f]_\Gamma\|^2_{L^2(M)}\le C\int_\Hb\osomp^2|S(t)w|^2\,d\nu$.

\emph{Transport.} $(S(t)w)(z)=\ue^{-t/2}w(\delta_tz)$, and $d(z,\delta_sz)=|s|$. Since $\psi$ is $1$-Lipschitz, $\osX(z)\le\ue^{|s|}\osX(\delta_sz)$ and
$\osL(z)\le\osL(\delta_sz)+|s|/\log2\le\osL(\delta_sz)(1+|s|/\log2)$; this is the ratio bound. Since $S(t)$ is unitary,
$\int\osomp^2|S(t)w|^2d\nu\le\ue^{t}(1+t/\log2)^{2c'}\int\osomp^2|w|^2d\nu$. Combining gives the claim.
\end{proof}

\begin{remark}[the surface object does not depend on the weight]\label{os:indep}
Let $\omega_1,\omega_2$ be two weights satisfying \textup{(W1)}, \textup{(W2)}. For $g\in L^2_{\omega_1}\cap L^2_{\omega_2}$ the element
$\mathcal V_sg=\ue^{-is/(2\hbar)}\cU(-s)[\ue^{i\varphi/\hbar}S(s)g]_\Gamma$ of $L^2(M)$ is the same for both weights, because the periodization is in both
cases the same almost everywhere absolutely convergent sum. In particular this applies to $\omega_1=\osX^{\gamma}$ with $\gamma>\tfrac12$ and
$\omega_2=\osomp$, and $L^2_{\osX^\gamma}\subset L^2_{\osomp}$.
\end{remark}

\subsection{Kernels of the Poisson and heat semigroups}

The two lemmas of this subsection are used in Section~\ref{subsec:cl} (Lemma~\ref{cl:kernel}). They are not used in the proof of
Theorem~\ref{pc:main}.
Let $\mu_\tau(s)=\tau(4\pi)^{-1/2}s^{-3/2}e^{-\tau^2/(4s)}$ (so that $e^{-\tau\osLam}=\int_0^\infty\mu_\tau(s)e^{s\Delta}\,ds$, by
the subordination identity $e^{-\tau\sqrt A}=\int_0^\infty\mu_\tau(s)e^{-sA}ds$, $A\ge0$), and let $p_s$ be the heat kernel of $\Hb$, which is
given by McKean's formula
\[
p_s(d)=\frac{\sqrt2\,e^{-s/4}}{(4\pi s)^{3/2}}\int_d^\infty\frac{\rho\,e^{-\rho^2/(4s)}}{(\cosh\rho-\cosh d)^{1/2}}\,d\rho
\]
(\cite[(1.7)]{GN98}). Let
$k_\tau(r)=\int_0^\infty\mu_\tau(s)p_s(r)\,ds$ be the radial kernel of $e^{-\tau\osLam}$.

\begin{lemma}[Poisson kernel representation and bounds]\label{os:poisson}
For $0<\tau\le\tau_0<\pi/2$ and $r\ge0$,
\[
k_\tau(r)=\frac{\sqrt2}{8\pi^2}\,\tau\int_r^\infty\frac{\rho\,K_2\bigl(\sqrt{\rho^2+\tau^2}/2\bigr)}{\rho^2+\tau^2}\,\frac{d\rho}{(\cosh\rho-\cosh r)^{1/2}}.
\]
Consequently there are $0<c_\tau,C_\tau<\infty$, depending only on $\tau_0$, with
\begin{align}
k_\tau(r)&\le C\,\tau\,e^{-r}\osbr{r}^{-3/2}\quad(r\ge1),\label{os:ub}\\
k_\tau(r)&\ge c\,\tau\,e^{-r}\osbr{r}^{-3/2}\quad(r\ge1),\label{os:lb}\\
c\,\tau(\tau^2+r^2)^{-3/2}\le k_\tau(r)&\le C\,\tau(\tau^2+r^2)^{-3/2}\quad(0\le r\le1).\label{os:small}
\end{align}
\end{lemma}
\begin{proof}
Insert the formula for $p_s$ into $k_\tau=\int\mu_\tau p_s\,ds$ and interchange the integrals (Tonelli). The $s$-integral is
\[
\int_0^\infty s^{-3}e^{-(\rho^2+\tau^2)/(4s)-s/4}\,ds=2\frac{1/4}{(\rho^2+\tau^2)/4}K_2\Bigl(\frac{\sqrt{\rho^2+\tau^2}}2\Bigr)
=\frac{2K_2(\sqrt{\rho^2+\tau^2}/2)}{\rho^2+\tau^2},
\]
by $\int_0^\infty s^{\nu-1}e^{-\beta s-\gamma/s}ds=2(\gamma/\beta)^{\nu/2}K_\nu(2\sqrt{\beta\gamma})$ (\cite{DLMF}, 10.32.10, and the
standard Bessel integral) with $\nu=-2$, $\beta=\tfrac14$, $\gamma=(\rho^2+\tau^2)/4$, and $K_{-2}=K_2$. The constant is
$\tfrac{\tau}{2\sqrt\pi}\cdot\sqrt2\cdot(4\pi)^{-3/2}\cdot2=\tfrac{\sqrt2}{8\pi^2}\tau$. This gives the representation.

For the bounds write $\rho=r+v$, so $\cosh\rho-\cosh r=2\sinh(r+v/2)\sinh(v/2)$. For $r\ge1$ the factor
$K_2(\sqrt{\rho^2+\tau^2}/2)$ is comparable to $e^{-\rho/2}\rho^{-1/2}$ uniformly for $\tau\le\tau_0$, the factor $\rho/(\rho^2+\tau^2)$
is at most $\rho^{-1}\le r^{-1}$, and for $0\le v\le1$ it is at least $((1+\tau_0^2)\,2r)^{-1}$; also $(\cosh\rho-\cosh r)^{-1/2}=(2\sinh(r+v/2)\sinh(v/2))^{-1/2}$ is at most $Ce^{-r/2}(\sinh(v/2))^{-1/2}$ for all $v>0$ and comparable to it for $0<v\le1$. The $v$-integral of
$e^{-v/2}(\sinh(v/2))^{-1/2}$ over $(0,\infty)$ is finite, which gives \eqref{os:ub}; its integral over $(0,1]$ is positive, and restricting to $0\le v\le1$ gives \eqref{os:lb}. For $r\le1$ split the
integral at $\rho=2$; there $\cosh\rho-\cosh r$ is comparable to $\rho^2-r^2$, and with $v=\sqrt{\rho^2-r^2}$ the main part is
comparable to $\int_0^{C}(r^2+\tau^2+v^2)^{-2}\,dv\asymp(r^2+\tau^2)^{-3/2}$, while the tail $\rho\ge2$ is bounded by the
exponential estimate. This gives \eqref{os:small}.
\end{proof}

Let $\osHK$ denote the continuation of $p_s$ to complex distance $\tilde d=a+i\theta$, $a\ge0$, $|\theta|\le\theta_0<\pi/2$.
McKean's formula for $p_s$ is first put in $d^2$-form; in the rest of this subsection $\tau$ denotes an integration variable, not the
parameter of Lemma~\ref{os:poisson}. With $g(w)=\sinh\sqrt w/(2\sqrt w)$ (an entire function of $w$) and
$F(w,\tau^2)=\int_0^1g(w+\lambda\tau^2)\,d\lambda=(\cosh\sqrt{w+\tau^2}-\cosh\sqrt w)/\tau^2$, substituting $\rho=\sqrt{d^2+\tau^2}$ gives
\[
p_s(d)=c_s\,e^{-d^2/(4s)}\int_0^\infty e^{-\tau^2/(4s)}F(d^2,\tau^2)^{-1/2}\,d\tau,\qquad c_s=\frac{\sqrt2e^{-s/4}}{(4\pi s)^{3/2}}.
\]

\begin{lemma}[Continued heat kernel]\label{os:hk}
Let $0<\theta_0<\pi/2$ and $c_{\theta_0}=\cos\theta_0\,e^{-\theta_0}$. For $a\ge0$, $|\theta|\le\theta_0$ and $s>0$,
\[
|\osHK_s(a+i\theta)|\le c_{\theta_0}^{-1/2}\,e^{\theta^2/(4s)}\,p_s(a).
\]
\end{lemma}
\begin{proof}
Put $d^2\mapsto\tilde d^2=a^2-\theta^2+2ia\theta$ in the $d^2$-representation. Let $x=a^2+\lambda\tau^2\ge a^2$ and $w=x-\theta^2+2ia\theta$,
and write $\sqrt w=A+iB$ with $A\ge0$. Then $|w|\le x+\theta^2$, so $B^2=(|w|-(x-\theta^2))/2\le\theta^2$; also $A^2=x-\theta^2+B^2\le x$.
Hence $0\le\sqrt x-A$, and if $x\ge\theta^2$ then $\sqrt x-A\le|\theta|$.

Now $\operatorname{Re}\frac{\sinh\zeta}{\zeta}=\frac{A\sinh A\cos B+B\cosh A\sin B}{A^2+B^2}$ for $\zeta=A+iB$. Since $|B|<\pi/2$ we have
$B\sin B\ge B^2\cos B$ (for either sign of $B$), and $\cosh A\ge\sinh A/A$, so $\operatorname{Re}(\sinh\zeta/\zeta)\ge\cos B\,\sinh A/A\ge\cos\theta_0\,\sinh A/A$. The
logarithmic derivative of $\sinh u/u$ lies in $[0,1]$, so $\sinh A/A\ge e^{-|\theta|}\sinh\sqrt x/\sqrt x$. Therefore
\[
\operatorname{Re}g(w)\ge c_{\theta_0}\,g(x)>0.
\]
(The case $x<\theta^2$ is handled directly, since then $\sinh A/A\ge1\ge(\sinh\sqrt x/\sqrt x)e^{-\theta_0}$.)

Hence $\operatorname{Re}F(\tilde d^2,\tau^2)\ge c_{\theta_0}F(a^2,\tau^2)>0$. For $\operatorname{Re}z>0$ one has $|z^{-1/2}|\le(\operatorname{Re}z)^{-1/2}$, so
$|F(\tilde d^2,\tau^2)^{-1/2}|\le c_{\theta_0}^{-1/2}F(a^2,\tau^2)^{-1/2}$. Finally
$|e^{-\tilde d^2/(4s)}|=e^{\theta^2/(4s)}e^{-a^2/(4s)}$, and the $\tau$-integral is dominated uniformly on compact sets of $\tilde d$, so the continuation is
holomorphic. This gives the bound.
\end{proof}

% Formalization status. \leanrepo is the public repository; \leancommit is a commit of it.
\providecommand{\lean}[1]{\texttt{\detokenize{#1}}}
\providecommand{\leancommit}{\texttt{9a4603a}}
\providecommand{\leanrepo}{\url{https://github.com/skypher/semiclassical-wave-propagation}}

\section{Formalization}\label{app:lean}

The principal $L^2$ approximation estimates, their optimal-truncation corollaries, and the supporting statements specified below
are formalized in Lean~4 and checked by its kernel. Table~\ref{tab:lean} gives an overview of direct counterparts,
paper deductions and variants, and results without formal proofs; the text explains the substantive differences.

\subsection{What is checked}
The formal development is the directory \texttt{lean/} of the repository \leanrepo, namespace \lean{Qschubert}, at commit \leancommit. It is built with Lean~4 \cite{Lean4} (toolchain \texttt{v4.35.0}), Mathlib \cite{Mathlib} at the tag \texttt{v4.35.0}, and \texttt{physicslib4} \cite{Phys} (commit \texttt{2337279}), which supplies the spectral theorem for unbounded self-adjoint operators. The development also depends on the project \emph{Prime Number Theorem And \dots} \cite{PNT} (commit \texttt{127b50e}), from which it takes the classical zero-free region of the Riemann zeta function; this is used only in the proof of Theorem~\ref{ev:mod}.
At this commit the build and axiom audits pass: every declaration depends only on
\lean{propext}, \lean{Classical.choice} and \lean{Quot.sound}, and none uses \lean{sorryAx}.
The repository's README gives the commands for reproducing the build and audits.

A kernel check certifies the proofs of the formal statements. Whether a formal statement says what the corresponding statement of this paper says is a matter of reading; the choices made are listed next. The formal statements were compared with the text in independent AI reviews of the statements, which did not repeat the build or the axiom check. These comparisons are distinct from kernel checking and do not constitute external peer review.

\subsection{Interpretation of the formal statements}
\begin{itemize}
\item \emph{Groups and surfaces.} A Fuchsian group is a subgroup $\Gamma\le\mathrm{PGL}(2,\mathbb R)$ that acts properly discontinuously on $\Hb$ and consists of classes of matrices of positive determinant. $M$ is the orbit space with the measure induced from a fundamental domain; no manifold structure on $M$ is used.
\item \emph{Operators.} $-\Delta$ on $L^2(\Hb)$ and on $L^2(M)$ is the Friedrichs extension from test functions, respectively from periodized test functions; $\Lambda$, $\ue^{\alpha\Lambda}$, $\cU(t)$ and $\Delta^N$ are functions of it by the spectral theorem.
\item \emph{Boundary points.} The base point $x_0\in\Hb$ is a parameter, and $b=g\,\infty$ with $g\in\mathrm{SL}(2,\mathbb R)$, $g\,x_0=x_0$. Every $b\in\partial\Hb$ arises, and the objects depend on $b$ only (\lean{exists_stab_smul_infty}, \lean{initState_eq_of_stab}, \lean{approxSolB_eq_of_stab}).
\item \emph{Amplitudes.} An amplitude is a function on $\Hb$ whose class has finite norm. The formal theorems produce its $C^\infty$ representative $a'$, equal to $a$ almost everywhere, before $\hbar,t,K,N,b$ are chosen, and $u^{(K)}$ is built from $a'$. Nonzero amplitudes of finite norm exist in both classes (\lean{printed_class_nonvacuous_beta}, \lean{swappedD_class_nonvacuous}).
\item \emph{Conclusions.} Each $L^2$ theorem asserts that the two orbit sums converge absolutely almost everywhere and define elements of $L^2(M)$, and then bounds $\norm{\Delta^N[\,\cdot\,]}$, an element of $[0,\infty]$ that equals $\infty$ outside the domain of $\Delta^N$. The pointwise declarations assert that the error lies in the domain of $\Delta^N$ and that $\Delta^NE$ has a continuous representative at every point of $M$.
\item \emph{The cap on $\delta$ in the optimal-truncation corollaries.} For Corollary~\ref{pc:cor15} the condition $\delta\le\kappa^{1/2}/\max(1,C)$ is part of the formal statement. For the reversed-weight and power-weight norms in Corollary~\ref{cor:dd15} the formal statements assert that some $C$ and $\delta>0$ exist for which the bound holds on the window; their proofs take $\delta\le\kappa^{1/2}/\max(1,C)$, and $\delta$ may always be decreased, since this only shortens the time window.
\item \emph{Ranges and grouped statements.} $0<\hbar\le1$; $N^{2N}=1$ at $N=0$.
Theorem~\ref{pc:main} states the bounded-radius result \lean{pcmain_all}; its proof first treats $\alpha\le\tfrac12$,
the range of \lean{pcmain_half}. The original-norm corollaries likewise use the declarations with suffix \lean{all}.
Proposition~\ref{os:data} retains the condition $\sigma m\le\tfrac12$, although its paper proof does not need it:
this is the range of its formal statement and suffices for its applications.
The formal statement of Proposition~\ref{prop:mean} has every $\hbar>0$ and every real $t$, and no hypothesis on the area:
its inequality is one in $[0,\infty]$ with $0\cdot\infty=0$.
\end{itemize}

\subsection{The formalized counterexample}\label{rem:g5-lean}
For Proposition~\ref{prop:g5}, the formal theorem \lean{printedThm13_false} proves the negation of the bound of \cite{S},
rather than the asymptotic formula.

The statement of \cite[Theorem~1.3]{S} is formalized as: for every Fuchsian $\Gamma$ and all $\alpha>0$, $\beta>1$ there are $C,\delta,\hbar_0>0$ such that
$\norm{u^{(0)}(t)_\Gamma-\cU(t)(a\ue^{i\varphi_b/\hbar})_\Gamma}_{L^2(M)}\le C\hbar\norm a_{\alpha,\beta}$ for all $0<\hbar\le\hbar_0$, all $0\le t\le\delta\hbar^{-1/2}$, all $b$ and all $C^\infty$ amplitudes of finite original norm. The restrictions to $\hbar\le\hbar_0$ and to $C^\infty$ amplitudes weaken that claim, and the formal theorem states that this weaker claim is false. Its proof uses $\Gamma=\{\mathrm{id}\}$, $\alpha=\tfrac12$, $\beta=2$, $b=\infty$ and the amplitude $a=\ue^{-\alpha\Lambda}\ue^{-3X_0}$, $X_0=\cosh d(i,\cdot)$. It takes the Dyson identity (ii) of Proposition~\ref{prop:dd-id} at $K=0$ and $K=1$, which gives the expansion in Proposition~\ref{prop:g5} with $R$ the second-order term, and uses that periodization over the trivial group is an isometry. The bound $\norm R\le\tfrac{\hbar^2t^2}8\sup_{r,s}\norm{\Delta(r)\Delta(s)a}$ comes from Lemma~\ref{lem:wb} at $N=0$, and, with $L:=Y^2-Y$, $\norm{La}\ge\norm a/4$ from $Y^*=1-Y$ on $\operatorname{Dom}(\Delta)$: $\langle a,La\rangle=-\norm{Ya}^2$ and $2\operatorname{Re}\langle a,Ya\rangle=\norm a^2$.

\subsection{Correspondence and exclusions}
The direct counterparts in Table~\ref{tab:lean} use the interpretations above.
The file \texttt{LEAN\_STATUS.md} in the repository gives the correspondence for every labelled statement,
including the supporting estimates. The differences in bounds, parameter ranges and conclusions are described below.

\paragraph{Paper deductions and variants.}
Corollaries~\ref{pc:cor16} and~\ref{cor:dd16} display the prefactor $(\alpha\hbar)^{-2N-2}$ obtained in
Section~\ref{subsec:pc-cor} from the approximation theorem at $N$ and $N+1$ and Lemma~\ref{pc:g43}.
The listed Lean pointwise declarations instead state the weaker prefactor $(\alpha\hbar)^{-2N-4}$;
no separate declaration of the sharper estimate is asserted. For the reversed-weight and power-weight norms, the listed declarations
cover the bounded-geometry specialization; the arbitrary-quotient estimate with $m(z)^{1/2}$ is a paper deduction
from the approximation theorems and Lemma~\ref{pc:g43}.

Corollary~\ref{ev:transfer} follows from the approximation bound at $N=0$ and the bounded-observable
estimate \lean{ev_transfer_core}. Its reversed-weight version is \lean{evTransfer_main}.
For the original norm, \lean{evTransfer_pc} covers $\alpha\le\tfrac12$; \lean{pcmain_all}
supplies the approximation bound on every bounded radius range.

The fixed-radius Corollary~\ref{thm:main} is presented here for $\beta>1/2$ and every fixed $\alpha>0$. It is deduced in the text from \lean{thm14Rmain_exp} at $K=N=0$, with a fixed radius and, for larger radii, monotonicity of the reversed-weight norm. The separately formalized \lean{thm13R_rep} has the narrower range $\beta>1$ and $\alpha\le\alpha_0<\pi/4$. The Mellin, scattering, zero-mode, initial-norm and outgoing-mass lemmas of
Section~\ref{sec:expect} are ingredients of the formal proof of Theorem~\ref{ev:mod};
their individual statements do not have separate exact formal counterparts.

\paragraph{Results without formal proofs.}
The classical comparison in Theorem~\ref{ev:mod}(iii) uses the cited equidistribution theorem for closed horocycles;
only parts (i),(ii) are covered by \lean{evmod}, whose statement is \lean{Stmt_evmod}.
The amplitude-class comparison (Proposition~\ref{pc:classes}), the bound on the mean (Corollary~\ref{cor:meanJ}) and the cusp
sharpness argument of Section~\ref{sec:cusp}, including Theorem~\ref{cl:necessary}, Corollary~\ref{cl:rational}, Remark~\ref{rem:g4}
and the kernel Lemmas~\ref{os:poisson} and~\ref{os:hk}, have no formal proofs.
For cusp sharpness the development contains only implications from a formal statement of Theorem~\ref{cl:main}(ii),
which is itself not proved formally.

The example-amplitude results, Lemmas~\ref{lem:gutzmer} and~\ref{lem:ex} and Corollary~\ref{cor:localized},
are also outside the formalization and are not used by the approximation theorems.
The status of Proposition~\ref{prop:g5} is explained in the formalized-counterexample discussion above.

\begin{table}[ht]
\centering\small
\caption{Formal counterparts and the principal exceptions.}\label{tab:lean}
\begin{tabular}{@{}p{0.36\linewidth}p{0.60\linewidth}@{}}
\hline
Statement & Lean theorem or scope\\
\hline
\multicolumn{2}{@{}l}{\emph{Direct counterparts in the reading above}}\\
Theorem~\ref{pc:main} & \lean{pcmain_all}\\
Corollary~\ref{pc:thm13} & \lean{pcthm13_all}\\
Corollary~\ref{pc:cor15} & \lean{pccor15_all}\\
Lemma~\ref{pc:g43} & \lean{sobolev_M_mult}\\
Proposition~\ref{prop:mean} & \lean{meanGrowth}\\
Theorem~\ref{ev:mod}(i),(ii) & \lean{evmod}\\
Theorem~\ref{thm:dd14} & \lean{thm14Rmain_exp}\\
Corollary~\ref{cor:dd15} & \lean{cor15Rmain_rep}, \lean{cor15R_rep}\\
Theorem~\ref{thm:dd-appx} & \lean{thm14R_rep}\\
Theorem~\ref{os:cw} & \lean{cw_abstract}\\
Lemma~\ref{os:transfer} & \lean{transfer}\\
\hline
\multicolumn{2}{@{}l}{\emph{Paper deductions and variants}}\\
Corollary~\ref{pc:cor16} & \lean{pccor16_all}, \lean{pccor16bg_all}: weaker exponent $-2N-4$\\
Corollary~\ref{cor:dd16} & \lean{cor16Rmain_rep}, \lean{cor16R_rep}: bounded geometry, weaker exponent $-2N-4$\\
Corollary~\ref{thm:main} & Specialization of \lean{thm14Rmain_exp}; see text\\
Corollary~\ref{ev:transfer} & \lean{evTransfer_main}; original norm uses \lean{pcmain_all} and \lean{ev_transfer_core}\\
Proposition~\ref{prop:g5} & \lean{printedThm13_false}: negation of Schubert's bound only\\
\hline
\multicolumn{2}{@{}l}{\emph{Results without formal proofs}}\\
Corollary~\ref{cor:localized} & Paper deduction from Lemma~\ref{lem:ex} and Theorem~\ref{thm:dd14}\\
Proposition~\ref{pc:classes}; Theorem~\ref{cl:necessary}; Section~\ref{sec:cusp} & Amplitude-class comparison and cusp sharpness\\
Corollary~\ref{cor:meanJ} & Paper deduction from Proposition~\ref{prop:mean} and Theorem~\ref{thm:dd14}\\
Theorem~\ref{ev:mod}(iii) & Cited horocycle equidistribution theorem\\
\hline
\end{tabular}
\end{table}


\begin{thebibliography}{DLMF}

\bibitem[AN07]{AN07}
N.~Anantharaman and S.~Nonnenmacher,
\emph{Half-delocalization of eigenfunctions for the Laplacian on an Anosov manifold},
Ann.\ Inst.\ Fourier (Grenoble) \textbf{57} (2007), no.~7, 2465--2523,
\href{https://doi.org/10.5802/aif.2340}{doi:\allowbreak 10.5802/aif.2340},
\href{https://arxiv.org/abs/math-ph/0610019}{arXiv:math-ph/0610019}.

\bibitem[Ana08]{Ana08}
N.~Anantharaman,
\emph{Entropy and the localization of eigenfunctions},
Ann.\ of Math.\ (2) \textbf{168} (2008), no.~2, 435--475,
\href{https://doi.org/10.4007/annals.2008.168.435}{doi:\allowbreak 10.4007/annals.2008.168.435}.

\bibitem[AR12]{AR12}
N.~Anantharaman and G.~Rivi\`ere,
\emph{Dispersion and controllability for the Schr\"odinger equation on negatively curved manifolds},
Anal.\ PDE \textbf{5} (2012), no.~2, 313--338,
\href{https://doi.org/10.2140/apde.2012.5.313}{doi:\allowbreak 10.2140/apde.2012.5.313},
\href{https://arxiv.org/abs/1007.4343}{arXiv:1007.4343}.

\bibitem[BB79]{BB79}
M.~V.~Berry and N.~L.~Balazs,
\emph{Evolution of semiclassical quantum states in phase space},
J.\ Phys.\ A \textbf{12} (1979), no.~5, 625--642,
\href{https://doi.org/10.1088/0305-4470/12/5/012}{doi:\allowbreak 10.1088/0305-4470/12/5/012}.

\bibitem[BdB00]{BdB00}
F.~Bonechi and S.~De Bi\`evre,
\emph{Exponential mixing and $|\ln\hbar|$ time scales in quantized hyperbolic
maps on the torus},
Commun.\ Math.\ Phys.\ \textbf{211} (2000), no.~3, 659--686,
\href{https://doi.org/10.1007/s002200050831}{doi:\allowbreak 10.1007/s002200050831},
\href{https://arxiv.org/abs/chao-dyn/9910014}{arXiv:chao-dyn/9910014}.

\bibitem[Ber77]{Ber77}
M.~V.~Berry,
\emph{Regular and irregular semiclassical wavefunctions},
J.\ Phys.\ A \textbf{10} (1977), no.~12, 2083--2091,
\href{https://doi.org/10.1088/0305-4470/10/12/016}{doi:\allowbreak 10.1088/0305-4470/10/12/016}.

\bibitem[BGP99]{BGP99}
D.~Bambusi, S.~Graffi and T.~Paul,
\emph{Long time semiclassical approximation of quantum flows: a proof of the
Ehrenfest time},
Asymptot.\ Anal.\ \textbf{21} (1999), no.~2, 149--160,
\href{https://doi.org/10.3233/ASY-1999-365}{doi:\allowbreak 10.3233/ASY-1999-365},
\href{https://arxiv.org/abs/math-ph/9805018}{arXiv:math-ph/9805018}.

\bibitem[BR02]{BR02}
A.~Bouzouina and D.~Robert,
\emph{Uniform semiclassical estimates for the propagation of quantum
observables},
Duke Math.\ J.\ \textbf{111} (2002), no.~2, 223--252,
\href{https://doi.org/10.1215/S0012-7094-02-11122-3}{doi:\allowbreak 10.1215/S0012-7094-02-11122-3}.

\bibitem[Bro22]{Bro22}
S.~Brooks,
\emph{Eisenstein quasimodes and semiclassical behavior of expectation values},
Ann.\ Henri Poincar\'e \textbf{23} (2022), no.~12, 4429--4448,
\href{https://doi.org/10.1007/s00023-022-01184-y}{doi:\allowbreak 10.1007/s00023-022-01184-y};
preprint under the title \emph{Eisenstein series and breakdown of semiclassical correspondence},
\href{https://arxiv.org/abs/1711.01473}{arXiv:1711.01473} (2017).

\bibitem[CR97]{CR97}
M.~Combescure and D.~Robert,
\emph{Semiclassical spreading of quantum wave packets and applications near
unstable fixed points of the classical flow},
Asymptot.\ Anal.\ \textbf{14} (1997), no.~4, 377--404,
\href{https://doi.org/10.3233/ASY-1997-14405}{doi:\allowbreak 10.3233/ASY-1997-14405}.

\bibitem[DJ18]{DJ18}
S.~Dyatlov and L.~Jin,
\emph{Semiclassical measures on hyperbolic surfaces have full support},
Acta Math.\ \textbf{220} (2018), no.~2, 297--339,
\href{https://doi.org/10.4310/ACTA.2018.v220.n2.a3}{doi:\allowbreak 10.4310/ACTA.2018.v220.n2.a3},
\href{https://arxiv.org/abs/1705.05019}{arXiv:1705.05019}.

\bibitem[DLMF]{DLMF}
\emph{NIST Digital Library of Mathematical Functions},
\href{https://dlmf.nist.gov/}{https://dlmf.nist.gov/}, Release 1.2.8 of 2026-09-15.
F.~W.~J.~Olver, A.~B.~Olde Daalhuis, D.~W.~Lozier, B.~I.~Schneider, R.~F.~Boisvert, C.~W.~Clark, B.~R.~Miller,
B.~V.~Saunders, H.~S.~Cohl, and M.~A.~McClain, eds.

\bibitem[Fau07]{Fau07}
F.~Faure,
\emph{Semi-classical formula beyond the Ehrenfest time in quantum chaos. (I) Trace formula},
Ann.\ Inst.\ Fourier (Grenoble) \textbf{57} (2007), no.~7, 2525--2599,
\href{https://doi.org/10.5802/aif.2341}{doi:\allowbreak 10.5802/aif.2341},
\href{https://arxiv.org/abs/nlin/0610004}{arXiv:nlin/0610004}.

\bibitem[GN98]{GN98}
A.~Grigor'yan and M.~Noguchi,
\emph{The heat kernel on hyperbolic space},
Bull.\ London Math.\ Soc.\ \textbf{30} (1998), no.~6, 643--650,
\href{https://doi.org/10.1112/S0024609398004780}{doi:\allowbreak 10.1112/S0024609398004780}.

\bibitem[Han26]{Han26}
X.~Han,
\emph{Schr\"odinger group along the horocycles and Egorov's theorem on compact hyperbolic surfaces},
\href{https://arxiv.org/abs/2611.04131}{arXiv:2611.04131} (2026).

\bibitem[HJ00]{HJ00}
G.~A.~Hagedorn and A.~Joye,
\emph{Exponentially accurate semiclassical dynamics: propagation, localization,
Ehrenfest times, scattering, and more general states},
Ann.\ Henri Poincar\'e \textbf{1} (2000), no.~5, 837--883,
\href{https://doi.org/10.1007/PL00001017}{doi:\allowbreak 10.1007/PL00001017},
\href{https://arxiv.org/abs/math-ph/9911036}{arXiv:math-ph/9911036}.

\bibitem[IR22]{IR22}
M.~Ingremeau and A.~Rivera,
\emph{How Lagrangian states evolve into random waves},
J.\ \'Ec.\ polytech.\ Math.\ \textbf{9} (2022) 177--212,
\href{https://doi.org/10.5802/jep.181}{doi:\allowbreak 10.5802/jep.181},
\href{https://arxiv.org/abs/2011.02943}{arXiv:2011.02943}.

\bibitem[Kr\"o09]{K}
B.~Kr\"otz,
\emph{Crown theory for the upper half plane},
in: Automorphic Forms and $L$-functions I. Global aspects,
Contemp.\ Math.\ \textbf{488} (2009), 147--182,
\href{https://doi.org/10.1090/conm/488/09568}{doi:\allowbreak 10.1090/conm/488/09568},
\href{https://arxiv.org/abs/0705.1243}{arXiv:0705.1243}.

\bibitem[Mat20]{Mathlib}
The mathlib Community,
\emph{The Lean mathematical library},
in: Proceedings of the 9th ACM SIGPLAN International Conference on Certified Programs and Proofs (CPP 2020),
ACM, 2020, pp.~367--381,
\href{https://doi.org/10.1145/3372885.3373824}{doi:\allowbreak 10.1145/3372885.3373824}.

\bibitem[MU21]{Lean4}
L.~de Moura and S.~Ullrich,
\emph{The Lean 4 theorem prover and programming language},
in: Automated Deduction -- CADE 28, Lecture Notes in Comput.\ Sci.\ 12699, Springer, 2021, pp.~625--635,
\href{https://doi.org/10.1007/978-3-030-79876-5_37}{doi:\allowbreak 10.1007/978-3-030-79876-5\_37}.

\bibitem[MV07]{MV07}
H.~L.~Montgomery and R.~C.~Vaughan,
\emph{Multiplicative Number Theory I. Classical Theory},
Cambridge Stud.\ Adv.\ Math.\ 97, Cambridge University Press, 2007,
ISBN 978-0-521-84903-6.

\bibitem[Olv97]{O}
F.~W.~J.~Olver,
\emph{Asymptotics and Special Functions},
A~K~Peters, Wellesley, MA, 1997.

\bibitem[Phy]{Phys}
\emph{physicslib4}, a Lean~4 library,
\url{https://github.com/physicslib/physicslib4}; used at the commit named in Appendix~\ref{app:lean}.

\bibitem[PNT]{PNT}
A.~Kontorovich, T.~Tao et al.,
\emph{Prime Number Theorem And \dots} (the PNT+ project), a formalization project in Lean~4,
\url{https://github.com/AlexKontorovich/PrimeNumberTheoremAnd}; used at the commit named in Appendix~\ref{app:lean}.

\bibitem[Sar81]{Sar81}
P.~Sarnak,
\emph{Asymptotic behavior of periodic orbits of the horocycle flow and Eisenstein series},
Comm.\ Pure Appl.\ Math.\ \textbf{34} (1981), no.~6, 719--739,
\href{https://doi.org/10.1002/cpa.3160340602}{doi:\allowbreak 10.1002/cpa.3160340602}.

\bibitem[Sch05]{Sch05}
R.~Schubert,
\emph{Semiclassical behaviour of expectation values in time evolved Lagrangian
states for large times},
Commun.\ Math.\ Phys.\ \textbf{256} (2005), no.~1, 239--254,
\href{https://doi.org/10.1007/s00220-005-1319-4}{doi:\allowbreak 10.1007/s00220-005-1319-4},
\href{https://arxiv.org/abs/math/0402038}{arXiv:math/0402038}.

\bibitem[Sch07]{S}
R.~Schubert,
\emph{Semiclassical wave propagation for large times},
\href{https://arxiv.org/abs/0705.0134v1}{arXiv:0705.0134v1} (2007);
withdrawn by the author in v2 (2011).

\bibitem[STV12]{STV12}
R.~Schubert, R.~O.~Vallejos and F.~Toscano,
\emph{How do wave packets spread? Time evolution on Ehrenfest time scales},
J.\ Phys.\ A \textbf{45} (2012), no.~21, 215307,
\href{https://doi.org/10.1088/1751-8113/45/21/215307}{doi:\allowbreak 10.1088/1751-8113/45/21/215307},
\href{https://arxiv.org/abs/1112.4056}{arXiv:1112.4056}.

\bibitem[TH91]{TH91}
S.~Tomsovic and E.~J.~Heller,
\emph{Semiclassical dynamics of chaotic motion: unexpected long-time accuracy},
Phys.\ Rev.\ Lett.\ \textbf{67} (1991), no.~6, 664--667,
\href{https://doi.org/10.1103/PhysRevLett.67.664}{doi:\allowbreak 10.1103/PhysRevLett.67.664}.

\bibitem[TH93]{TH93}
S.~Tomsovic and E.~J.~Heller,
\emph{Long-time semiclassical dynamics of chaos: the stadium billiard},
Phys.\ Rev.\ E \textbf{47} (1993), no.~1, 282--299,
\href{https://doi.org/10.1103/PhysRevE.47.282}{doi:\allowbreak 10.1103/PhysRevE.47.282}.

\bibitem[Yak08]{Y}
S.~B.~Yakubovich,
\emph{Uncertainty principles for the Kontorovich--Lebedev transform},
Math.\ Model.\ Anal.\ \textbf{13} (2008), no.~2, 289--302,
\href{https://doi.org/10.3846/1392-6292.2008.13.289-302}{doi:\allowbreak 10.3846/1392-6292.2008.13.289-302}.

\bibitem[Zag81]{Zag81}
D.~Zagier,
\emph{Eisenstein series and the Riemann zeta-function},
in: Automorphic Forms, Representation Theory and Arithmetic (Bombay, 1979),
Tata Inst.\ Fund.\ Res.\ Studies in Math.\ 10, Springer, 1981, pp.~275--301,
\href{https://doi.org/10.1007/978-3-662-00734-1_10}{doi:\allowbreak 10.1007/978-3-662-00734-1\_10}.

\end{thebibliography}

\end{document}
